\documentclass[11pt,a4paper]{article}

\usepackage[T1]{fontenc}
\usepackage{lmodern}
\usepackage{amsmath,amssymb,amsthm}
\usepackage{booktabs}
\usepackage{needspace}
\usepackage{microtype}
\usepackage[margin=2.7cm]{geometry}
\usepackage{xcolor}
\usepackage[colorlinks=true,linkcolor=blue!50!black,citecolor=blue!50!black,urlcolor=blue!50!black]{hyperref}
\usepackage{siunitx}

\newtheorem{theorem}{Theorem}
\newtheorem{proposition}[theorem]{Proposition}
\newtheorem{lemma}[theorem]{Lemma}
\newtheorem{corollary}[theorem]{Corollary}
\theoremstyle{definition}
\newtheorem{conjecture}[theorem]{Conjecture}
\newtheorem{remark}[theorem]{Remark}
\newtheorem{definition}[theorem]{Definition}
\newtheorem{observation}[theorem]{Numerical observation}

\DeclareMathOperator{\arrgg}{arg}
\newcommand{\Rr}{R}
\newcommand{\Kk}{K}
\newcommand{\eee}{\mathrm{e}}
\newcommand{\ii}{\mathrm{i}}
\newcommand{\chat}{\hat{c}_1}

\title{A first-mode law for the separation between\\
the Kneser continuation and the regular iteration on $(1,\eta)$}

\author{%
  Guanghao Li\thanks{Correspondence: \texttt{wsmslgh@gmail.com}.  Code, data and
  a Chinese-language laboratory report accompany this note; see
  Section~\ref{sec:repro}.}%
}

\date{26 September 2026}

\begin{document}
\maketitle

\begin{abstract}
For $1<b<\eta=\eee^{1/\eee}$, we investigate the separation between the
regular super-exponential $\Rr_b$ and the boundary value $\Kk_b$ of a
Kneser-type continuation in the base.  With both solutions normalised by
$F(0)=1$, write $\Kk_b=\Rr_b\circ(\mathrm{id}+P_b)$, where $P_b$ is
$1$-periodic.  The measured separation is dominated by its first Fourier mode:
\[
 D_b(z):=\Kk_b(z)-\Rr_b(z)
 \approx c_1(b)\Rr_b'(z)(\eee^{2\pi\ii z}-1).
\]
Across the sampled bases, $|\hat c_1|=|c_1/\Lambda|$ varies by less than $5\%$,
where $\Lambda=\exp(4\pi^2/\log\lambda)$ and $\lambda$ is the attracting
multiplier.  The elementary parameter relation $\log b=\lambda\eee^{-\lambda}$
proves that this candidate scale has saddle-node exponent
$\log\Lambda\sim-C_\eta/\sqrt{\eta-b}$, with
$C_\eta=20.3508008\ldots$; its relation to the separation remains conjectural.
A separate computation of the parabolic horn map for $u\mapsto\eee^u-1$
gives $|A_1|=0.08905843641221$ in a fixed Fatou-coordinate normalisation.
The measured $|\hat c_1|$ increases toward that value.  Matching the direction
of the transition suggests the inverse horn map as the limiting object:
its coefficients satisfy $B_1=-A_1$ and $B_2=-A_2+2\pi\ii A_1^2$.
At six sampled bases, $|c_2/c_1^2|$ increases toward the candidate target
$|B_2/B_1^2|=1.57442268552971\ldots$.
All ratios $c_n/c_1^n$ are invariant under time translations; a distinct,
larger formal conjugacy relation singles out the quadratic coefficient and
iterative residue of the zero-constant-mode cylinder germ.
We then show that the Kneser continuation can be eliminated altogether:
a welding identity expresses every $c_n$, up to an additive $O(\Lambda)$, through
the Fourier modes $t_n$ of the transition map $T_b=\Rr_b^{-1}\circ S_b$
between the regular solutions at the two real fixed points,
$c_n/\Lambda^n=t_n\eee^{-2\pi\ii nt_0}+O(\Lambda)$.  This reproduces the
ladder values, arguments included, to their full precision, and reduces the
horn-map conjectures to a confluence statement for the saddle-node unfolding.
We prove that confluence statement: the Kneser-free invariants converge to
the inverse horn-map coefficients for every $n$.  A unisingular
parabolic-renormalisation theorem proves $B_1\ne0$, giving the limiting phase
and normalised ratios.  On the Kneser side, we prove the decay estimate for the
sewing solution with explicit constants.  That solution is characterised uniquely by its two regular
representations and the height of its seam.  The base-continued family has it
as boundary value whenever its $\theta$-representations stay uniformly
bounded along the ladder.
The sewn family varies holomorphically near every real $b\in(1,\eta)$;
near $\eta$ its imaginary part at real heights has an explicit leading
$\Lambda$-profile and is not identically zero on any interval about $0$.
Computing $T_b$ directly up to $\eta-b=3\cdot10^{-5}$ confirms the limits for
$n=1,2,3$ with deficits $\propto|\log\lambda|^2$.
Numerical consistency checks of the separation coefficients are reported
separately from exact identities; their quoted decimal comparisons
have no certified error bounds.
An interval-certified quasiconformal argument extends the exterior
Kneser family across a non-cusp Shell--Thron boundary subarc on a common
complex neighbourhood of the full height segment $[0,1]$.
Identification with the inner sewn family remains open.
\end{abstract}

\section{Introduction}

Tetration is the solution of the Abel-type functional equation
\begin{equation}\label{eq:abel}
  F(z+1)=b^{F(z)},\qquad F(0)=1 .
\end{equation}
For $b>\eta:=\eee^{1/\eee}$ the two fixed points of $w\mapsto b^{w}$ are a
complex conjugate pair, and Kneser's construction~\cite{Kneser1950} produces
the solution that is real-analytic on $(-2,\infty)$ and tends to the two fixed
points as $\operatorname{Im}z\to\pm\infty$.  Its conformal uniformisation
is uniquely characterised by the initial-region criterion
of~\cite{TrappmannKouznetsov2011}; high-precision algorithms for it are
well developed~\cite{Kouznetsov2009,KouznetsovTrappmann2010}.

For $1<b<\eta$ the situation is different.  The map $E_b(w)=b^{w}$ has two
\emph{real} fixed points
\begin{equation}\label{eq:fps}
  L=\eee^{\lambda}\ \ (\text{attracting},\ 0<\lambda<1),
  \qquad
  L_2=\eee^{\lambda_2}\ \ (\text{repelling},\ \lambda_2>1),
\end{equation}
and the classical Koenigs--Schr\"oder construction at the attracting point
yields a real-analytic solution $\Rr_b$ of~\eqref{eq:abel} by elementary means.
This is the solution usually called ``the'' tetration in that range; it is one
of the four regular super-exponentials studied for $b=\sqrt2$
in~\cite{KouznetsovTrappmann2010}, and it is the reason why the interior of the
Shell--Thron region is explicitly excluded from~\cite{PaulsenCowgill2017}.

There is, however, a second natural solution.  The two-fixed-point construction
is holomorphic in the base across the Shell--Thron boundary, so one may take
\begin{equation}\label{eq:K}
  \Kk_b:=\lim_{\varepsilon\to0^{+}}\kappa_{b+\ii\varepsilon},
\end{equation}
the boundary value on $(1,\eta)$ of the Kneser-type family continued from the
upper half of the region.  Paulsen~\cite{Paulsen2019} proved that $\Kk_b$ is
\emph{not} the regular solution: there is no open subinterval of $(1,\eta)$ on
which $\Kk_b=\Rr_b$.  His proof is an analytic contradiction argument and gives
no rate; the only quantitative datum in that paper is a single number for
$b=\sqrt2$, obtained from a $120$-digit contour computation,
\begin{equation}\label{eq:paulsen}
  \operatorname{Im}\kappa_{\sqrt2^{+}}(1/2)
  \approx -1.18899697185401045226976795872715599664\cdot10^{-48},
\end{equation}
reported there as evidence that the limit is not real, and not compared with
any scale.

This note is about the size and the shape of $D_b:=\Kk_b-\Rr_b$.  The
statements we can prove are elementary (Section~\ref{sec:setup} and
Proposition~\ref{prop:firstmode}); the substance is numerical, and we have tried
to be explicit throughout about which is which (Section~\ref{sec:claims}).

\paragraph{Summary of findings.}
\begin{enumerate}
\item The measured separation is dominated by one complex coefficient
  (Section~\ref{sec:firstmode}).  With $P_b(0)=0$, the leading profile is
  $c_1(\eee^{2\pi\ii z}-1)$; omitting the constant term creates a
  $z$-dependent error.
\item The normalised modulus $|\hat c_1|=|c_1/\Lambda|$ stays near
  $0.0850$--$0.0889$ across the reported data (Section~\ref{sec:invariant}).
  This comparison fixes the imaginary origin: the modulus is invariant under
  real, but not arbitrary complex, time translations.
\item The resolved low-order spectrum is compatible with
  $c_n=\kappa_n c_1^n$ for bounded low-order $\kappa_n$
  (Section~\ref{sec:spectrum}).  It does not determine the large-$n$ growth
  or the nearest singularity.  A uniform strip-bound proof remains open
  (Section~\ref{sec:nostrip}).
\item The scale $\Lambda$, independently of its conjectured relation to $D_b$,
  has the proved saddle-node exponent of Proposition~\ref{prop:saddle}.
  The measured $|\hat c_1|$ increases toward the first parabolic horn
  coefficient in the chosen normalisation (Section~\ref{sec:horn}).
\item The reversed transition gives the second-mode target
  $|B_2|=0.012487\ldots$, and the six measured ratios
  $|\kappa_2|/|\kappa_2^\infty|$ rise from $0.873$ to $0.979$
  (Section~\ref{sec:mode2}).  This is evidence for a limit, not its proof.
\item Translation invariance permits comparisons at every order.  For the
  distinct problem of classifying a simple parabolic cylinder germ under
  tangent-to-identity formal conjugacy, the parameters are its quadratic
  coefficient $a_2$ and residue $\rho$, with
  $\kappa_2=2\pi\ii(\frac12-\rho)$ (Section~\ref{sec:onlytwo}).
  This classification does not prohibit testing higher Fourier coefficients.
\item The separation is a transition-map invariant (Section~\ref{sec:transition}):
  $\hat c_n=t_n\eee^{-2\pi\ii nt_0}+O(\Lambda)$ for the Kneser-free map
  $T_b=\Rr_b^{-1}\circ S_b$ (Theorem~\ref{thm:welding}); the conjectures follow for
  every $n$ from a confluence statement (Theorem~\ref{thm:reduction}), and the
  limits are confirmed numerically to $\eta-b=3\cdot10^{-5}$.
\end{enumerate}

\section{Setup}\label{sec:setup}

Throughout, $1<b<\eta$ and $E(w)=b^{w}$.

\begin{lemma}[Exact parametrisation of the family]\label{lem:param}
Let $L$ be a fixed point of $E$ and $\lambda=E'(L)=L\log b$ its multiplier.
Then
\begin{equation}\label{eq:param}
  L=\eee^{\lambda},\qquad \log b=\lambda\,\eee^{-\lambda}.
\end{equation}
The map $\lambda\mapsto \lambda\eee^{-\lambda}$ increases on $(0,1)$ and
decreases on $(1,\infty)$ with maximum $\eee^{-1}=\log\eta$ at $\lambda=1$;
consequently $\lambda\in(0,1)$ parametrises $(1,\eta)$ through the attracting
fixed point, $\lambda_2\in(1,\infty)$ parametrises it through the repelling one,
and the two branches meet at the parabolic base $b=\eta$.
\end{lemma}

\begin{proof}
$L=b^{L}$ gives $\log L=L\log b=\lambda$, hence $L=\eee^{\lambda}$ and
$\log b=\lambda/L=\lambda\eee^{-\lambda}$.  The monotonicity statements are
those of $\lambda\eee^{-\lambda}$.
\end{proof}

We write $\lambda$ for the multiplier at the attracting fixed point and set
\begin{equation}\label{eq:Lambda}
  h:=\frac{2\pi}{|\log\lambda|},\qquad
  \Lambda:=\exp\!\Bigl(\frac{4\pi^{2}}{\log\lambda}\Bigr)=\eee^{-2\pi h}.
\end{equation}
The quantity $h$ is the modulus of the imaginary period of the regular
solution: if $\sigma$ denotes the Koenigs coordinate at $L$, normalised by
$\sigma(E(w))=\lambda\sigma(w)$, then $\Rr(z)=\sigma^{-1}(c\lambda^{z})$ for a
suitable constant $c$, and $\lambda^{z}$ has period $2\pi\ii/\log\lambda$, of
modulus $h$.  Note in particular that, on the real segment, $\Rr$ is
\emph{exactly periodic} in the imaginary direction and therefore has no limit as
$\operatorname{Im}z\to\pm\infty$; the two-sided limit characterisation of the
Kneser solution degenerates precisely on $(1,\eta)$, which is the structural
reason the interval is delicate.

\section{The first-mode law}\label{sec:firstmode}

\begin{proposition}\label{prop:firstmode}
Let $\Rr$ and $\Kk$ be holomorphic solutions of~\eqref{eq:abel}, both
normalised by $F(0)=1$.  Assume a branch of $\Rr^{-1}$ is defined on the
required images of $\Kk$ and is compatible with the shift:
$\Rr^{-1}(E(w))=\Rr^{-1}(w)+1$ wherever used.  Put
$P=\Rr^{-1}\circ\Kk-\mathrm{id}$.  Then $P(z+1)=P(z)$ on the common
domain and $P(0)=0$, with the inverse branch fixed by $\Rr^{-1}(1)=0$.

If $P=\sum_{n\ge0}c_n\eee^{2\pi\ii nz}$ converges locally uniformly
near $0$ and the points under consideration, then
$c_0=-\sum_{n\ge1}c_n$.  Define
$Q(z)=\sum_{n\ge2}c_n(\eee^{2\pi\ii nz}-1)$.  Taylor expansion gives
\begin{equation}\label{eq:firstmode}
 D(z)=c_1\Rr'(z)(\eee^{2\pi\ii z}-1)+\Rr'(z)Q(z)+O(P(z)^2),
\end{equation}
locally where the second derivative of $\Rr$ is bounded along the Taylor
segments.  Omitting $Q$ at quadratic accuracy additionally requires
$\Rr'Q=O(P^2)$.
\end{proposition}

\begin{proof}
The functional equation and the compatible inverse branch imply
$(\Rr^{-1}\circ\Kk)(z+1)=(\Rr^{-1}\circ\Kk)(z)+1$.
The normalisation gives $P(0)=0$; evaluating the convergent series at zero
yields the identity for $c_0$.  Thus
$P=c_1(\eee^{2\pi\ii z}-1)+Q$; substitute this into
$\Rr(z+P(z))-\Rr(z)=\Rr'(z)P(z)+O(P(z)^2)$.
\end{proof}

The numerical question is whether the tail $Q$ and the Taylor remainder
are small enough for the leading profile to dominate.  We test this through
\begin{equation}\label{eq:q}
 q(z):=\frac{D(z)}{\Rr'(z)(\eee^{2\pi\ii z}-1)},
\end{equation}
away from zeros of the denominator.  First-mode dominance predicts an
approximately constant $q$.  Table~\ref{tab:qz} shows $|q|$ at four values
of $z$, including $z=\tfrac12-\tfrac{\ii}{2}$, where the first positive
mode is amplified by $\eee^\pi\approx23.14$ and the first negative mode
is suppressed by the same factor.

\begin{table}[htbp]
\centering
\caption{The first-mode law~\eqref{eq:firstmode}, tested through~\eqref{eq:q}.
The last column is $\max|q|/\min|q|$ over the four $z$.  For comparison, the
naive normalisation $D/(\Rr'\eee^{2\pi\ii z})$, which omits the constant mode,
has spread $1.917$ at every point of the data set.}
\label{tab:qz}
\small\setlength{\tabcolsep}{4pt}
\begin{tabular}{lccccc}
\toprule
base & $z=\tfrac12$ & $z=-\tfrac12$ & $z=\tfrac14$ & $z=\tfrac12-\tfrac{\ii}{2}$ & spread \\
\midrule
$1.1+0.0005\ii$ & \num{1.987636e-9} & \num{1.987636e-9} & \num{1.987636e-9} & \num{1.987637e-9} & $1.0000$\\
$1.05+0.003\ii$ & \num{1.422053e-7} & \num{1.422053e-7} & \num{1.422053e-7} & \num{1.422061e-7} & $1.0000$\\
\bottomrule
\end{tabular}
\end{table}

\section{The normalised first coefficient $\chat=c_1/\Lambda$}\label{sec:invariant}

Since $\Kk$ is defined as a boundary value~\eqref{eq:K}, we compute $c_1$ at
$b+\ii y$ for a ladder of $y\downarrow0$ and extrapolate.  The relevant
normalisation is by the scale $\Lambda$ of~\eqref{eq:Lambda}:
\begin{equation}
  \chat(b):=\lim_{y\to0^{+}}\frac{c_1(b+\ii y)}{\Lambda(b+\ii y)} .
\end{equation}

\begin{remark}[Translation invariance]\label{rem:invariance}
A simultaneous shift $\Kk_\delta(z)=\Kk(z+\delta)$,
$\Rr_\delta(z)=\Rr(z+\delta)$ gives
$P_\delta(z)=P(z+\delta)$ and $c_n\mapsto c_n\eee^{2\pi\ii n\delta}$.
Thus every $\kappa_n=c_n/c_1^n$, for $c_1\ne0$, is invariant under
complex translations.  In contrast,
$|c_1|\mapsto |c_1|\eee^{-2\pi\operatorname{Im}\delta}$:
$|\hat c_1|$ is invariant only under real shifts.  Its comparison with a
parabolic horn coefficient requires a fixed imaginary normalisation.
Arguments of coefficients are meaningful once an origin has been specified.
Independent source and target translations additionally change the constant
mode, while leaving the ratios of positive modes invariant.
\end{remark}

Table~\ref{tab:chat} collects the values.  The first six rows are computed
here; the seventh is Paulsen's number~\eqref{eq:paulsen} divided by
$2\Rr_{\sqrt2}'(1/2)\Lambda(\sqrt2)$; the last two use values of $\Kk$ published
from \texttt{fatou.gp} runs, i.e.\ a completely independent implementation in
PARI/GP, at two complex bases in the interior of the Shell--Thron region.

\begin{table}[htbp]
\centering
\caption{The near-invariant $|\chat|$.  $|D|$ ranges over forty-five orders of
magnitude; $|\chat|$ varies by $4.6\%$.  The column $|D|$ is the value of
$|D(\tfrac12)|$ at the lowest rung of the ladder actually computed, not the
limit.  The last column compares with the horn-map coefficient
$|A_1|=0.08905843641221$ of the parabolic limit computed in
Section~\ref{sec:horn}.  Rows marked $\dagger$ are complex bases, where $\chat$
is a sample of the same function off the real segment and the last column is not
meaningful.}
\label{tab:chat}
\small\setlength{\tabcolsep}{3.5pt}
\begin{tabular}{lcccccl}
\toprule
$b$ & $\lambda$ (or $|\lambda|$) & $\Lambda$ & $|D|$ & $|\chat|$ & $|\chat|/|A_1|$ & source \\
\midrule
$1.05$      & $0.051362$ & \num{1.68e-6}  & \num{1.01e-8}  & $0.085017$ & $0.95462$ & this work, $20$ digits\\
$1.10$      & $0.105964$ & \num{2.30e-8}  & \num{3.24e-10} & $0.086427$ & $0.97045$ & this work, $20$ digits\\
$1.15$      & $0.164802$ & \num{3.10e-10} & \num{7.08e-12} & $0.087205$ & $0.97919$ & this work, $20$ digits\\
$1.20$      & $0.229312$ & \num{2.28e-12} & \num{7.30e-14} & $0.087732$ & $0.98511$ & this work, $20$ digits\\
$1.25$      & $0.301735$ & \num{4.91e-15} & \num{2.04e-16} & $0.088125$ & $0.98952$ & this work, $24$ digits\\
$1.30$      & $0.385935$ & \num{9.82e-19} & \num{1.16e-19} & $0.088434$ & $0.99299$ & this work, $32$ digits\\
$\sqrt2$    & $0.693147$ & \num{1.66e-47} & \num{1.19e-48} & $0.088941$ & $0.99868$ & \cite{Paulsen2019}, $120$ digits\\
\midrule
$0.8+0.4\ii^{\dagger}$ & $0.387654$ & \num{1.05e-3} & \num{9.86e-5} & $0.085592$ & --- & \texttt{fatou.gp}, $80$ digits\\
$1+\ii^{\dagger}$      & $0.710560$ & \num{2.13e-2} & \num{4.41e-3} & $0.087272$ & --- & \texttt{fatou.gp}, $80$ digits\\
\bottomrule
\end{tabular}
\end{table}

\begin{remark}[Where the model stops]
The published \texttt{fatou.gp} set contains a third interior base, $b=2+\ii$,
with $|\lambda|=0.98767$ and hence $|\Lambda|=0.806$.  There the ``exponentially
small'' quantity is not small at all, the truncation to one mode is not
justified, and indeed the same computation returns $|\chat|=0.1065$.  We regard
this as a useful marker of the boundary of validity rather than as a
counterexample.
\end{remark}

\paragraph{Internal consistency.}
At $b=1.10$ the value $|\chat|=0.08642738$ is obtained from six
settings (working precision $16$ and $20$ digits, $\theta$ sample height
$0.10$ and $0.15$) and, separately, from a finer ladder reaching
$y=\num{5e-4}$, which gives $0.08642720$; the absolute difference is $1.8\cdot10^{-7}$ (about $2.1$ parts per million).  At
$b=1.20$ and $b=1.25$ the same quantity computed at working precisions $20$,
$24$ and $32$ digits agrees to all seven digits printed, with signal-to-noise
ratios up to $2\cdot10^{21}$.

\section{The Fourier spectrum of $P$}\label{sec:spectrum}

Rather than fitting, one can measure $P$ directly: for $z$ on a horizontal line
$\operatorname{Im}z=y_0$ compute $P(z)=\Rr^{-1}(\Kk(z))-z$, choosing the branch
of the logarithm in $\Rr^{-1}$ nearest to $z$, and transform.  Lowering $y_0$
amplifies the mode $n$ by $\eee^{-2\pi ny_0}$ and so makes higher modes
observable; raising $y_0$ does the same for negative modes.

\begin{table}[htbp]
\centering
\caption{Spectrum of $P$ at $b=1.1+\num{5e-4}\ii$
($\Lambda=\num{2.3001e-8}$, $h=2.7992$).  The values of $|c_0|,|c_1|,|c_2|$ are
stable across the three sampling lines; $|c_3|$ is not resolved.  The
mode-to-mode ratio is $\num{2.86e-9}$, an order of magnitude below $\Lambda$.}
\label{tab:spectrum}
\begin{tabular}{lcccc}
\toprule
$y_0$ & $|c_0|$ & $|c_1|$ & $|c_2|$ & $|c_3|$\\
\midrule
$0.0$  & \num{1.9876e-9} & \num{1.9876e-9} & \num{5.6899e-18} & noise\\
$-0.5$ & \num{1.9876e-9} & \num{1.9876e-9} & \num{5.6899e-18} & \num{1.77e-26}\\
$-1.0$ & \num{1.9876e-9} & \num{1.9876e-9} & \num{5.6899e-18} & \num{1.32e-26}\\
\bottomrule
\end{tabular}
\end{table}

Three facts come out of Table~\ref{tab:spectrum} and its analogues at
$b=1.05,1.15$.

\begin{observation}\label{obs:spectrum}
\ \begin{enumerate}
\item $c_0/c_1=-1$ to ten significant digits, confirming the normalisation
  structure of Proposition~\ref{prop:firstmode}.
\item Sampling on $\operatorname{Im}z=+0.5$ and $+1.0$ gives $|c_{-1}|$ equal to
  $\num{4.0e-23}$ and $\num{1.3e-23}$ respectively.  A genuine coefficient does
  not depend on the sampling line; the measured values decrease as the line is
  raised, i.e.\ they are the numerical floor.  Hence
  $|c_{-1}|/|c_1|\lesssim5\cdot10^{-15}$: the hypothesis that $P$ has no
  negative modes, which expresses the boundedness of $\Kk$ in the upper
  half-plane, is verified over fifteen orders of magnitude.
\item Writing $\kappa_n=c_n/c_1^{\,n}$, the measured $\kappa_2$ is $O(1)$:
  \[
    |\kappa_2|=1.36643,\ 1.44023,\ 1.47531
    \quad\text{and}\quad
    \arg\kappa_2=1.41804,\ 1.47802,\ 1.51415
  \]
  at $b=1.05,1.10,1.15$.  By Remark~\ref{rem:invariance} these are translation
  invariants; they need not survive arbitrary nonlinear changes of time coordinate.
\end{enumerate}
\end{observation}

\subsection{What the spectrum does not prove}\label{sec:nostrip}

A periodic function bounded by $M$ on a horizontal line at depth $d$ satisfies
$|c_n|\le M\eee^{-2\pi nd}$ for positive modes.  A bound at
$d=h=2\pi/|\log\lambda|$ with $M$ uniform as $b\to\eta^-$ would therefore
establish $|c_1|=O(\Lambda)$.  It would not, on its own, establish a nonzero
limit of $c_1/\Lambda$.

The few resolved coefficients do not rule out this approach.  For example,
if the positive-mode series is written as
$P-c_0=\sum_{n\ge1}\kappa_n(c_1\xi)^n$, $\xi=\eee^{2\pi\ii z}$,
then its radius in $\xi$ is determined by
$L_\kappa=\limsup_{n\to\infty}|\kappa_n|^{1/n}$.
Only under the additional, unproved condition $L_\kappa=1$ would its
singularity depth be
\begin{equation}\label{eq:depth}
 d=\frac{-\log|c_1|}{2\pi}
   =h+\frac{\log(1/|\chat|)}{2\pi}.
\end{equation}
For general finite positive $L_\kappa$, subtract
$\log L_\kappa/(2\pi)$ from the right-hand side.  The final equality
in~\eqref{eq:depth} is an algebraic identity from $c_1=\Lambda\chat$;
computing its two sides is not an independent measurement of a singularity.
Moreover, a singularity deeper than $h$ would be compatible with analyticity
on the desired strip.  The unresolved issue is a uniform bound obtained
independently of the unknown coefficient.

In a domain where the Koenigs coordinate $\sigma$ and a logarithm branch
exist, $p(z)=\lambda^{-z}\sigma(\Kk(z))$ is periodic and
$P=\log(p/c)/\log\lambda$.  Zeros of $p$ can create logarithmic
singularities, but other obstructions to continuation also require control.
The present low-order spectrum neither locates all such obstructions nor
settles the origin of the scale $\Lambda$.

\section{The parabolic end}\label{sec:parabolic}

By Lemma~\ref{lem:param}, $\log b=\lambda\eee^{-\lambda}$.  Writing
$\lambda=1-\varepsilon$ and expanding,
\[
  \log b=\eee^{-1}\Bigl(1-\frac{\varepsilon^{2}}{2}-\frac{\varepsilon^{3}}{3}
  +O(\varepsilon^{4})\Bigr),
  \qquad\text{so}\qquad
  \varepsilon=\sqrt{2\eee\,(\log\eta-\log b)}\,\bigl(1+O(\varepsilon)\bigr),
\]
and $|\log\lambda|=\varepsilon+\varepsilon^{2}/2+\cdots$.  (Numerically, the
ratio of $|\log\lambda|$ to $\sqrt{2\eee(\log\eta-\log b)}$ is
$1.0386,\,1.0108,\,1.0033,\,1.0003$ for
$\eta-b=10^{-2},10^{-3},10^{-4},10^{-6}$.)  Substituting into~\eqref{eq:Lambda}
and using $\log\eta-\log b\sim(\eta-b)/\eta$:

\begin{proposition}\label{prop:saddle}
As $b\to\eta^{-}$,
\begin{equation}\label{eq:saddle}
  \Lambda(b)=\exp\!\Bigl(-\frac{C_\eta}{\sqrt{\eta-b}}\,(1+o(1))\Bigr),
  \qquad
  C_\eta=\frac{4\pi^{2}}{\sqrt{2\eee^{\,1-1/\eee}}}=20.3508008\ldots
\end{equation}
\end{proposition}

This is the classical shape $\exp(-C/\sqrt{\varepsilon})$ of exponentially small
splitting at a saddle-node bifurcation.  It is the right shape for structural
reasons: at $b=\eta$ the two real fixed points~\eqref{eq:fps} merge into the
parabolic point $w=\eee$ of $\eee^{w/\eee}$, and the two regular
super-exponentials of the parabolic base~\cite{KouznetsovTrappmann2012} are the
attracting- and repelling-petal Fatou parametrisations.  The decay of $\Lambda$ motivates a vanishing-separation conjecture.  It does
not establish it: control of $\hat c_1$, higher modes, and $\Rr_b'$ is
also needed to transfer the scale to $D_b$.

For unfoldings of a parabolic germ into a pair of fixed points, the analytic
classification is by an ``unfolded'' \'Ecalle--Voronin modulus which converges,
as the parameter tends to the parabolic value, to the modulus of the limiting
germ~\cite{MRR2004,RousseauTeyssier2008}; exponentially small differences
between the two natural parametrisations are measured by the horn maps of that
limit~\cite{Ecalle1981,Voronin1981,Lavaurs1989,Shishikura2000,Oudkerk1999}.
Our data are consistent with such a statement, and we can now test it directly.

\subsection{The horn map of the parabolic limit}\label{sec:horn}

At $b=\eta$ put $w=\eee(1+u)$; the map becomes the parabolic germ
\begin{equation}\label{eq:germ}
  f(u)=\eee^{u}-1=u+\tfrac{u^{2}}{2}+\tfrac{u^{3}}{6}+\cdots,
\end{equation}
with attracting petal on $\operatorname{Re}u<0$ and repelling petal on
$\operatorname{Re}u>0$, and with the (divergent) formal Abel function
\begin{equation}\label{eq:alpha}
  \alpha(u)=-\frac{2}{u}+\frac{\log(-u)}{3}+\sum_{k\ge1}c_ku^{k},
  \qquad \alpha(f(u))=\alpha(u)+1 .
\end{equation}
Both Fatou coordinates are normalised by the \emph{same} $\alpha$:
\[
  \Phi_{\mathrm{att}}(u)=\lim_{n\to\infty}\bigl[\alpha(f^{n}(u))-n\bigr],
  \qquad
  \Phi_{\mathrm{rep}}(u)=\lim_{n\to\infty}\bigl[\alpha(f^{-n}(u))+n\bigr],
\]
the branch of $\log(-u)$ being continued along the orbit.  On the upper gate the
transition $h=\Phi_{\mathrm{rep}}\circ\Phi_{\mathrm{att}}^{-1}$ commutes with
$z\mapsto z+1$ and tends to the identity as $\operatorname{Im}z\to+\infty$
(both coordinates being asymptotic to the same $\alpha$), so
\begin{equation}\label{eq:horn}
  h(z)-z=\sum_{n\ge1}A_n\,\eee^{2\pi\ii nz},
\end{equation}
which is the lifted \'Ecalle--Voronin horn map in the chosen normalisation.
The modulus $|A_1|$ is unchanged under simultaneous real translations, but
changes under imaginary translations; the specified $\alpha$ fixes that freedom.  Computing
$h(z)-z$ at $N$ equally spaced points of a horizontal line and transforming
gives
\begin{equation}\label{eq:A1}
  \boxed{\;|A_1|=0.08905843641221\;}\qquad |A_2|=0.037349258777,\qquad
  |A_3|=0.02246384 .
\end{equation}
These values are unchanged, to all digits shown, when the working precision runs
over $30$--$50$ digits, the number of terms of~\eqref{eq:alpha} over $20$--$60$,
the matching radius over $0.005$--$0.02$, the transform size over $N=16,24,32$,
and the sampling height over $\operatorname{Im}z=1,1.5,2$; the constant term
$A_0$ comes out as zero to $10^{-47}$ or better, as it must.

Comparing with Table~\ref{tab:chat}: the hyperbolic-side quantity $|\chat(b)|$
increases toward $|A_1|$ on the sampled bases, the ratio being $0.95462$ at
$b=1.05$ and $0.99868$ at $b=\sqrt2$.  The value at $\sqrt2$ is not ours: it is
Paulsen's $120$-digit contour computation~\eqref{eq:paulsen}, so the $0.13\%$
agreement there is between two computations sharing neither code, method, nor
author.

\begin{conjecture}\label{conj:horn}
$\displaystyle\lim_{b\to\eta^{-}}|\chat(b)|=|A_1|$, the first horn-map
coefficient~\eqref{eq:A1} of the parabolic germ~\eqref{eq:germ}.
\end{conjecture}

This is the expected shape of an unfolded \'Ecalle--Voronin
statement~\cite{MRR2004,RousseauTeyssier2008}: the exponentially small splitting
of the two natural parametrisations of the perturbed map, measured in the units
forced by the multiplier, converges to the modulus of the parabolic limit.  The numerical evidence concerns $n=1$; the limit and the required
normalisation match have not been established.

\subsection{The second invariant: the reversed transition map}\label{sec:mode2}

The obvious extension of Conjecture~\ref{conj:horn}, namely
$c_n/\Lambda^{n}\to A_n$, is inconsistent with the presently measured $n=2$ values by a factor of about three:
$|\hat c_2|=|\kappa_2|\,|\chat|^{2}$ measures $0.00993$, $0.01076$, $0.01125$ at
$b=1.05,1.10,1.15$ against $|A_2|=0.037349$.  The reason is a direction, not a
rescaling, and reversing it yields a substantially closer numerical target.
Neither comparison alone establishes a limiting relation.

The hyperbolic-side object is $P=\Rr^{-1}\circ\Kk-\mathrm{id}$, the attracting
(regular) time as a function of the repelling (Kneser) time, i.e.\ the
transition map $\alpha_{\mathrm{att}}\circ\alpha_{\mathrm{rep}}^{-1}$;
the $h$ of~\eqref{eq:horn} is $\Phi_{\mathrm{rep}}\circ\Phi_{\mathrm{att}}^{-1}$,
the other way round.  Subject to establishing this identification of the
Kneser time with the repelling time, the candidate limiting object is
$h^{-1}$, and inverting a $1$-periodic map tangent to the identity mixes the
modes: writing $h^{-1}(z)-z=\sum_{n\ge1}B_n\eee^{2\pi\ii nz}$,
\begin{equation}\label{eq:Bn}
  B_1=-A_1,\qquad B_2=-A_2+2\pi\ii A_1^{2},\qquad
  B_3=-\Bigl(A_1\bigl(2\pi\ii B_2+\tfrac{(2\pi\ii)^{2}}{2}B_1^{2}\bigr)
  +4\pi\ii A_2B_1+A_3\Bigr),
\end{equation}
so that $n=1$ is unaffected in modulus --- which is why
Conjecture~\ref{conj:horn} was insensitive to the direction --- while at $n=2$
the term $2\pi\ii A_1^{2}$, carrying the factor $2\pi$, changes the order of
magnitude.  Computing $\Phi_{\mathrm{att}}\circ\Phi_{\mathrm{rep}}^{-1}$
\emph{directly} by the algorithm of Section~\ref{sec:horn} run backwards, rather
than through~\eqref{eq:Bn}, gives the same numbers, with the constant term
$B_0$ vanishing to $10^{-59}$:
\begin{equation}\label{eq:B12}
  |B_1|=0.089058436412213331563\;(=|A_1|),\qquad
  |B_2|=0.012487384111564700811,
\end{equation}
\begin{equation}\label{eq:kappainf}
\begin{aligned}
 \kappa_2^\infty:=\frac{B_2}{B_1^2}
   &=-0.02532049309975755123+1.5742190652319516267\ii,\\
 |\kappa_2^\infty|&=1.5744226855297069\ldots
\end{aligned}
\end{equation}
unchanged in all digits shown over working precisions $40$--$60$ digits,
transform sizes $N=16,24,32$ and sampling heights $\operatorname{Im}z=1,1.5,2$.

\begin{proposition}\label{prop:horn-nonzero}
The first inverse upper horn coefficient satisfies $B_1\ne0$.
\end{proposition}
\begin{proof}
For $f(u)=\eee^u-1$, the negative real ray belongs to its immediate
parabolic basin $A$, since $u<f(u)<0$ there.  The only finite singular
value of $f$ is $-1$, an asymptotic value along that ray; $f$ has no
critical points.  By~\cite[Lemma~87]{Cheritat2022}, the restriction
$f|_A:A\to A$ has no other singular values.  Thus $f$ belongs to the
one-petal class $S_\infty$ of~\cite[Definition~8]{Cheritat2022}.
By~\cite[Theorem~9 and Corollary~11]{Cheritat2022}, its upper parabolic
renormalisation has exactly one attracting petal at $0$.  This
renormalisation is the exponential projection of the upper inverse horn
map, up to nonzero linear scalings of source and target
\cite[Appendix~A]{Cheritat2022}.  In our normalisation that projection is
$H(\zeta)=\zeta+2\pi\ii B_1\zeta^2+O(\zeta^3)$.
After normalising the multiplier to one, such linear scalings preserve
whether the quadratic coefficient vanishes.  One attracting petal makes
it nonzero, whence $B_1\ne0$.
\end{proof}
The hyperbolic-side $\kappa_2$, extrapolated to $y=0$ along the ladders, is then

\begin{center}
\begin{tabular}{llll}
\toprule
$b$ & $\kappa_2$ & $|\kappa_2|/|\kappa_2^{\infty}|$ & $|\hat c_2|/|B_2|$\\
\midrule
$1.05$ & $1.374049\,\eee^{+1.413420\ii}$ & $0.87273$ & $0.79532$\\
$1.10$ & $1.440851\,\eee^{+1.477690\ii}$ & $0.91516$ & $0.86188$\\
$1.15$ & $1.479111\,\eee^{+1.511255\ii}$ & $0.93946$ & $0.90077$\\
$1.20$ & $1.505578\,\eee^{+1.533318\ii}$ & $0.95627$ & $0.92800$\\
$1.25$ & $1.525582\,\eee^{+1.549420\ii}$ & $0.96898$ & $0.94875$\\
$1.30$ & $1.541516\,\eee^{+1.561918\ii}$ & $0.97910$ & $0.96538$\\
\bottomrule
\end{tabular}
\end{center}

\noindent
monotonically increasing, with the deficit $1-|\hat c_2|/|B_2|$ a varying multiple
$4.5$--$4.9$ of the $n=1$ deficit on the same bases, and with $\arg\kappa_2$
approaching $\arg\kappa_2^{\infty}=1.586879$ from below
($1.4134,1.4777,1.5113,1.5333,1.5494,1.5619$).  The last three rows are runs
made after the truncation of Remark~\ref{rem:n2} was removed; there the two
sampling lines agree to $10^{-6}$ at $b=1.25$ and $10^{-14}$ at $b=1.30$.  For $b=1.05$ and $b=1.10$ two mutually independent
ladders agree on $\kappa_2$ to four and six digits.

\begin{conjecture}\label{conj:horn2}
$\displaystyle\lim_{b\to\eta^{-}}\kappa_2(b)=\kappa_2^{\infty}$
of~\eqref{eq:kappainf}, the second coefficient of the reversed horn map of the
parabolic germ~\eqref{eq:germ}, normalised by the first.
\end{conjecture}

\subsection{How fast the limit is approached}\label{sec:deficit}

In the fixed imaginary normalisation, write
$\delta_1(b)=1-|\chat(b)|/|B_1|$ for the first-mode modulus deficit.
Over nine bases --- eight computed here and $b=\sqrt2$
from~\eqref{eq:paulsen} --- $\delta_1$ runs from $4.53\cdot10^{-2}$ at $b=1.05$
to $1.22\cdot10^{-3}$ at $b=\sqrt2$, and the natural variable for it is the multiplier scale
$\varepsilon:=|\log\lambda|$:

\begin{center}
\begin{tabular}{lrrrr}
\toprule
$b$ & $\varepsilon$ & $\eta-b$ & $\delta_1$ & $\delta_1/\varepsilon^{2}$\\
\midrule
$1.05$ & $2.96886$ & $0.39467$ & $0.0453065$ & $0.005140$\\
$1.10$ & $2.24465$ & $0.34467$ & $0.0295431$ & $0.005863$\\
$1.15$ & $1.80301$ & $0.29467$ & $0.0208085$ & $0.006401$\\
$1.20$ & $1.47267$ & $0.24467$ & $0.0148953$ & $0.006868$\\
$1.25$ & $1.19820$ & $0.19467$ & $0.0104955$ & $0.007310$\\
$1.30$ & $0.95209$ & $0.14467$ & $0.0070327$ & $0.007758$\\
$1.35$ & $0.71389$ & $0.09467$ & $0.0042041$ & $0.008249$\\
$1.40$ & $0.45443$ & $0.04467$ & $0.0018301$ & $0.008862$\\
$\sqrt2$ & $0.36651$ & $0.03045$ & $0.0012202$ & $0.009084$\\
\bottomrule
\end{tabular}
\end{center}

The historical exploratory fits to these nine values report a pure-power
exponent $p=1.737$ (log-RMS $0.040$).  For models
$\log\delta_1=\log C+p\log\varepsilon+\sum_{j=1}^m a_j\varepsilon^j$,
free fits with $m=1,2,3$ give $p=1.928,1.974,1.99893$ and log-RMS
$0.0032,0.00058,0.00017$.  These sample residuals suggest, but do not
identify, the candidate asymptotic law
\begin{equation}\label{eq:deficit}
  \delta_1(b)=C\varepsilon^2(1+O(\varepsilon)),\qquad
  \varepsilon=|\log\lambda|,\qquad C\approx0.01013.
\end{equation}
The previously printed uncertainty $\pm0.00002$ was neither a confidence
interval nor a certified error bound and is withdrawn.  The reported fixed
$p=2$ fit used three correction terms, whereas the reported $p=7/4$ fit used
one.  Their residuals are not a fair exponent comparison.  Matched correction
orders, sensitivity to omitted points, and data closer to $\varepsilon=0$
are required before excluding alternatives.  Adding correction parameters
also reduces residuals without necessarily improving extrapolation.

The parameter relation gives
$\eta-b\sim\varepsilon^2/(2\eee^{1-1/\eee})$ and $1-\lambda\sim\varepsilon$.
If the normalised deficit were differentiable in the unfolding parameter
with a nonzero linear term, a quadratic correction in $\varepsilon$ would
follow.  Those hypotheses have not been proved.  The saddle-node geometry
alone does not establish regularity of the normalised observable.
The measured range reaches only $\varepsilon\approx0.367$ and finite-range
fits change substantially with the choice of variable.

Paulsen publishes only $\operatorname{Im}\Kk_{\sqrt2}(1/2)$.  Using
$\operatorname{Im}D(1/2)\approx-2|c_1|\Rr'(1/2)\sin(\arg c_1)$ to
recover a modulus requires a phase estimate.  The corrected deficit
$0.0012202$ uses a phase extrapolated from our nearest two bases; it is not
an independent modulus measurement.  Omitting this correction gives
$0.0013210$.  The one-correction exponents for the corrected, uncorrected,
and omitted external point are $1.928,1.853,1.919$, respectively.
These are sensitivity checks, not proof of an exponent.  The reproducible
\texttt{deficit\_audit.py} audit separates eight internal points, nine points
with the phase correction, and nine without it.  Their cubic-correction
free exponents are approximately $1.99053$, $1.99889$, and $1.40683$.
More seriously, the two archived ladders at $b=1.35$ differ in
$|\hat c_1|$ by about $8.1\cdot10^{-5}$, about $22\%$ of the deficit
when expressed in deficit units.  Substituting archived alternative readings
substantially changes the fits.  This observed spread is a sensitivity
warning, not a probability model or a certified uncertainty interval.

Writing $\delta_\kappa=1-|\kappa_2|/|\kappa_2^\infty|$, the measured ratios are
\[
 \delta_\kappa/\delta_1
 =2.8090,\ 2.8717,\ 2.9093,\ 2.9356,\ 2.9557,\ 2.9718
 \quad (b=1.05,\ldots,1.30).
\]
Aitken extrapolation gives $3.04$, with no error bound or established
convergence model.  The exact identity
$\delta_2=1-(1-\delta_\kappa)(1-\delta_1)^2$ implies a difference of $2$
between the ratios only asymptotically, if $\delta_\kappa/\delta_1$ remains
bounded and both deficits vanish.

An earlier explanation posited the common relative change
$a_2\mapsto a_2(1+s)$, $a_3\mapsto a_3(1+s)$.  Its algebra is
\[
\begin{aligned}
 1-\frac{\kappa_2(s)}{\kappa_2^\infty}&=P\frac{s}{1+s},
 &P&=\frac{1-\rho}{1/2-\rho}=2.99514-0.03209\ii,\\
 \delta_1&=1-|1+s|=-\operatorname{Re}s+O(|s|^2).
\end{aligned}
\]
Consequently, for $q=(1-\kappa_2/\kappa_2^\infty)/\delta_1$,
$\operatorname{Re}(q/P)$ should approach $-1$, provided
$|s|^2=o(\delta_1)$.  The previous version mistakenly used $+1$.
The reported values $1.0224$, $1.0082$, $1.0021$, $0.9985$, $0.9960$, $0.9942$ do not
support that model.  Moreover, the modulus-deficit ratio depends on
$\operatorname{Re}(Ps)/(-\operatorname{Re}s)$, not generally on $|P|$.
The proposed target $|P|=2.995307$ and the phase-limit estimate derived
from this model are therefore withdrawn.  Direct phase measurements remain
usable, but a controlled limiting phase alignment remains open.

\begin{remark}
The fits above are superseded by~\eqref{eq:kappa1}: in the unfolding parameter
$p=|\log\lambda|\log\lambda_2$, the deficit is analytic with
$\delta_1=0.0101990063\,p+O(p^2)$.
\end{remark}

\subsection{Translation gauge and formal conjugacy}\label{sec:onlytwo}

Standard Fatou coordinates on petals with the prescribed asymptotic form
are unique up to additive constants, and their transition maps change by
source and target translations~\cite[Sections 2.2--2.3]{BuffEpstein2002}.
All $\kappa_n$ are invariant under these translations
(Remark~\ref{rem:invariance}).  Arbitrary periodic changes of Abel coordinate
need not preserve those asymptotic conditions; separate changes on the two
sides do not in general act by a simultaneous conjugacy.

A different question arises after setting the transition's constant mode to
zero.  Its cylinder germ is then
\begin{equation}\label{eq:germzeta}
 F(\zeta)=\zeta\exp\Bigl(2\pi\ii\sum_{n\ge1}B_n\zeta^n\Bigr)
 =\zeta+a_2\zeta^2+a_3\zeta^3+\cdots,\quad
 a_2=2\pi\ii B_1,\quad
 a_3=2\pi\ii B_2+\tfrac{(2\pi\ii)^2}{2}B_1^2.
\end{equation}
The multiplier is $F'(0)=1$; $a_2$ is the quadratic coefficient.
By Proposition~\ref{prop:horn-nonzero}, this is a simple parabolic germ, and formal conjugacy tangent to the identity
preserves $a_2$ and the iterative residue $\rho=1-a_3/a_2^2$.
Allowing a nonzero linear change also permits normalising $a_2$ to $1$;
see~\cite[Definition 1 and Appendix]{BuffEpstein2002}.
Substitution gives the exact identity
\begin{equation}\label{eq:kapparho}
 \kappa_2=\frac{B_2}{B_1^2}
 =2\pi\ii\Bigl(\frac12-\rho\Bigr),\qquad
 \rho=0.249455254258832-0.004029881638351\ii.
\end{equation}
If the constant mode $c_0$ is retained, the multiplier instead equals
$\eee^{2\pi\ii c_0}$; the tangent-to-identity formula requires its removal.
Subtracting $c_0$ in the output time coordinate does so without changing
positive Fourier coefficients.

The gauge-check script conjugates $h^{-1}$ simultaneously
by $\psi(z)=z+\sum_{n\ge1}s_n\eee^{2\pi\ii nz}$, which descends to a
tangent-to-identity cylinder change.  With the three tested coefficient lists,
$|c_1|$ and $\kappa_2$ remain unchanged to twelve printed digits, while
$\kappa_3$ becomes $22.9-2.7\ii$, $-26.2+81.9\ii$, and $121.2-63.8\ii$.
This checks invariance under that larger group; its rounded inputs do not
certify the underlying horn coefficients to twelve digits.

It does not follow that higher Fourier comparisons are meaningless.
Once both sides use compatible normalised coordinates,
$\kappa_n(b)\to B_n/B_1^n$ remains a testable proposal for every $n$.
The present $n\ge3$ readings lack reliable numerical resolution; that is
why they are not used as convergence evidence.  Formal coordinate changes
also need not converge, so finite formal classification does not eliminate
the analytic information carried by a horn map.

\begin{remark}[Precision and storage]\label{rem:n2}
An earlier harness stored the $O(1)$ Taylor coefficients of $\Kk$ with
$24$ significant digits.  This produced a floor near $10^{-25}$ in
reconstructed spectra even when the working precision was increased.
The previously reported $|\hat c_2|=0.01203$ at $b=1.20$ was not a reliable
measurement.  The table in Section~\ref{sec:mode2} uses the corrected
storage runs for $b=1.15,1.20,1.25,1.30$; the former warning that these
reruns were pending is obsolete.  Agreement across sampling lines is a
consistency test, not an error bound.  The first-mode values extracted from
separately stored $D$ avoid this particular coefficient-storage floor,
without thereby becoming certified estimates.
\end{remark}

\section{The separation is a transition-map invariant}\label{sec:transition}

Everything so far compares $\Kk_b$, a boundary value of a family built by the
$\theta$-mapping, with $\Rr_b$.  This section removes $\Kk_b$ from the
problem.  The separation coefficients turn out to be determined, up to a
relative error of order $\Lambda$, by the transition map between the two
\emph{regular} solutions at the two real fixed points --- an object built from
two convergent Schr\"oder series in a fraction of a second, at any $b<\eta$.

\subsection{The transition map}
Let $\tau$ be the Koenigs coordinate at the repelling point $L_2$,
$\tau(E(w))=\lambda_2\tau(w)$, $\tau'(L_2)=1$, and put
\[
 S(z)=\tau^{-1}\bigl(-\lambda_2^{\,z}\bigr),
\]
continued to an entire function by $S(z+1)=E(S(z))$.  It maps $\mathbb R$
decreasingly onto the segment $(L,L_2)$ and has period $2\pi\ii/\log\lambda_2$.
Let $\Rr^{-1}(w)=\log(\sigma(w)/\sigma(1))/\log\lambda$ be the attracting Abel
coordinate, $\Rr^{-1}(1)=0$, which is multivalued with period
$\ii h=2\pi\ii/|\log\lambda|$.  On the horizontal strip where $S$ takes values
in the immediate basin of $L$ and a continuous branch is chosen, put
\[
 T(z):=\Rr^{-1}(S(z)),\qquad T(z+1)=T(z)+1,\qquad
 T(z)-z=\sum_{n\in\mathbb Z}t_n\eee^{2\pi\ii nz}.
\]

\begin{lemma}[Real structure]\label{lem:realT}
For real $b\in(1,\eta)$, $\operatorname{Im}T(x)\equiv h/2\pmod h$ for every
real $x$.  Consequently, on the branch with $\operatorname{Im}t_0=h/2$,
$T(x)-x-\ii h/2$ is real, $t_{-n}=\overline{t_n}$, and the line
$\operatorname{Im}w=-h/2$ of attracting time is the image of $\mathbb R$ under
the branch $T-\ii h$.
\end{lemma}

\begin{proof}
$S(x)\in(L,L_2)$, where $\sigma>0$, while $1<L$ gives $\sigma(1)<0$.  Hence
$\sigma(S(x))/\sigma(1)$ is negative and
$\operatorname{Im}\log(\sigma(S(x))/\sigma(1))\in\pi+2\pi\mathbb Z$; divide
by $\log\lambda=-2\pi/h$.
\end{proof}

Translations of $S$ ($S\mapsto S(\cdot+\delta)$, $\delta\in\mathbb C$) send
$t_0\mapsto t_0+\delta$, $t_n\mapsto t_n\eee^{2\pi\ii n\delta}$, so
\begin{equation}\label{eq:taun}
 \tau_n(b):=t_n\,\eee^{-2\pi\ii n t_0}\qquad(\operatorname{Im}t_0=h/2)
\end{equation}
does not depend on the normalisation of $S$; it depends on that of $\Rr$
only through the real translation fixed by $\Rr(0)=1$.  The condition
$\operatorname{Im}t_0=h/2$ refers to the real normalisation of $S$ and fixes the
branch of $\Rr^{-1}$; a branch shift by $\ii kh$ would multiply $\tau_n$ by
$\Lambda^{-nk}$.  The branch is fixed once in the real normalisation and then
carried along.

\subsection{The welding identity}

\begin{theorem}\label{thm:welding}
Let $\Kk$ be a solution of~\eqref{eq:abel}, holomorphic on the sets below
where it is represented, and suppose there are
$1$-periodic holomorphic maps $\varphi_\pm=\mathrm{id}+(\text{periodic})$ with
\begin{enumerate}
\item $\Kk=\Rr\circ\varphi_+$ on an upper half-plane $U_+$, where
  $P=\varphi_+-\mathrm{id}=\sum_{n\ge0}c_n\eee^{2\pi\ii nz}$ converges on $U_+$;
\item $\Kk=S\circ\varphi_-$ on a lower half-plane $U_-$, where $\varphi_-$ is
  univalent, $G:=\varphi_-^{-1}-\mathrm{id}$ is bounded on
  $V_-=\varphi_-(U_-)$, which contains a lower half-plane, and
  $G(w)\to g_0$ as $\operatorname{Im}w\to-\infty$;
\item $U_+\cap U_-$ contains a horizontal line $\ell$ with $\varphi_-(\ell)$
  inside the strip of $T$.
\end{enumerate}
Then, with the branch of $T$ fixed by continuity from $\varphi_+\circ\varphi_-^{-1}$,
for every $n\ge1$
\begin{equation}\label{eq:welding}
 t_n=\sum_{m\ge n}c_m\bigl[\eee^{2\pi\ii mG}\bigr]_{n-m},\qquad
 t_0=g_0+\sum_{m\ge0}c_m\bigl[\eee^{2\pi\ii mG}\bigr]_{-m},
\end{equation}
where $[F]_k$ is the $k$-th Fourier coefficient; in particular
$[\eee^{2\pi\ii mG}]_0=\eee^{2\pi\ii mg_0}$ and $[\eee^{2\pi\ii mG}]_k=0$ for $k>0$.
\end{theorem}

\begin{proof}
On $\ell$, $\Rr\circ\varphi_+=S\circ\varphi_-=\Rr\circ T\circ\varphi_-$, so
$\varphi_+-T\circ\varphi_-$ takes values in the discrete period lattice of
$\Rr$ and is constant; absorb it into the branch of $T$.  Hence
$T=\varphi_+\circ\varphi_-^{-1}$ on the curve $\varphi_-(\ell)$, i.e.\
$T(w)-w=G(w)+P(w+G(w))$ there.  The Fourier coefficients are computed by
Cauchy's theorem on this curve.  The region below it lies in $V_-$, and
$\sum c_m\eee^{2\pi\ii m\varphi_-^{-1}(w)}$ converges uniformly on the curve
$\varphi_-(\ell)$ itself, because $\ell\subset U_+$.  Since $G$ is bounded and holomorphic on $V_-$, with
limit $g_0$ at $-\ii\infty$, so is $\eee^{2\pi\ii mG}$, whose Fourier series
therefore contains only modes $k\le0$ and has constant term
$\eee^{2\pi\ii mg_0}$.  Expanding
$P(w+G)=\sum_mc_m\eee^{2\pi\ii mw}\eee^{2\pi\ii mG(w)}$ and collecting modes
gives~\eqref{eq:welding}; the Fourier coefficients of $T-\mathrm{id}$ do not depend
on the line inside its strip of holomorphy.
\end{proof}

\begin{corollary}\label{cor:cn}
Under the hypotheses of Theorem~\ref{thm:welding}, for every $n\ge1$,
\[
 t_n\,\eee^{-2\pi\ii ng_0}
 =c_n+\eee^{-2\pi\ii ng_0}
   \sum_{m>n}c_m[\eee^{2\pi\ii mG}]_{n-m}.
\]
If $c_n\ne0$, this may be written as
\[
 t_n\,\eee^{-2\pi\ii ng_0}=c_n\bigl(1+\rho'_n\bigr),\qquad
 |\rho'_n|\le\sum_{m>n}\frac{|c_m|}{|c_n|}\,
 \bigl|[\eee^{2\pi\ii mG}]_{n-m}\bigr|\,\bigl|\eee^{-2\pi\ii ng_0}\bigr|,
\]
and $g_0=t_0-c_0-\sum_{m\ge1}c_m[\eee^{2\pi\ii mG}]_{-m}$.  Heuristically, if the
$c_m$ decay like $\Lambda^m$ and $G$ is bounded on a line at height comparable to
$h$ above the gluing line, both corrections are $O(\Lambda)$.  The precise
statement for the sewing solution is Theorem~\ref{thm:W2}(4).  On the
branch of Lemma~\ref{lem:realT}, $g_0=t_0-\ii h+O(\Lambda)$, so that
for the sewing solution, by Theorem~\ref{thm:W2}(4),
\begin{equation}\label{eq:chat-tau}
 \hat c_n:=\frac{c_n}{\Lambda^n}=\tau_n(b)+O(\Lambda).
\end{equation}
Since $B_1\ne0$, also $c_1\ne0$ for $b$ sufficiently close to $\eta$ and
$\kappa_n:=c_n/c_1^{\,n}=t_n/t_1^{\,n}+O(\Lambda)$ for each fixed $n$.
\end{corollary}

The identity $\eee^{-2\pi\ii n(t_0-\ii h)}=\Lambda^n\eee^{-2\pi\ii nt_0}$
is where the scale $\Lambda$ enters: \emph{the $\Lambda$ law is the statement
that the gluing line sits at attracting time $-h/2$, half a period below the
real axis, while the gate carrying the positive modes of $T$ sits a further
half period below it.}  Nothing in the argument uses Kneser's construction.

\subsection{Welding construction}

\begin{proposition}\label{prop:weld}
For $b\in(1,\eta)$ glue the half-cylinders
$C_-=\{\operatorname{Im}w'\le0\}/\mathbb Z$ (repelling time) and
$C_+=\{\operatorname{Im}w\ge-h/2\}/\mathbb Z$ (attracting time) along their
boundaries by $w=T(w')-\ii h$.  The result is conformally $\mathbb C^*$; its
universal-cover coordinate $z$ defines a solution
\[
 \Kk^{W}(z)=\Rr(w(z))\ \text{on }C_+,\qquad \Kk^{W}(z)=S(w'(z))\ \text{on }C_-,
\]
unique up to a translation in $z$, which satisfies the hypotheses of
Theorem~\ref{thm:welding}.  It is holomorphic except where $w(z)$ meets the
singular set $\bigcup_{k}((-\infty,-2]+\ii kh)$ of $\Rr$.
\end{proposition}

\noindent
Unlike $S$, which is entire (the Poincar\'e function of a repelling point),
$\Rr$ is not: its inverse Koenigs map extends only through inverse branches of
$E$.  So $\Kk^W$ carries, besides the cut near $(-\infty,-2]$, the copies
near $(-\infty,-2]+\ii kh$, $k\ge1$, lying in its $\Rr$-represented region.
The same holds for the Kneser-type solutions at complex bases inside the
Shell--Thron region, whose upper representation is through the attracting
regular superfunction.  An earlier draft wrote ``entire'' here, which is
wrong.

\begin{proof}
By Lemma~\ref{lem:realT} the gluing map is a real-analytic diffeomorphism
between the boundary circles that extends holomorphically to a neighbourhood,
so the glued surface $X$ is a Riemann surface; each end is a punctured disc,
hence $X\cup\{\text{two points}\}$ is a simply connected compact Riemann
surface and $X\cong\mathbb C^*$.  The definitions agree on the seam because
$\Rr\circ(T-\ii h)=S$.  The generator of $\pi_1(X)$ lifts to $z\mapsto z+1$
and acts as $w\mapsto w+1$, $w'\mapsto w'+1$, whence
$\Kk^W(z+1)=E(\Kk^W(z))$.  Near each puncture the change of coordinate is a
holomorphic map of punctured discs fixing the puncture, which is the
statement that $w-z$ (resp.\ $w'-z$) has only non-negative (resp.\
non-positive) Fourier modes.  Uniqueness is the uniqueness of the uniformisation.
\end{proof}

\begin{proposition}[Local parameter dependence of the sewn family]
\label{prop:weld-parameter}
For every $b_0\in(1,\eta)$, the normalised geometric sewing of
Proposition~\ref{prop:weld} extends to a holomorphic family over a complex
disc $\Delta_{b_0}$ about $b_0$.  On every common regular domain of the
universal-cover coordinate, $\Kk_b^W(z)$ is jointly holomorphic in $(b,z)$.
These local families agree on overlaps and therefore define the inner sewn
family on an open complex neighbourhood of the entire real interval
$(1,\eta)$.  No uniform width of that neighbourhood is asserted.
\end{proposition}

\begin{proof}
Fix $b_0$.  Both fixed points and their Koenigs coordinates vary
holomorphically near $b_0$.  The compact set
$S_{b_0}([0,1])\subset(L,L_2)$ lies in the attracting basin and avoids its
fixed point.  Thus the chosen branch of
$T_b=\Rr_b^{-1}\circ S_b$ is jointly holomorphic on a parameter disc times
a narrow strip about $[0,1]$; the relation $T_b(w'+1)=T_b(w')+1$
extends it to the whole strip.  The real circle map is orientation
preserving: $\tau_b'(S_b(x))>0$ by inverse iteration to $L_2$,
$\sigma_b'(S_b(x))>0$ by forward iteration to $L$, and hence
$S_b'(x)<0$ and
$(\Rr_b^{-1})'(S_b(x))=\sigma_b'(S_b(x))/
 (\sigma_b(S_b(x))\log\lambda_b)<0$.
Consequently $T_{b_0}'(x)>0$.  Compactness of the circle and local
injectivity give a smaller strip on which $T_{b_0}$ is injective modulo
$1$; shrinking the parameter disc preserves this property and makes the
image of the real circle an analytic Jordan curve in the cylinder.

Put $U_b=T_b+\varpi_b$, where
$\varpi_b=2\pi\ii/\log\lambda_b$; at $b_0$ this is $T_{b_0}-\ii h$.
Use $U_b$ to identify a thin open annular collar in the repelling
cylinder with its image in the attracting cylinder.  More explicitly,
for a sufficiently small $\delta>0$ take the repelling chart
$\{\operatorname{Im}w'<\delta/2\}/\mathbb Z$ and the attracting chart
above the curve $U_b(\{\operatorname{Im}w'=-\delta/2\})$; identify them
on $|\operatorname{Im}w'|<\delta/2$.  After shrinking the parameter disc,
these charts and the transition $(b,w')\mapsto(b,U_b(w'))$ form a
Hausdorff holomorphic family of cylinders.  At $b=b_0$ it is the
analytic-collar version of the boundary welding in
Proposition~\ref{prop:weld}.  Adjoin the two ends using the parameter
independent local coordinates $\eee^{2\pi\ii w}$ and
$\eee^{-2\pi\ii w'}$.  This yields a proper holomorphic submersion
whose fibres are compact Riemann surfaces of genus zero.  The two added
points and the upper-chart point $w=0$ are three disjoint holomorphic
sections (the last remains above the seam after shrinking the disc).

Local holomorphic triviality of a proper family with mutually
biholomorphic compact fibres~\cite{FischerGrauert1965}, followed by the
unique fibrewise M\"obius map sending these three sections to $0,\infty,1$,
gives a jointly holomorphic coordinate $q_b$ on the compactified fibres.
Choose the local lift $z=(2\pi\ii)^{-1}\log q_b$ for which the marked
upper point $w=0$ has $z=0$.  On the lift of the overlap,
$\Rr_b(U_b(w'))=S_b(w')$ because $\varpi_b$ is a period of $\Rr_b$.
The two regular superfunctions therefore define $\Kk_b^W$, jointly
holomorphic wherever their represented branches are regular, and
$\Kk_b^W(0)=\Rr_b(0)=1$.  On the real parameter slice this is exactly
the geometric welding.  Its marked coordinate makes the local families
unique, so the identity theorem patches them on overlapping discs.
\end{proof}


\begin{corollary}[Non-reality at fractional real heights]
\label{cor:sewn-nonreal}
For every real $b<\eta$ sufficiently close to $\eta$, and every
$\delta>0$ sufficiently small, some real $x$ with $0<|x|<\delta$
satisfies $\operatorname{Im}\Kk_b^W(x)\ne0$.  In particular the sewn
solution is not real-valued on any real interval about $z=0$.
\end{corollary}

\begin{proof}
Suppose that $\Kk_b^W(x)$ were real for all real $x$ in a neighbourhood
of $0$.  Shrink that neighbourhood so its values lie in a real interval
about $1$ on which the attracting Koenigs coordinate $\sigma_b$ is
holomorphic and real, with $\sigma_b(\Kk_b^W(x))/\sigma_b(1)>0$.
The upper representation of Theorem~\ref{thm:W2} gives
\[
 \frac{\sigma_b(\Kk_b^W(x))}{\sigma_b(1)}
   =\lambda_b^{\,x+P_b(x)},\qquad P_b(0)=0.
\]
Since $\log\lambda_b$ is real and nonzero, continuity at $x=0$ makes
$\operatorname{Im}P_b(x)=0$ on this interval.  The function $P_b$ is
holomorphic and periodic on a strip containing the whole real line,
so reflection and the identity theorem make it real on that line.
But its Fourier expansion has only non-negative modes.  A real-valued
periodic function has conjugate positive and negative coefficients;
therefore all its nonconstant coefficients vanish.  This contradicts
$c_1(b)\ne0$ for $b$ sufficiently close to $\eta$, established in
Corollary~\ref{cor:cn} using Proposition~\ref{prop:horn-nonzero}.
\end{proof}

This analytic non-reality statement concerns the sewn family near $\eta$.
It is consistent with the tiny non-real part observed numerically by
Paulsen at $b=\sqrt2$~\cite[\S6]{Paulsen2019}, but does not identify his
base-continued family with $\Kk^W$.

\begin{corollary}[Complex cusp profile and its imaginary trace]
\label{cor:imaginary-profile}
Let $A=B_1\eee^{2\pi\ii a}\ne0$ and put $f=F_\eta^{\mathrm{att}}$.
On some fixed complex disk about $z=0$, as $b\to\eta^-$ through real bases,
\[
 \frac{\Kk_b^W(z)-\Rr_b(z)}{\Lambda(b)}
 \longrightarrow f'(z)A(\eee^{2\pi\ii z}-1)
\]
locally uniformly, hence with all $z$-derivatives.  In particular, on
a fixed real interval about $0$,
\[
 \frac{\operatorname{Im}\Kk_b^W(x)}{\Lambda(b)}
 \longrightarrow
 f'(x)
   \operatorname{Im}\!\left[A(\eee^{2\pi\ii x}-1)\right]
\]
uniformly on compact subintervals.  The right-hand side is not
identically zero on any interval about $0$.
\end{corollary}

\begin{proof}
By Theorem~\ref{thm:W2}(4), with
$B=\eee^{2\pi(Y_0+D+2)}$, there is a constant $C$ independent of $b$ and
$n$ such that $|c_n(b)|\le C(\Lambda B)^n$ for all $n\ge1$ and $b$ sufficiently
close to $\eta$.  Here we used $|\tau_n|\le M\eee^{2\pi n(Y_0+1)}$ and
absorbed the factor $\Lambda$ in the error bound into $C$.
Since $P_b(0)=0$, its constant mode is $-\sum_{n\ge1}c_n$.
Choose a disk in the common upper representation supplied by
Corollary~\ref{cor:sewn-cusp-local}.  The Fourier tail is geometric
on every smaller disk once $b$ is close enough to $\eta$, and there
\[
 P_b(z)=c_1(b)(\eee^{2\pi\ii z}-1)+O(\Lambda^2)
\]
locally uniformly.  Theorem~\ref{thm:W2}(4) and
Corollary~\ref{cor:tauconv} give $c_1/\Lambda\to A$.
On a fixed disk about $0$, $\Rr_b\to F_\eta^{\mathrm{att}}$ with its
derivatives by Corollary~\ref{cor:sewn-cusp-local} and its proof.
Taylor's formula gives
$\Kk_b^W(z)=\Rr_b(z)+\Rr_b'(z)P_b(z)+O(\Lambda^2)$
on a smaller disk.  This proves the complex limit, with derivative
convergence by Cauchy's formula.  For real $x$, $\Rr_b(x)$ and $f'(x)$
are real, giving the imaginary trace.
If its last factor vanished on an interval, both the first and second
derivatives at $x=0$ would vanish, forcing $\operatorname{Re}A=
\operatorname{Im}A=0$, contrary to $A\ne0$.
\end{proof}

\begin{corollary}[Obstruction at the real cusp]
\label{cor:cusp-real-obstruction}
For every sufficiently small complex neighbourhood $V$ of $0$, one can
choose a real $x_*\in V\setminus\{0\}$ and constants $c,\delta>0$ such that
\[
 |\operatorname{Im}\Kk_b^W(x_*)|\ge c\Lambda(b)
 \qquad(\eta-\delta<b<\eta,\ b\in\mathbb R).
\]
There is therefore no jointly holomorphic family $F(b,z)$ on a product
of a disc about $b=\eta$ and $V$ which is real at real heights for real
$b>\eta$ and agrees with $\Kk_b^W$ for all real $b<\eta$ sufficiently
close to $\eta$.  In particular, direct holomorphic continuation of a
real-normalised Kneser family through the real cusp cannot identify it
with the inner sewn branch.
\end{corollary}

\begin{proof}
The limiting profile in Corollary~\ref{cor:imaginary-profile} is not
identically zero on any real interval about $0$.  Choose $x_*$ where it
is nonzero.  Since $\Lambda(b)>0$ for real $b<\eta$, the uniform limit
gives the stated lower bound after reducing $\delta$.
For fixed real $x_*$, Schwarz reflection in the parameter says that a
holomorphic function $b\mapsto F(b,x_*)$ which is real on a real interval
above $\eta$ is also real on the real interval below $\eta$ in the same
disc.  This contradicts the lower bound if $F=\Kk^W$ there.
\end{proof}

\begin{corollary}[Conjugate inner branches and their jump]
\label{cor:conjugate-sewn-jump}
For real $b\in(1,\eta)$ put
$\Kk_b^{W,\#}(z)=\overline{\Kk_b^W(\bar z)}$ on the reflected regular
domain.  This reflected family is locally jointly holomorphic in
$(b,z)$ near every such real base.  Both branches obey the tetration
equation, are normalised at $z=0$, agree at every nonnegative integer
height where defined, and have a simple zero at $z=-1$ with the same
logarithmic monodromy $2\pi\ii/\log b$ at $z=-2$.
For $x_*$ as in Corollary~\ref{cor:cusp-real-obstruction},
\[
 |\Kk_b^W(x_*)-\Kk_b^{W,\#}(x_*)|\ge2c\Lambda(b)
 \qquad(\eta-\delta<b<\eta,\ b\in\mathbb R).
\]
More precisely, on a fixed real interval about $0$ the quotient of
their difference by $\Lambda(b)$ tends locally uniformly to
$2\ii(F_\eta^{\mathrm{att}})'(x)
\operatorname{Im}[A(\eee^{2\pi\ii x}-1)]$.
\end{corollary}

\begin{proof}
For complex $b$ near a real base, define the reflected family by
$\Kk_b^{W,\#}(z)=\overline{\Kk_{\bar b}^W(\bar z)}$.
Proposition~\ref{prop:weld-parameter} makes it jointly holomorphic.
Reflection preserves the equation for real $b$ and therefore, by the
identity theorem, on its local complex parameter disc.  The integer
values follow from the equation and normalisation.  The zero at $-1$
and its simplicity follow from Proposition~\ref{prop:backward-log-monodromy},
as does the monodromy.  On real heights the difference is exactly
$2\ii\operatorname{Im}\Kk_b^W(x)$, so the lower bound and the profile
follow from Corollaries~\ref{cor:cusp-real-obstruction} and
\ref{cor:imaginary-profile}.
\end{proof}

\begin{proposition}[The cusp splitting is beyond all algebraic orders]
\label{prop:cusp-no-puiseux}
On a sufficiently small real interval $J$ about $0$, put
\[
 M(x)=(F_\eta^{\mathrm{att}})'(x)
       \operatorname{Im}[A(\eee^{2\pi\ii x}-1)].
\]
The interval $J$ may be chosen so that $M(x)\ne0$ for every
$x\in J\setminus\{0\}$.  For every such $x$
and every integer $q\ge1$, the germ along $t>0$ of
\[
 t\longmapsto\Kk^W_{\eta-t^q}(x)
\]
has no meromorphic extension through $t=0$.  In particular, no finite
ramification of the base parameter makes this value meromorphic at
the real cusp; it has no convergent Puiseux--Laurent germ there.
If it admits a single-valued holomorphic continuation on a full
punctured $t$-disc, the singularity at $t=0$ is essential.
\end{proposition}

\begin{proof}
The profile $M$ is real analytic, vanishes at $0$, and, by
Corollary~\ref{cor:imaginary-profile}, is not identically zero on any
interval about $0$.  Thus $0$ is an isolated zero, and we shrink $J$
accordingly.
Suppose such a meromorphic germ $G(t)$ existed.  Its reflection
$G^\#(t)=\overline{G(\bar t)}$ would also be meromorphic, so
$D(t)=G(t)-G^\#(t)$ would be meromorphic near $0$.
For positive real $t$, the imaginary profile and the saddle-node
scale of $\Lambda$ give
\[
 D(t)=2\ii\operatorname{Im}\Kk^W_{\eta-t^q}(x)
 =O\!\left(\exp(-c_0t^{-q/2})\right)=o(t^N)
 \qquad\text{for every }N\ge0
\]
with some $c_0>0$.  A meromorphic Laurent germ with this decay on the
positive axis is identically zero.  This contradicts the nonzero
limit $D(t)/\Lambda(\eta-t^q)\to2\ii M(x)$.
\end{proof}

\begin{corollary}[A two-jet witness at the normalisation height]
\label{cor:cusp-two-jet}
Put $f=F_\eta^{\mathrm{att}}$ and $A=B_1\eee^{2\pi\ii a}\ne0$.
As $b\to\eta^-$ through real bases,
\[
 \frac{1}{\Lambda(b)}
 \begin{pmatrix}
   \operatorname{Im}[(\Kk_b^W)'(0)]\\
   \operatorname{Im}[(\Kk_b^W)''(0)]
 \end{pmatrix}
 \longrightarrow
 \begin{pmatrix}
   2\pi f'(0)\operatorname{Re}A\\
   4\pi f''(0)\operatorname{Re}A
     -4\pi^2 f'(0)\operatorname{Im}A
 \end{pmatrix}\ne\begin{pmatrix}0\\0\end{pmatrix}.
\]
Consequently at least one of the first two $z$-derivatives at $0$,
with the same choice of derivative for every $q\ge1$, has no
meromorphic base germ through $t=0$ after $b=\eta-t^q$.
The corresponding derivative of the reflected inner branch differs
from that of $\Kk_b^W$ by at least $c\Lambda(b)$ in modulus for some
$c>0$ and all sufficiently close real $b<\eta$.
\end{corollary}

\begin{proof}
The complex limit of Corollary~\ref{cor:imaginary-profile} and
Cauchy's formula give convergence of the first two derivatives.
The derivatives of $\Rr_b$ at $0$ are real for real $b$, giving the
displayed vector.  Since $f'(0)\ne0$ by
Corollary~\ref{cor:sewn-cusp-local}, vanishing of its first entry
forces $\operatorname{Re}A=0$, and then vanishing of its second forces
$\operatorname{Im}A=0$, a contradiction.
For a nonzero entry, the reflection argument in the proof of
Proposition~\ref{prop:cusp-no-puiseux} applies to that derivative.
Reflection doubles its imaginary part, which also gives the lower bound.
\end{proof}

\begin{remark}[Conjugate bypasses of the cusp]
Suppose the real-normalised Kneser germ admits analytic continuation
along two conjugate parameter paths bypassing $\eta$ and reaching a
common real interval below $\eta$.  If the continuation along one path
agrees there with $\Kk^W$, uniqueness of analytic continuation and
Schwarz reflection make the continuation along the reflected path
$\overline{\Kk^W_{\bar b}(\bar z)}$.  At the fixed real height $x_*$ of
Corollary~\ref{cor:cusp-real-obstruction}, the two values differ in
modulus by at least $2c\Lambda(b)$.  The existence of such
continuations is not asserted here.
\end{remark}

Two remarks place Proposition~\ref{prop:weld} relative to $\Kk_b$.  For
non-real $b$ with $\operatorname{Im}b>0$ the same sewing (along the curve
$T_b(\mathbb R)$) yields a solution tending to $L$ and $L_2$ as
$\operatorname{Im}z\to\pm\infty$.  The geometric sewn family depends
holomorphically on $b$ near every real point of $(1,\eta)$
(Proposition~\ref{prop:weld-parameter}).  These two limits alone do not
characterise the solution (Remark~\ref{rem:paulsen}); identifying it with
the base-continued $\Kk_b$ requires the seam condition of
Theorem~\ref{thm:char}.  The numerical comparison in Table~\ref{tab:tau}
shows that $\tau_1$ and $t_2/t_1^2$ reproduce the ladder
values of $\hat c_1$ and $\kappa_2$ --- moduli \emph{and} arguments --- to the
precision of the ladder extrapolation, e.g.\ eight digits at $b=1.10$.

\begin{table}[htbp]
\centering
\caption{Kneser ladders (Sections~\ref{sec:invariant},~\ref{sec:mode2}, extrapolated
to $\operatorname{Im}b=0$) against the Kneser-free invariants
$\tau_1=t_1\eee^{-2\pi\ii t_0}$ and $t_2/t_1^2$ (script \texttt{transition\_map.py}).}
\label{tab:tau}
\small\setlength{\tabcolsep}{3.5pt}
\begin{tabular}{lllll}
\toprule
$b$ & $\hat c_1$ (ladder) & $\tau_1$ & $\kappa_2$ (ladder) & $t_2/t_1^2$\\
\midrule
$1.05$ & $0.0850235\,\eee^{1.534136\ii}$ & $0.0850226\,\eee^{1.534165\ii}$ & $1.374049\,\eee^{1.413420\ii}$ & $1.374336\,\eee^{1.413645\ii}$\\
$1.10$ & $0.08642720\,\eee^{1.55109214\ii}$ & $0.08642720\,\eee^{1.55109214\ii}$ & $1.440851\,\eee^{1.477690\ii}$ & $1.440851\,\eee^{1.477690\ii}$\\
$1.15$ & $0.0872053\,\eee^{1.561012\ii}$ & $0.0872052\,\eee^{1.561013\ii}$ & $1.479111\,\eee^{1.511255\ii}$ & $1.479107\,\eee^{1.511258\ii}$\\
$1.30$ & $0.088434$ & $0.0884321\,\eee^{1.577580\ii}$ & $1.541516\,\eee^{1.561918\ii}$ & $1.541515\,\eee^{1.561919\ii}$\\
\bottomrule
\end{tabular}
\end{table}

\subsection{Reduction of the conjectures to confluence}

Write $a:=\Phi_{\mathrm{att}}(1/\eee-1)=3.0292972144180360989\ldots$ for the
$\alpha$-normalised attracting Fatou coordinate of the base point $w=1$ at
$b=\eta$ (real, since the orbit of $u=1/\eee-1$ stays on the negative axis).

\begin{theorem}\label{thm:reduction}
Assume the following confluence statement \textup{(C)}: there are constants
$s(b)\in\mathbb C$ such that, as $b\to\eta^-$,
\[
 \Rr_b^{-1}\bigl(\eee(1+u)\bigr)\to\Phi_{\mathrm{att}}(u)-a,\qquad
 S_b^{-1}\bigl(\eee(1+u)\bigr)+s(b)\to\Phi_{\mathrm{rep}}(u),
\]
uniformly on compact subsets of the upper gate, with the branches continued as
in Section~\ref{sec:horn}.  Then for every $n\ge1$
\[
 \tau_n(b)\longrightarrow B_n\,\eee^{2\pi\ii na}.
\]
Consequently, also
$t_n/t_1^{\,n}\to B_n/B_1^{\,n}$ for every $n\ge1$.
Together with Corollary~\ref{cor:cn}, and subject to the identification of
$\Kk_b$ discussed in Sections~\ref{sec:W1}--\ref{sec:unif}, this gives
$\hat c_n\to B_n\eee^{2\pi\ii na}$ for all $n$.  The limiting phase
$\arg\hat c_1\to\arg B_1+2\pi a=1.5865330757\ldots\pmod{2\pi}$
and the ratio limit of Conjecture~\ref{conj:horn2} follow for the sewing solution.
\end{theorem}

\begin{proof}
Replace $S_b$ by $S_b(\cdot-s(b))$, which leaves $\tau_n$ unchanged.  By (C),
$T_b\to\Phi_{\mathrm{att}}\circ\Phi_{\mathrm{rep}}^{-1}-a=h^{-1}-a$ uniformly
on $[0,1]+\ii Y_g$ for fixed $Y_g$ in the gate, so the Fourier coefficients on
that line converge: $t_n\to B_n$ ($n\ge1$) and $t_0\to-a$, the latter because
$B_0=0$ in the $\alpha$ normalisation.  Hence
$\tau_n=t_n\eee^{-2\pi\ii nt_0}\to B_n\eee^{2\pi\ii na}$.
\end{proof}

Statement (C) is a confluence theorem for the saddle-node unfolding
$u\mapsto \eee^{\lambda\eee^{-\lambda}\eee(1+u)}/\eee-1$ in the direction where the
two fixed points are real (the ``Glutsyuk'' direction): the Koenigs
coordinates at the two hyperbolic points converge, on the gate, to the
sectorial Fatou coordinates of the parabolic germ.  Results of this type are
due to Glutsyuk~\cite{Glutsyuk2001} and, for the full unfolded modulus,
Marde\v{s}i\'c--Roussarie--Rousseau~\cite{MRR2004} and
Christopher--Rousseau~\cite{ChristopherRousseau}.  Rather than matching their
normalisations, we prove (C) directly for this family in
Section~\ref{sec:proofC}, so Theorem~\ref{thm:reduction} holds unconditionally.

\subsection{Proof of the confluence statement (C)}\label{sec:proofC}

We prove (C) directly for the family at hand; no normalisation from the
literature is needed.  In $u$-coordinates ($w=\eee(1+u)$, $\mu=\eee\log b$)
\[
 f_b(u)=\eee^{\mu-1+\mu u}-1,\qquad f_\eta(u)=\eee^u-1 .
\]
For $b<\eta$ close to $\eta$ the map has, in a fixed disc $D_r$, exactly the two
fixed points $u_1=\eee^{\lambda-1}-1<0<u_2=\eee^{\lambda_2-1}-1$ (Rouch\'e), and
\begin{equation}\label{eq:prepared}
 f_b(u)-u=(u-u_1)(u-u_2)\,k_b(u),\qquad k_b\to k_\eta,\quad
 k_\eta(u)=\tfrac12+\tfrac u6+\cdots,
\end{equation}
uniformly on $D_r$, with $k_b$ zero-free there.  Put $A_j=1/\log\lambda_j$,
so $A_1<0<A_2$.

\begin{lemma}[Residues]\label{lem:res}
As $b\to\eta^-$: $A_1(u_1-u_2)\to2$, $A_1+A_2\to\nu:=1-4k_\eta'(0)=\tfrac13$, hence
$A_1u_1+A_2u_2\to2$ and $A_ju_j^2\to0$.
\end{lemma}

\begin{proof}
$\lambda_1=1+(u_1-u_2)k_b(u_1)$ and $\lambda_2=1+(u_2-u_1)k_b(u_2)$.  With
$d=u_1-u_2\to0$ and $1/\log(1+x)=1/x+\tfrac12+O(x)$,
$A_1+A_2=\frac{k_b(u_2)-k_b(u_1)}{d\,k_b(u_1)k_b(u_2)}+1+O(d)\to1-4k_\eta'(0)$,
and $A_1d\to1/k_\eta(0)=2$.
\end{proof}

\noindent
Numerically $A_1+A_2=0.3333331469$ and $A_1(u_1-u_2)=1.99676$ at $\lambda=0.99$.
Define the \emph{model time}
\[
 \Psi_b(u)=A_1\log(u-u_1)+A_2\log(u-u_2),
\]
which is exactly the time of the vector field $(u-u_1)(u-u_2)/(\ldots)$ whose
time-one map has multipliers $\lambda_1,\lambda_2$, and let
$\alpha_2(u)=-2/u+\tfrac13\log(-u)$.  By Lemma~\ref{lem:res}, for $u,p$ in a
simply connected region at positive distance from $0$, with logarithms
continued along the same path,
\begin{equation}\label{eq:psilimit}
 \Psi_b(u)-\Psi_b(p)\longrightarrow\alpha_2(u)-\alpha_2(p).
\end{equation}

\begin{lemma}[Uniform one-step error]\label{lem:F}
$F_b:=\Psi_b\circ f_b-\Psi_b-1$ is single-valued and holomorphic on a fixed disc
$D_{r'}$, $F_b\to F_\eta:=\alpha_2\circ f_\eta-\alpha_2-1$ uniformly there, and
\[
 |F_b(u)|\le M\,|u-u_1|\,|u-u_2|\qquad(u\in D_{r'/2})
\]
with $M$ independent of $b$.
\end{lemma}

\begin{proof}
By~\eqref{eq:prepared}, $(f(u)-u_1)/(u-u_1)=1+(u-u_2)k(u)$ and
$(f(u)-u_2)/(u-u_2)=1+(u-u_1)k(u)$, so
\[
 F=A_1\log\frac{1+(u-u_2)k}{1+(u-u_1)k}+(A_1+A_2)\log\bigl(1+(u-u_1)k\bigr)-1,
\]
holomorphic for $|u|<r'$ small.  The first logarithm is $O(|u_1-u_2|)$
uniformly, and $A_1(u_1-u_2)$, $A_1+A_2$ are bounded (Lemma~\ref{lem:res}); so
$|F|\le M_0$ on $|u|=r'$ uniformly, and the displayed formula converges.  At
$u=u_1$ the first term is $A_1\log\lambda_1=1$ and the second vanishes, so
$F(u_1)=0$; likewise $F(u_2)=0$.  The maximum principle applied to
$F/((u-u_1)(u-u_2))$ on $D_{r'}$ gives the bound, and letting $b\to\eta$ gives
$F_\eta=O(u^2)$.
\end{proof}

\begin{lemma}[Uniform petals and summability]\label{lem:petal}
Put $Z_b(u)=A_1\log\frac{u-u_1}{u-u_2}$ and
$\Omega_{R,b}=\{u:\operatorname{Re}Z_b(u)>R\}$ (a disc of Apollonius about
$u_1$).  There is $R_0$ such that for $R\ge R_0$ and all $b$ near $\eta$:
$\Omega_{R,b}\subset D_{r'/2}$, $f_b(\Omega_{R,b})\subset\Omega_{R,b}$,
$\operatorname{Re}Z_b(f_b^ku)\ge\operatorname{Re}Z_b(u)+\tfrac34k$, and
\[
 \sum_{k\ge N}|f_b^ku-u_1|\,|f_b^ku-u_2|\le\frac{C}{R+\tfrac34N}
 \qquad(u\in\Omega_{R,b}),
\]
with $C$ independent of $b$ and $u$.  As $b\to\eta$, $\Omega_{R,b}$ converges on
compacta to the petal $\{\operatorname{Re}(-2/u)>R\}$.
\end{lemma}

\begin{proof}
$Z(f(u))-Z(u)=A_1\log\frac{1+(u-u_2)k}{1+(u-u_1)k}=1+O(|u|+|u_1|+|u_2|)$
uniformly, by the expansion in the proof of Lemma~\ref{lem:F}; on
$\Omega_{R,b}\subset D_{C/R+|u_1|}$ this is within $\tfrac14$ of $1$ once $R$ is large,
which gives invariance and the growth of $\operatorname{Re}Z$.  With
$q=(u-u_1)/(u-u_2)$, $|q|=\eee^{-\sigma}$, $\sigma=|\log\lambda_1|\operatorname{Re}Z$,
one has $u-u_2=(u_2-u_1)/(q-1)$, whence
\[
 |u-u_1||u-u_2|=\frac{|q|\,|u_2-u_1|^2}{|1-q|^2}
 \le\Bigl(\frac{|u_2-u_1|}{|\log\lambda_1|}\Bigr)^2
 \frac{(1+\sigma)^2\eee^{-\sigma}}{(\operatorname{Re}Z)^2},
\]
using $1-\eee^{-\sigma}\ge\sigma/(1+\sigma)$.  The first factor is bounded
(Lemma~\ref{lem:res}) and $(1+\sigma)^2\eee^{-\sigma}\le4$; summing
$(R+\tfrac34k)^{-2}$ over $k\ge N$ gives the bound.  The last assertion follows
from $\operatorname{Re}Z_b(u)\to\operatorname{Re}(-2/u)$, i.e.\
$A_1(u_2-u_1)\to-2$.
\end{proof}

\begin{theorem}[Confluence]\label{thm:C}
Statement \textup{(C)} holds.  More precisely, with branches continued along
paths in the respective parabolic basins:
\begin{enumerate}
\item $\Rr_b^{-1}(\eee(1+u))\to\Phi_{\mathrm{att}}(u)-a$ uniformly on compact
  subsets of the parabolic basin of $f_\eta$;
\item for a fixed $p'$ in the gate (below), putting
  $s(b)=\Phi_{\mathrm{rep}}(p')-\mathcal S_b(p')$,
  $\mathcal S_b(u)+s(b)\to\Phi_{\mathrm{rep}}(u)$ uniformly on compact
  subsets of the repelling basin (backward orbits tending to $0$).  Here
  $\mathcal S_b$ is the real-normalised repelling time, i.e.\ the branch of
  $S_b^{-1}$ that is real on the segment $(u_1,u_2)$ and is continued through
  the band of Lemma~\ref{lem:band}, then along backward-invariant paths.
\end{enumerate}
The \emph{gate} of height $Y_*$ is $\mathcal G(Y_*)=\Phi_{\mathrm{rep}}^{-1}
(\{Y_*\le\operatorname{Im}z\le Y_*+1,\ -1\le\operatorname{Re}z\le2\})$.  For $Y_*$ large it is a compact
subset of both parabolic basins of $f_\eta$.  For $b$ near $\eta$ it lies in
the band $\mathcal B_b(Y_*-2)$, because $Z_b\to-2/u$ uniformly on it
(Lemma~\ref{lem:res}) and $\operatorname{Im}(-2/u)=\operatorname{Im}\alpha(u)+O(1)
=\operatorname{Im}\Phi_{\mathrm{rep}}(u)+O(1)$ there.
\end{theorem}

\begin{proof}
(1) On $\Omega_{R,b}$ put
$\mathcal A_b(u)=\lim_n[\Psi_b(f_b^nu)-n]=\Psi_b(u)+\sum_{k\ge0}F_b(f_b^ku)$;
the series converges absolutely by Lemmas~\ref{lem:F} and~\ref{lem:petal}.
Since $f^nu\to u_1$ and the Koenigs map satisfies
$\sigma(f^nu)=\lambda_1^n\sigma(u)$,
$\Psi_b(f^nu)-n=A_1\log\sigma(u)+A_2\log(u_1-u_2)+o(1)$, because
$A_1\log\lambda_1=1$.  Thus $\mathcal A_b$ is the Koenigs Abel coordinate up to
an additive constant, and
$\Rr_b^{-1}(\eee(1+u))=\mathcal A_b(u)-\mathcal A_b(u^*)$, $u^*=1/\eee-1$.
Let $K$ be a compact subset of the petal
$P_R=\{\operatorname{Re}(-2/u)>R\}$ and $U\supset K\cup\{f_\eta^{N_0}u^*\}$ a
simply connected region in $P_R$ at positive distance from $0$.  For $b$ near
$\eta$, $U\subset\Omega_{R,b}$ and $u_1,u_2\notin U$.  For $u,p\in U$,
\[
 \mathcal A_b(u)-\mathcal A_b(p)=\Psi_b(u)-\Psi_b(p)
 +\sum_{k\ge0}\bigl[F_b(f_b^ku)-F_b(f_b^kp)\bigr].
\]
The first term converges by~\eqref{eq:psilimit}.  In the sum, each term
converges because $f_b^k\to f_\eta^k$ and $F_b\to F_\eta$ uniformly.  The tails
beyond $N$ are $\le 2MC/(R+\tfrac34N)$ uniformly in $b$ (Lemma~\ref{lem:petal}).
Hence the difference converges uniformly to
$\alpha_2(u)-\alpha_2(p)+\sum_k[F_\eta(f_\eta^ku)-F_\eta(f_\eta^kp)]
 =\Phi_{\mathrm{att}}(u)-\Phi_{\mathrm{att}}(p)$.  The last identity holds
because $\Phi_{\mathrm{att}}=\lim[\alpha(f^nu)-n]$ and
$\alpha-\alpha_2=O(u)\to0$ along orbits.  A general compact $K$ of the basin
satisfies $f_\eta^{N}(K)\Subset P_R$ for some $N$.  Then
$f_b^N(K)\Subset P_R$ for $b$ near $\eta$, and
$\mathcal A_b=\mathcal A_b\circ f_b^N-N$.  Taking $p=f_b^{N_0}u^*$, which
tends to $f_\eta^{N_0}u^*\in U$, and using uniform convergence on a
neighbourhood gives (1).  (On $\Omega_{R,b}$ the function $\mathcal A_b$ is
multivalued, since this punctured Apollonius disc surrounds $u_1$.
It is single-valued on $U$, which
avoids $u_1$.  On $K$, $\sigma_b\ne0$ because $f_b^N(K)$ avoids $u_1$.  On
non-immediate components of the basin the branch is defined by
$\mathcal A_b\circ f_b^N-N$.  The Koenigs map is normalised by
$\sigma_b'(u_1)=1$.)

(2) Apply (1)'s argument to $G_b(v)=-f_b^{-1}(-v)$, the principal inverse
branch reflected.  Then $G_\eta(v)=-\log(1-v)=v+\tfrac{v^2}2+\tfrac{v^3}3+\cdots$,
which is of the same prepared form with $k_\eta'(0)=\tfrac13$ and residue
$-\tfrac13$.  The attracting fixed point of $G_b$ is $-u_2$, with multiplier
$1/\lambda_2$.  Its Koenigs Abel coordinate is $-S_b^{-1}(\eee(1-v))$ up to a
constant, and its parabolic Fatou coordinate is $-\Phi_{\mathrm{rep}}(-v)$ up
to a constant (since $\alpha_{f}(-v)=-\alpha_{G}(v)+\text{const}$).  The free
constant is fixed by $p'$.
\end{proof}

With Theorem~\ref{thm:C}, Theorem~\ref{thm:reduction} becomes unconditional.
Four points need checking; each is settled here.  The proof of
Lemma~\ref{lem:band} uses only Lemmas~\ref{lem:res}--\ref{lem:petal}, so there is
no circularity.

\emph{(a) Inversion on the gate line.}  Fix $Y_*\ge Y_B+6$ and the line
$\ell=[0,1]+\ii Y_*$.  By Rouch\'e, the solution $v_b(z)$ of
$\mathcal S_b(v)+s(b)=z$ near $\Phi_{\mathrm{rep}}^{-1}(z)$ exists and converges
uniformly on $\ell$.  Moreover $S_b(z-s(b))=\eee(1+v_b(z))$, because $\mathcal S_b$
is a branch of $S_b^{-1}$.

\emph{(b) The gate line lies in the strip of holomorphy of $T$.}  The points
$v_b(z)$, $z\in\ell$, lie within $o(1)$ of the gate, hence in the band $\mathcal B_b(Y_B+3)$.
By Lemma~\ref{lem:band}(ii) their $S$-times $z-s(b)$ satisfy
$|\operatorname{Im}(z-s(b))|\le h/2-Y_B-3+\pi/3+C/Y_B+C\varepsilon<h/2-Y_B-2$.  So the
shifted line $\ell-s(b)$ lies in the strip $\Sigma$ of Lemma~\ref{lem:band}(iii),
where $T$ is holomorphic.  Therefore the Fourier coefficients of
$\tilde T=T(\cdot-s)$ on $\ell$ are $t_n\eee^{-2\pi\ii ns}$ and
$\tilde t_0=t_0-s$, with the $t_n$ those of $T$ on $\mathbb R$.  No prior
knowledge of $\operatorname{Im}s$ is used.

\emph{(c) The branch of $\Rr_b^{-1}$.}  Continue $\Rr_b^{-1}$ from $u^*$ along
the negative axis to $f_\eta^{N_0}u^*$, where it is real.  Then go inside
$P_R\setminus D_\delta(0)$ through the upper half-plane to the gate, and then
inside the band down to the segment $(u_1,u_2)$.  This path passes above $u_1$
and does not wind around $u_2$.  Along it $\Rr_b^{-1}=\Psi_b+\sum_kF_b\circ f_b^k+c_R$, so the change of
$\operatorname{Im}\Rr_b^{-1}$ is
$A_1\Delta\arg(u-u_1)+A_2\Delta\arg(u-u_2)+O(C/Y+\varepsilon)=h/2+0+\text{small}$.
Lemma~\ref{lem:realT} forces the exact value, so on the segment
$\operatorname{Im}\Rr_b^{-1}=+h/2$.  The branch in
Theorem~\ref{thm:C}(1) is therefore the one with $\operatorname{Im}t_0=h/2$ used
in~\eqref{eq:taun}.

\emph{(d) $B_0=0$.}  At $b=\eta$, on the gate,
$\Phi_{\mathrm{att}}-\alpha_2=\sum_{k\ge0}F_\eta(f^ku)$ and
$\Phi_{\mathrm{rep}}-\alpha_2=-\sum_{k\ge1}F_\eta(f^{-k}u)$, with the same
principal $\log(-u)$ in $\alpha_2$.  That logarithm stays continuous along both
orbits, since $-f^{\pm k}u$ remains in the lower half-plane.  By
Lemma~\ref{lem:F} and $f^{\pm k}u\approx-2/(x\pm k+\ii Y)$, both sums are
$O(1/Y)$.  Hence $\Phi_{\mathrm{att}}\circ\Phi_{\mathrm{rep}}^{-1}(z)-z\to0$ as
$\operatorname{Im}z\to+\infty$, so its constant Fourier mode vanishes:
$B_0=0$.  Numerically $|B_0|<10^{-59}$.

By (a)--(d) and Theorem~\ref{thm:C}, $\tilde T\to h^{-1}-a$ uniformly on
$\ell$, where $h^{-1}$ denotes the inverse horn map (not the period $h$).
So $\tilde t_n\to B_n$ and $\tilde t_0\to-a$, and
$\tau_n=\tilde t_n\eee^{-2\pi\ii n\tilde t_0}$ gives:

\begin{corollary}\label{cor:tauconv}
For every $n\ge1$, $\tau_n(b)\to B_n\eee^{2\pi\ii na}$ as $b\to\eta^-$.
Since $B_1\ne0$, also $t_n/t_1^n\to B_n/B_1^n$ for every $n\ge1$.
\end{corollary}

\noindent
A direct pointwise check of Theorem~\ref{thm:C} on two gate points
(\texttt{confluence\_check.py}) gives errors
$0.086,\,0.021,\,0.0050,\,0.00079$ (attracting side) at
$\lambda=0.8,0.9,0.95,0.98$, decreasing like $\varepsilon^2$.

\subsection{The Kneser side I: proof of the decay estimate (W2)}\label{sec:W2}

Throughout, $b<\eta$ is real and close to $\eta$, $\varepsilon=|\log\lambda|$,
and $S$ has the real normalisation of Section~\ref{sec:transition}, so
$T(z)-z=\sum t_n\eee^{2\pi\ii nz}$ with $\operatorname{Im}t_0=h/2$ and
$t_{-n}=\overline{t_n}$ (Lemma~\ref{lem:realT}).

\begin{lemma}[The band]\label{lem:band}
Let $Z_b(u)=A_1\log\frac{u-u_1}{u-u_2}$ with the branch
$\operatorname{Im}Z_b\in(0,h)$ on $\mathbb C\setminus((-\infty,u_1]\cup[u_2,\infty))$.
This branch has $\operatorname{Im}Z_b=h/2$ on $(u_1,u_2)$ and
$\operatorname{Im}Z_b<h/2$ on $\operatorname{Im}u>0$.  For $Y>0$ let
\[
 \mathcal B_b(Y)=\{u:\ Y<\operatorname{Im}Z_b(u)<h-Y\},
\]
a lens bounded by two circular arcs through $u_1,u_2$.  There is an absolute
$Y_B$ such that for $Y\ge Y_B$ and all $b$ near $\eta$:
\begin{enumerate}
\item[(i)] every $u\in\mathcal B_b(Y)$ has its forward orbit in $D_{r'}$ tending
  to $u_1$ and its principal backward orbit in $D_{r'}$ tending to $u_2$, with
  $\sum_{k\in\mathbb Z}|F_b(f_b^ku)|\le C/Y+C\varepsilon$;
\item[(ii)] the real-normalised times $\mathcal R_b=\Psi_b+\sum_{k\ge0}F_b\circ f_b^k+c_R$
  and $\mathcal S_b=\Psi_b-\sum_{k\ge1}F_b\circ f_b^{-k}+c_S$ are holomorphic and
  injective on $\mathcal B_b(Y)$.  They are branches of $\Rr_b^{-1}$ and
  $S_b^{-1}$, with $\mathcal S_b$ real on $(u_1,u_2)$ and $\mathcal R_b$ as in
  item (c) above, and
  \[
   \bigl|\operatorname{Im}\mathcal S_b(u)-(\operatorname{Im}Z_b(u)-h/2)\bigr|
   \le\tfrac\pi3+C/Y+C\varepsilon ;
  \]
\item[(iii)] $S_b$ maps the strip $\Sigma=\{|\operatorname{Im}w'|<h/2-Y-2\}$
  into $\mathcal B_b(Y)$ with $\mathcal S_b\circ S_b=\mathrm{id}$.  Consequently
  $T=\mathcal R_b\circ S_b$ is holomorphic on $\Sigma$, and
  \[
   T(w')-w'=(c_R-c_S)+\sum_{k\in\mathbb Z}F_b\bigl(f_b^k(S_b(w'))\bigr).
  \]
\end{enumerate}
\end{lemma}

\begin{proof}
(i) Put $q=\eee^{Z/A_1}=(u-u_1)/(u-u_2)$ and $\zeta=Z/A_1=-\varepsilon Z$, so
$\operatorname{Im}\zeta\in(-2\pi+\varepsilon Y,-\varepsilon Y)$ on the band.  On
$-2\pi\le\operatorname{Im}\zeta\le0$ one has
$|\eee^\zeta-1|\ge c\min(1,\delta(\zeta))\max(1,|\eee^\zeta|)$ with
$\delta(\zeta)=\min(|\zeta|,|\zeta+2\pi\ii|)$ and an absolute $c>0$; the only
zeros are $0$ and $-2\pi\ii$.  On the band
$\delta(\zeta)=\varepsilon\min(|Z|,|Z-\ii h|)>\varepsilon Y$.  Since
$u-u_2=(u_2-u_1)/(q-1)$ and $|u_2-u_1|/\varepsilon\to2$
(Lemma~\ref{lem:res}), this gives
\begin{equation}\label{eq:bandsize}
 |u-u_2|\le C\bigl(\varepsilon+Y^{-1}\bigr),\qquad
 |u-u_1||u-u_2|=\frac{|q|\,|u_2-u_1|^2}{|q-1|^2}\le
 \begin{cases}C\,\min(|Z|,|Z-\ii h|)^{-2},&\delta(\zeta)\le1,\\
 C\varepsilon^2\eee^{-\varepsilon|\operatorname{Re}Z|},&\delta(\zeta)>1,\end{cases}
\end{equation}
using $|q|/\max(1,|q|)^2=\eee^{-|\operatorname{Re}\zeta|}$.  So
$\mathcal B_b(Y)\subset D_{C/Y+|u_1|+|u_2|}\subset D_{r'/2}$ for $Y$ large.  By the
proof of Lemma~\ref{lem:petal}, $Z\circ f-Z=1+O(|u|+\varepsilon)$, so
$\operatorname{Re}Z$ increases by at least $\tfrac34$ per step while the orbit is
in $D_{r'/2}$.

The imaginary part is controlled through
$\Psi_b=Z_b+(A_1+A_2)\log(u-u_2)$, with step error
$|F_b|\le M|u-u_1||u-u_2|$ (Lemma~\ref{lem:F}).  Summing~\eqref{eq:bandsize}
along $\operatorname{Re}Z_k\ge\operatorname{Re}Z_0+\tfrac34k$ bounds
$\sum_k|F_b(u_k)|$ by $C'/Y+C'\varepsilon$.  Points with
$\delta(\zeta_k)\le1$ contribute $\sum_k(x_k^2+Y^2)^{-1}\le C/Y$, and the others
$\varepsilon^2\sum\eee^{-\varepsilon|x_k|}\le C\varepsilon$.  Now
$\operatorname{Im}Z-\operatorname{Im}\Psi=-(A_1+A_2)\arg(u-u_2)$ with
$\arg(u-u_2)\in(0,2\pi)$ on the lens.  So $\operatorname{Im}Z$ drifts by at most
$C'/Y+C'\varepsilon+\tfrac{2\pi}3+o(1)$ along the orbit, which is $<Y/2$ for
$Y\ge Y_B\ge8$.  By induction the orbit stays in
$\mathcal B_b(Y/2)\subset D_{r'/2}$, and $\operatorname{Re}Z_k\to+\infty$,
i.e.\ $u_k\to u_1$.  The continued values of $Z$ along the orbit have
$\operatorname{Im}\in(Y/2,h-Y/2)\subset(0,h)$, so they agree with the fixed
slit-plane branch.  The backward orbit is the forward orbit of the reflected
inverse family $G_b(v)=-f_b^{-1}(-v)$, to which the same argument applies
(its prepared form is checked in the proof of Theorem~\ref{thm:C}(2)).  Its
period is $h_2=h+2\pi/3+o(1)$, so $\mathcal B_b(Y)$ is contained in the
corresponding band of $G_b$ at height $Y(1+O(\varepsilon))-1$.  This is harmless
for $Y\ge Y_B$.

(ii) As in the proof of Theorem~\ref{thm:C}(1), $\Psi_b(f^nu)-n\to A_1\log\sigma_b(u)+\mathrm{const}$
along forward orbits.  So $\mathcal R_b$ is a branch of the attracting Abel
coordinate, and $\mathcal S_b$ likewise of the repelling one.  Choose the
constants $c_R$ (by item (c) above) and $c_S$ so that $\mathcal S_b$ is real on
the segment.  Both $F_b$ and $f_b$ are real, so the sums are real there.  On
the segment $\operatorname{Im}Z=h/2$ and $\arg(u-u_2)=\pi$, which fixes
$\operatorname{Im}c_S$.  The displayed estimate then follows from
$\operatorname{Im}\Psi=\operatorname{Im}Z+(A_1+A_2)\arg(u-u_2)$,
$A_1+A_2\to\tfrac13$ and (i).  Injectivity: if $\mathcal S_b(u)=\mathcal S_b(u')$
then the principal backward orbits satisfy
$\tau_b(u_{-n})=\tau_b(u'_{-n})$ near $u_2$, where $\tau_b$ is injective, so
$u_{-n}=u'_{-n}$.  Moreover $f_b$ is injective on $D_{r'}$ (its non-injectivity
has period $2\pi\ii/\log b$ in $w$, far larger than the disc), so $u=u'$.  The
same argument with $\sigma_b$ and forward orbits proves injectivity of
$\mathcal R_b$.

(iii) Fix $x\in\mathbb R$.  $S_b(x)\in(u_1,u_2)\subset\mathcal B_b(Y)$ and
$\mathcal S_b(S_b(x))=x$, since both sides are the real-normalised repelling
time.  Move $y'$ from $0$ towards $\pm(h/2-Y-2)$.  As long as $S_b(x+\ii y')$
stays in the lens, continuation gives $\mathcal S_b(S_b(x+\ii y'))=x+\ii y'$.
It cannot leave through the corners $u_1,u_2$, where
$|\operatorname{Re}\mathcal S_b|\to\infty$.  It cannot leave through the arcs
either: there $\operatorname{Im}Z\in\{Y,h-Y\}$, so by (ii)
$|\operatorname{Im}\mathcal S_b|\ge h/2-Y-\pi/3-o(1)>|y'|$.  Hence
$S_b(\Sigma)\subset\mathcal B_b(Y)$, $T=\mathcal R_b\circ S_b$ on $\Sigma$, and
$T(w')-w'=\mathcal R_b(u)-\mathcal S_b(u)$ with $u=S_b(w')$ is the displayed
bi-infinite sum.
\end{proof}

\begin{lemma}[Uniform mode bound]\label{lem:modes}
Put $Y_0:=Y_B+3$.  There is $M$, which can be taken $\le1$, and $b_1<\eta$ such
that for $b\in(b_1,\eta)$ and all $n\ne0$,
\[
 |t_n|\le M\eee^{-2\pi|n|d},\qquad d:=\tfrac h2-Y_0-1 .
\]
\end{lemma}

\begin{proof}
Take $Y=Y_B+1$ in Lemma~\ref{lem:band}(iii), so that the estimate (ii) is
available on the closure of the lens.  $\Sigma$ has half-width
$h/2-Y_B-3=d+1$, so the lines $\operatorname{Im}w'=\pm d$ lie in $\Sigma$.  On them, by (iii) and (i),
$T-\mathrm{id}-t_0=\sum_kF_b(f_b^kS_b)-\text{mean}$ has modulus at most $2(C/Y_B+C\varepsilon)$.  Let $M$ be the supremum of this over
$b\in(b_1,\eta)$.  Cauchy's estimate on these lines gives the bound
for $n>0$ and $n<0$.  This does not use Theorem~\ref{thm:C}.
\end{proof}

For $H>0$ let $\mathfrak A_H$ be the space of $1$-periodic
$g=\sum g_n\eee^{2\pi\ii nz}$ with
$\|g\|_H=\sum|g_n|\eee^{2\pi|n|H}<\infty$.  It is a Banach algebra
(the weight is submultiplicative), its elements are holomorphic and bounded by
$\|g\|_H$ on $|\operatorname{Im}z|\le H$, and the projections
$\Pi_{<0},\Pi_{\ge0}$ onto negative and non-negative modes have norm $1$.  Also
$\|\eee^{g}\|_H\le\eee^{\|g\|_H}$ and
$\|\eee^{g_1}-\eee^{g_2}\|_H\le\|g_1-g_2\|_H\,\eee^{\max\|g_j\|_H}$.

\begin{theorem}[W2]\label{thm:W2}
Let $M,d$ be as in Lemma~\ref{lem:modes}, $\rho=\tfrac14$, and choose $D\ge1$
with $4\pi M\eee^{-2\pi D}\le\rho(1-\eee^{-2\pi D})^2$.  Put $H=d-D-\rho$ and
$\tau=T-\mathrm{id}-t_0$.  Assume $H>1$, which holds for $b$ near $\eta$ since
$d\sim\pi/\varepsilon$, and that $b$ is so close to $\eta$ that
$C_1\Lambda\eee^{4\pi(Y_0+D+1)}<\tfrac12$, where $C_1=4M+2\rho$.  Then:
\begin{enumerate}
\item The map $\mathcal T(Q)=-\Pi_{<0}[\tau\circ(\mathrm{id}+Q)]$ is a
  contraction, with constant $\le\tfrac12$, of the ball
  $\{Q\in\Pi_{<0}\mathfrak A_H:\|Q\|_H\le\rho\}$ into itself.  Its fixed point
  $Q^*$ and $P':=\Pi_{\ge0}[\tau\circ(\mathrm{id}+Q^*)]$ satisfy
  \[
   T\circ(\mathrm{id}+Q^*)=\mathrm{id}+t_0+P'\qquad\text{on }|\operatorname{Im}z|<H.
  \]
\item $\varphi_-=\mathrm{id}+Q^*$ is univalent on $\{\operatorname{Im}z<H-1\}$
  and $\varphi_+=\mathrm{id}+t_0+P'$ on $\{\operatorname{Im}z>-(H-1)\}$.  The
  function equal to $\Rr\circ(\varphi_+-\ii h)$ above and $S\circ\varphi_-$ below
  is a solution of~\eqref{eq:abel}; it is the sewing solution $\Kk^W_b$ of
  Proposition~\ref{prop:weld}.
\item The equation $\varphi_+(z_0)=\ii h$ has a unique solution with
  $|z_0-(\ii h-t_0)|<1$, and $|z_0-(\ii h-t_0)|\le C_1\Lambda\eee^{4\pi(Y_0+D+1)}$.
  After the translation $z\mapsto z+z_0$ one has $\Kk^W_b(0)=1$ and
  $P=\Rr^{-1}\circ\Kk^W_b-\mathrm{id}$ with $P(0)=0$.
\item For every $n\ge1$ the Fourier coefficients of this $P$ satisfy
  \[
   \Bigl|\frac{c_n}{\Lambda^n}-\tau_n(b)\Bigr|\le C_n\Lambda,
   \qquad C_n=C\,M(M+1)\,\eee^{2\pi(n+2)(Y_0+D+2)},
  \]
  with an absolute constant $C$.
\end{enumerate}
\end{theorem}

\begin{proof}
(1) Write $\tau\circ(\mathrm{id}+Q)=\sum_{m\ne0}t_m\eee^{2\pi\ii mz}\eee^{2\pi\ii mQ}$.
Then
\[
 \|\tau\circ(\mathrm{id}+Q)\|_H\le\sum_{m\ne0}|t_m|\eee^{2\pi|m|(H+\rho)}
 \le2M\sum_{m\ge1}\eee^{-2\pi mD}\le\rho
\]
by Lemma~\ref{lem:modes} and the choice of $D$.  Similarly
\[
 \|\tau\circ(\mathrm{id}+Q_1)-\tau\circ(\mathrm{id}+Q_2)\|_H\le
 \sum_{m\ne0}|t_m|\eee^{2\pi|m|(H+\rho)}2\pi|m|\,\|Q_1-Q_2\|_H
 \le\frac{4\pi M\eee^{-2\pi D}}{(1-\eee^{-2\pi D})^2}\|Q_1-Q_2\|_H ,
\]
which is at most $\tfrac12\|Q_1-Q_2\|_H$.  Since $|\operatorname{Im}(z+Q)|\le H+\rho<d$
on the strip, the Fourier series of $T$ converges at $z+Q(z)$ and equals $T$
there (Lemma~\ref{lem:band}).  The identity follows by splitting
$\tau\circ(\mathrm{id}+Q^*)=P'+\Pi_{<0}[\ldots]=P'-Q^*$.

(2) The series of $Q^*$ has only negative modes and converges on
$\operatorname{Im}z\le H$, so it converges on the whole lower half-plane
$\operatorname{Im}z<H$.  There $|{Q^*}'|\le2\pi\eee^{-2\pi}\rho\,/(1-\eee^{-2\pi})^2<1$ on
$\operatorname{Im}z\le H-1$, so $\operatorname{Re}\varphi_-'>0$ and $\varphi_-$
is univalent (Noshiro--Warschawski).  The same holds for $\varphi_+$ above.  On
the strip, $\Rr\circ(\varphi_+-\ii h)=\Rr\circ(T\circ\varphi_--\ii h)
=\Rr\circ T\circ\varphi_-=S\circ\varphi_-$, because $\ii h$ is a period of
$\Rr$ and $\Rr\circ T=S$.  The functional equation holds because $\varphi_\pm$
commute with $z\mapsto z+1$.  The resulting function realises the sewing of
Proposition~\ref{prop:weld}: the images of $\varphi_\pm$ are the two
half-cylinders and the seam in the $z$-plane is $\Gamma=\varphi_-^{-1}(\mathbb R)$, a curve in
$|\operatorname{Im}z|\le\rho$.  Since $\varphi_\pm-\mathrm{id}$ is bounded, the
argument principle shows that the restrictions of $\varphi_\pm$ to the two sides of
$\Gamma$ map onto the two half-cylinders.  By uniqueness of the uniformisation
it is $\Kk^W_b$ up to translation.  (The identity $\Rr\circ T=S$ is used on the
band, where orbits stay in $D_{r'}$ and $\mathcal R_b$ inverts $\Rr$;
Lemma~\ref{lem:band}.)

(3) For $\operatorname{Im}z\ge h/2-1$ we have
$|P'(z)-p'_0|\le\sum_{n\ge1}|p'_n|\eee^{-2\pi n(h/2-1)}$ and
$|p'_n|\le\rho\,\eee^{-2\pi nH}$.  Moreover
$p'_0=\sum_{m\ge1}t_m[\eee^{2\pi\ii mQ^*}]_{-m}$, and
$|[\eee^{2\pi\ii mQ^*}]_{-m}|\le\eee^{2\pi m\rho}\eee^{-2\pi mH}$, so
$|p'_0|\le2M\eee^{-2\pi(d+H-\rho)}\le2M\Lambda\eee^{2\pi(2Y_0+D+2+2\rho)}$, and
$\sum_{n\ge1}|p'_n|\eee^{-2\pi n(h/2-1)}\le2\rho\Lambda\eee^{2\pi(Y_0+D+2+\rho)}$.
Both are $\le C_1\Lambda\eee^{4\pi(Y_0+D+1)}/2$.  Since
$\operatorname{Im}(\ii h-t_0)=h/2$ and $|{P'}'|\le\rho$ there, Rouch\'e's theorem
on the disc of radius $1$ gives the unique root and the bound, using the
smallness hypothesis.  Then
$\Kk^W_b(z_0)=\Rr(\varphi_+(z_0)-\ii h)=\Rr(0)=1$, and after translation
$P(z)=\varphi_+(z+z_0)-\ii h-z$, so $P(0)=0$.

(4) For $n\ge1$ the $n$-th mode of $P$ is $c_n=p'_n\eee^{2\pi\ii nz_0}$.
Because $Q^*$ has no constant or positive modes,
$[\eee^{2\pi\ii mQ^*}]_0=1$ and $[\eee^{2\pi\ii mQ^*}]_k=0$ for $k>0$.  Hence
$p'_n=t_n+\sum_{m>n}t_m[\eee^{2\pi\ii mQ^*}]_{n-m}$, and
\[
 |p'_n-t_n|\le\sum_{m>n}M\eee^{-2\pi md}\eee^{2\pi m\rho}\eee^{-2\pi(m-n)H}
 \le2M\eee^{-2\pi n(d-\rho)}\eee^{-2\pi(d+H-\rho)} .
\]
With $z_0=\ii h-t_0+\zeta$, $|\zeta|\le C_1\Lambda\eee^{4\pi(Y_0+D+1)}$ by (3), we have
$\eee^{2\pi\ii nz_0}=\Lambda^n\eee^{-2\pi\ii nt_0}\eee^{2\pi\ii n\zeta}$ and
$|t_n\eee^{-2\pi\ii nt_0}|=|\tau_n|$, so
$|t_n|=|\tau_n|\eee^{-\pi nh}$.  Dividing,
\[
 \frac{c_n}{\Lambda^n}-\tau_n=\tau_n\bigl(\eee^{2\pi\ii n\zeta}-1\bigr)
 +(p'_n-t_n)\eee^{-2\pi\ii nt_0}\eee^{2\pi\ii n\zeta}.
\]
In the first term $|\tau_n|=|t_n|\eee^{\pi nh}\le M\eee^{2\pi n(Y_0+1)}$ and
$|\eee^{2\pi\ii n\zeta}-1|\le4\pi n|\zeta|$ when $2\pi n|\zeta|\le1$; in
the second, $|\eee^{2\pi\ii n\zeta}|\le\eee$ in that case.  If instead
$2\pi n|\zeta|>1$, then
$\Lambda\ge(2\pi nC_1)^{-1}\eee^{-4\pi(Y_0+D+1)}$, and the trivial bound
$|c_n|\Lambda^{-n}+|\tau_n|\le C'M\eee^{2\pi n(Y_0+D+2)}$ is already
$\le C_n\Lambda$ (this uses $D\ge1$).  In the second term,
$|(p'_n-t_n)\eee^{-2\pi\ii nt_0}|\le2M\eee^{2\pi n(Y_0+1+\rho)}\Lambda
\eee^{2\pi(2Y_0+2+D+2\rho)}$, by $d=h/2-Y_0-1$ and $H=d-D-\rho$.  Adding the two
gives the stated bound, since
$n(Y_0+1+\rho)+2Y_0+2D+2.5\le(n+2)(Y_0+D+2)$.
\end{proof}

\begin{corollary}\label{cor:Kw}
For the sewing solution, $\hat c_n(\Kk^W_b)\to B_n\eee^{2\pi\ii na}$ for every
$n\ge1$ as $b\to\eta^-$.  Moreover,
$\kappa_n(\Kk^W_b)\to B_n/B_1^n$ and
$\lim\arg\hat c_1=\arg B_1+2\pi a$.
\end{corollary}

\begin{proof}
Theorem~\ref{thm:W2}(4) and Corollary~\ref{cor:tauconv}: $C_n$ does not depend
on $b$ and $\Lambda\to0$.  Proposition~\ref{prop:horn-nonzero} makes $c_1$
nonzero for $b$ near $\eta$, justifying division and the limiting phase.
\end{proof}

\begin{corollary}[Local parabolic limit of the sewn solution]
\label{cor:sewn-cusp-local}
Put $u^*=1/\eee-1$ and $a=\Phi_{\mathrm{att}}(u^*)$.  On some fixed disk
about $z=0$, the function
\[
 F_\eta^{\mathrm{att}}(z)
 =\eee\bigl(1+\Phi_{\mathrm{att}}^{-1}(z+a)\bigr)
\]
is holomorphic, satisfies $F_\eta^{\mathrm{att}}(0)=1$, and has
$(F_\eta^{\mathrm{att}})'(0)\ne0$.  As $b\to\eta^-$ through real bases,
$\Kk_b^W\to F_\eta^{\mathrm{att}}$ locally uniformly on that disk,
with all $z$-derivatives.  Consequently there are $r>0$ and $b_2<\eta$
such that, for $b_2<b<\eta$, the only preimage of $1$ under $\Kk_b^W$
in $B(0,r)$ is $0$; both $\Kk_b^W$ and $1/\Kk_b^W$ are uniformly
bounded there.
\end{corollary}

\begin{proof}
Theorem~\ref{thm:C}(1) gives locally uniform convergence of
$A_b(u):=\Rr_b^{-1}(\eee(1+u))$ to
$A_\eta(u):=\Phi_{\mathrm{att}}(u)-a$ near $u^*$, with $A_b(u^*)=0$.
The derivative $A_\eta'(u^*)$ is nonzero: the series in the proof of
Theorem~\ref{thm:C}, Lemmas~\ref{lem:F} and~\ref{lem:petal}, and Cauchy's
estimate give $\Phi_{\mathrm{att}}'(u)=2/u^2+O(1/u)$ on a narrower
attracting petal.  The Abel identity and $f_\eta'(u)\ne0$ propagate
nonvanishing back along the orbit of $u^*$.  Uniform
convergence of these local inverses therefore gives
$\Rr_b\to F_\eta^{\mathrm{att}}$ on a fixed disk about $0$.

In Theorem~\ref{thm:W2}, write
$P_b(z)=\zeta_b+P_b'(z+z_0(b))$, where
$\zeta_b=z_0(b)+t_0(b)-\ii h(b)\to0$.  The proof of that theorem gives
$p'_0(b)\to0$ and $|p'_n(b)|\le\rho\eee^{-2\pi nH(b)}$ for $n\ge1$,
while $H(b)\to\infty$ and $\operatorname{Im}z_0(b)=h(b)/2+o(1)$.
Hence $P_b\to0$ uniformly on every fixed bounded $z$-disk on which
the upper representation is defined.  That representation contains a
fixed disk about $0$ for all sufficiently large $b$, and
$\Kk_b^W(z)=\Rr_b(z+P_b(z))$ there.  This proves the asserted convergence;
Cauchy's formula gives convergence of derivatives.  Choose a small
closed disk on which $F_\eta^{\mathrm{att}}$ is nonzero and
$F_\eta^{\mathrm{att}}-1$ has only its simple zero at $0$.
Rouch\'e's theorem and $\Kk_b^W(0)=1$ give the final assertions.
\end{proof}

\noindent
The construction is directly computable (\texttt{weld\_solve.py}).  At
$b=1.10$ the iteration contracts by $\sim10^{-4}$ per step and gives
$\hat c_1=0.0864272029\,\eee^{1.5510921403\ii}$, against
$\tau_1=0.0864272022\,\eee^{1.5510921408\ii}$.  The relative difference is
$8.8\cdot10^{-9}=0.38\Lambda$, and at $b=1.30$ it is $3.9\cdot10^{-19}=0.40\Lambda$.
At these bases $h$ is too small for the hypotheses of Theorem~\ref{thm:W2}, so
the agreement is not an instance of the theorem.  It does match the first-order
expansion of the correction,
$c_1/(\Lambda\tau_1)-1\approx\Lambda\bigl[-4\pi^2|\tau_1|^2-2\pi\ii\tau_1
-4\pi\ii\tau_2\bar\tau_1/\tau_1\bigr]$, whose modulus is $0.3834\Lambda$ and
$0.3984\Lambda$ respectively (derived independently by a referee).

\subsection{The Kneser side II: identification (W1)}\label{sec:W1}

It remains to identify $\Kk^W_b$ with the boundary value $\Kk_b$
of~\eqref{eq:K}.  We prove the local part and isolate the one global input.

\begin{theorem}[W1, local]\label{thm:W1}
Let $b_0\in(b_1,\eta)$.  There is a disc $\Delta\ni b_0$ in the base plane such
that:
\begin{enumerate}
\item For $b'\in\Delta$, with $L(b'),L_2(b')$ the continued fixed points,
  $\Rr_{b'},S_{b'}$ the regular superfunctions, $\varpi(b')=2\pi\ii/\log\lambda(b')$
  the period of $\Rr_{b'}$, and $\hat T_{b'}=\Rr_{b'}^{-1}\circ S_{b'}+\varpi/2$,
  the contraction of Theorem~\ref{thm:W2} works with $\rho$ and $D$
  unchanged, $2M$ in place of $M$, and $H-\tfrac12$ in place of $H$.  The solution is
  $\Kk^W_{b'}=\Rr_{b'}\circ(\varphi_++\varpi)$ above, where
  $\varphi_+=T_{b'}\circ\varphi_-$, and $S_{b'}\circ\varphi_-$ below.  It is
  normalised by the root of $\varphi_+(z_0)=-\varpi$ in the disc of radius $1$
  about $-\varpi-t_0$.  (For real $b'$, $\varpi=-\ii h$.)
  Its fixed point, the resulting solution $\Kk^W_{b'}$, and the normalising root
  $z_0(b')$ are holomorphic in $b'$.
\item For $b'\in\Delta$ with $\operatorname{Im}b'>0$ one has $\arg\lambda>0$,
  $\arg\lambda_2<0$, and $\Kk^W_{b'}(z)\to L(b')$ as $\operatorname{Im}z\to+\infty$,
  $\Kk^W_{b'}(z)\to L_2(b')$ as $\operatorname{Im}z\to-\infty$, uniformly for
  $\operatorname{Re}z$ in compacts.  In the upper and lower half-planes (of
  the translated coordinate) it has the representations
  $\Rr_{b'}\circ(\mathrm{id}+\theta_{\mathrm{up}})$ and
  $S_{b'}\circ(\mathrm{id}+\theta_{\mathrm{dn}})$.  These are the defining
  representations of the complex-base construction in \texttt{\_cbuild.py}.
\item (Local uniqueness.)  If a solution $F$ with $F(0)=1$ admits such
  representations and, after the translation that puts its seam
  representation in the form $T\circ(\mathrm{id}+Q)=\mathrm{id}+t_0+P'$, has
  $\|Q\|_H\le\rho$ and its normalising root in that disc, then
  $F=\Kk^W_{b'}$.  (Without the root condition the statement is false: the
  translates of $\Kk^W_{b'}$ by the other roots $c_j$, $j\ge1$, have the same $Q$.)
\end{enumerate}
\end{theorem}

\begin{proof}
(1) By Lemma~\ref{lem:band}(iii), $T$ is holomorphic on the strip of half-width $d+1$,
so there is slack $1$ in $d$; we use half of it.  The Koenigs maps at $L(b')$, $L_2(b')$ and the principal logarithms
defining $\Rr^{-1}$ depend holomorphically on $b'$.  To choose one
parameter disc for the unbounded strip, first work on the compact fundamental
rectangle $0\le\operatorname{Re}z\le1$,
$|\operatorname{Im}z|\le d(b_0)+\tfrac12$.  Its image under $S_{b_0}$
is compactly contained in the common domain of the two chosen Abel charts;
their local parameter continuations therefore give a single disc $\Delta$
on this rectangle.  The identity $T_{b'}(z+1)=T_{b'}(z)+1$ extends the
result to the whole strip.  Thus $\hat T_{b'}-\mathrm{id}$ is jointly
holomorphic on $\Delta\times\{|\operatorname{Im}z|<d(b_0)+\tfrac12\}$.
Shrinking $\Delta$ again,
its modes satisfy $|t_n(b')|\le2M\eee^{-2\pi|n|(d(b_0)-1/2)}$ (Cauchy), and the
estimates of Theorem~\ref{thm:W2}(1) hold with $2M$ in place of $M$:
the original choice of $D$ gives a Lipschitz constant at most
$2\rho=\tfrac12$ and a ball-image norm at most
$\rho(1-\eee^{-2\pi D})/\pi<\rho$.  The iterates $\mathcal T^k(0)$ are holomorphic in $b'$
and converge uniformly, so $Q^*$ is holomorphic in $b'$.  So are $P'$ and,
by the implicit function theorem (Rouch\'e), $z_0$.

(2) From $\lambda\eee^{-\lambda}=\log b'$,
$\partial\lambda/\partial b=\eee^{\lambda}/(b(1-\lambda))$ is positive for
$\lambda<1$ and negative for $\lambda_2>1$.  Hence $\operatorname{Im}b'>0$
gives $\operatorname{Im}\lambda>0>\operatorname{Im}\lambda_2$ for $\Delta$
small.  Above, $\Kk^W=\Rr(z+c+o(1))$, and
$\Rr(w)=\sigma^{-1}(c'\lambda^w)\to L$ when
$\operatorname{Re}(w\log\lambda)\to-\infty$.  This is the case as
$\operatorname{Im}w\to+\infty$, since
$\operatorname{Re}(w\log\lambda)=\operatorname{Re}w\log|\lambda|-\operatorname{Im}w\arg\lambda$.
Below, the same argument with $\arg\lambda_2<0$ applies.  The representations
are those of Theorem~\ref{thm:W2}(2), transported by $z\mapsto z+z_0$.

(3) The fixed point of a contraction is unique in its ball.  $P'$ and $z_0$
are then determined, the latter by $F(0)=1$ together with the root
condition.
\end{proof}

\begin{corollary}\label{cor:W1}
Suppose that, for all $b'$ in a nonempty open subset of
$\Delta\cap\{\operatorname{Im}b'>0\}$, the Kneser-type solution $\kappa_{b'}$
continued from $b>\eta$ satisfies the hypothesis of Theorem~\ref{thm:W1}(3), and
that $\kappa$ is jointly holomorphic in $(b',z)$ off a closed singular set whose
complement is connected in $z$.  Then $\kappa_{b'}=\Kk^W_{b'}$ on
$\Delta\cap\{\operatorname{Im}b'>0\}$.  The boundary value $\Kk_b$ exists on
$\Delta\cap\mathbb R$ and equals $\Kk^W_b$.  The upper family continues
analytically across $\Delta\cap\mathbb R$.  The conclusions of
Corollary~\ref{cor:Kw} then hold for $\Kk_b$, including the first-mode
and ratio limits of Conjectures~\ref{conj:horn} and~\ref{conj:horn2}.
\end{corollary}

\begin{proof}
On the open set, Theorem~\ref{thm:W1}(3) gives $\kappa=\Kk^W$.  Both families
are holomorphic in $(b',z)$ for $z$ in a lower half-plane, where $\Kk^W=S\circ\varphi_-$
has no cuts.  So they agree on $\Delta\cap\{\operatorname{Im}b'>0\}$ by the
identity theorem, and then everywhere by continuation in $z$.  The rest is Theorem~\ref{thm:W1}(1) and
Corollary~\ref{cor:Kw}.
\end{proof}

The hypothesis of Corollary~\ref{cor:W1} is an inequality on an open set of bases.  It
concerns only the lower correction of $\kappa$ near the seam.  The ladders
are consistent with it: they reproduce $\tau_1$ and $t_2/t_1^2$ to their
precision (Table~\ref{tab:tau}), and the explicit sewing solution matches the
ladder value of $\hat c_1$ at $b=1.10$ to eight digits.  This check is
numerical and not certified.  A proof free of numerics would need a global uniqueness theorem
for Kneser-type solutions at complex bases inside the Shell--Thron region,
one that selects the no-winding solution: the class defined only by the limits
at $\pm\ii\infty$ also contains solutions of the form
$\Rr\circ(z+2\pi\ii mz/\log\lambda+\ldots)$.  The base continuation from
$b>\eta$ passes the Shell--Thron boundary, where $|\Lambda|=1$ and the present
contraction does not apply.  This is the one input we do not supply.

\subsection{A characterisation of the sewing solution, and boundary values}\label{sec:char}

Corollary~\ref{cor:W1} needs a quantitative input about $\kappa$ at one base.
Here we replace it with a qualitative characterisation.  From it follows a
boundary-value theorem whose only hypothesis is a uniform bound on the two
$\theta$-representations along a ladder $b+\ii y$, $y\downarrow0$.

\begin{definition}\label{def:class}
For real $b\in(1,\eta)$ and $Y>0$ let $\mathcal C_b(Y)$ be the set of
solutions $F$ of~\eqref{eq:abel} with $F(0)=1$ such that
(``solution'' means holomorphic off a closed singular set and satisfying the
equation where defined)
\begin{enumerate}
\item[(a)] $F=\Rr\circ(\mathrm{id}+P)$ on $\{\operatorname{Im}z>-2Y\}$, with $P$
  holomorphic, $1$-periodic and bounded there, and $P(0)=0$ (the branch
  $\Rr^{-1}(1)=0$);
\item[(b)] $F=S\circ(\mathrm{id}+Q)$ on $\{\operatorname{Im}z<-Y\}$, with $Q$
  holomorphic, $1$-periodic and bounded there, and
  $|\operatorname{Im}Q-h/2|\le1$ on the strip $O=\{-2Y<\operatorname{Im}z<-Y\}$.
\end{enumerate}
\end{definition}

Condition (b) says that the real line of repelling time, i.e.\ the segment
$(L,L_2)$, sits half a period of attracting time below the real axis of $F$.
Heuristically, some such condition is necessary.  For the sewing solution the
roots of $\varphi^W_+(c)\in\ii h\mathbb Z$ near the axis are
$c_j=c_0+\ii jh+O(\Lambda)$.  The translates $\Kk^W(\cdot+c_j)$ keep an
$\Rr$-representation above and an $S$-representation on some lower half-plane,
and give $|c_1|\asymp\Lambda^{1+j}$.  Condition (b) excludes $j\ne0$, as the
proof shows.

\begin{theorem}[Characterisation]\label{thm:char}
There are $b_1<\eta$ and $Y_B$ such that for $b\in(b_1,\eta)$ and
$Y_0+D+3<Y<(h-Y_0-D-3)/2$ (with $Y_0,D$ as in Theorem~\ref{thm:W2}),
\[
 \mathcal C_b(Y)=\{\Kk^W_b\},
\]
with $\Kk^W_b$ normalised as in Theorem~\ref{thm:W2}(3).  In particular
$\hat c_n\to B_n\eee^{2\pi\ii na}$ for every $n$ holds for the unique element
of $\mathcal C_b(Y)$.
\end{theorem}

\begin{proof}
\emph{$\Kk^W_b\in\mathcal C_b(Y)$.}  In the normalised coordinate,
$\varphi_-(z)=z+z_0+Q^*(z+z_0)$ is defined for
$\operatorname{Im}(z+z_0)<H$, i.e.\ for $\operatorname{Im}z<-(Y_0+1+D+\rho)+O(\Lambda)$,
which contains $\{\operatorname{Im}z<-Y\}$.  There $|Q-z_0|\le\rho$, so
$|\operatorname{Im}Q-h/2|\le\rho+O(\Lambda)<1$.  The $\Rr$-representation holds
for $\operatorname{Im}(z+z_0)>-(H-1)$, i.e.\ above
$-(h-Y_0-D-2-\rho)$, which contains $\{\operatorname{Im}z>-2Y\}$.  $P$ and $Q$
are bounded there.

\emph{Uniqueness.}  Let $F\in\mathcal C_b(Y)$, $\varphi_+=\mathrm{id}+P$ and
$\varphi_-=\mathrm{id}+Q$.  Work with the strip
$\Sigma_q=\{|\operatorname{Im}w'|<d-D\}$ of $S$-time.  By Lemma~\ref{lem:modes} and
the choice of $D$, on $\overline{\Sigma_q}$
\[
 |T-\mathrm{id}-t_0|\le\frac{2M\eee^{-2\pi D}}{1-\eee^{-2\pi D}}\le\frac{\rho}{2\pi},\qquad
 |T'-1|\le\rho<1 ,
\]
so $T$ is univalent on the convex set $\Sigma_q$ (Noshiro--Warschawski),
continuous up to the boundary, and within $\rho/2\pi$ of a translation.

(i) \emph{The overlap.}  Let $O=\{-2Y<\operatorname{Im}z<-Y\}$.  For $z\in O$,
$\operatorname{Im}\varphi_-(z)\in(h/2-2Y-1,\,h/2-Y+1)\subset(-(d-D),d-D)$,
because $2Y<h-Y_0-D-2$ and $Y>Y_0+D+2$.  Since
$\sigma\circ\Rr(w)=\sigma(1)\lambda^w$ wherever $\Rr$ is defined,
$\Rr\circ\varphi_+=S\circ\varphi_-=\Rr\circ T\circ\varphi_-$ gives
$\lambda^{\varphi_+}=\lambda^{T\circ\varphi_-}$.  Hence
$\varphi_+=T\circ\varphi_--\ii h+k\ii h$ on $O$ for a constant $k\in\mathbb Z$.
No injectivity is needed for this step.

(ii) \emph{The lifted sewn surface.}  Let $V_-=\{\operatorname{Im}w'<d-D\}$
(repelling time) and let $V_+$ be the region above the curve
$T(\{\operatorname{Im}w'=-(d-D)\})-\ii h$ (attracting time).  Let $\tilde X$ be
$V_-\sqcup V_+$ modulo $w=T(w')-\ii h$ for $w'\in\Sigma_q$.  Since $T$ is a
homeomorphism of $\overline{\Sigma_q}$ onto its image, the identification
graph is closed and $\tilde X$ is a Hausdorff Riemann surface.  Translation by
$1$ acts freely on it, and $X=\tilde X/\mathbb Z\cong\mathbb C^*$, both ends
being punctured discs.  The function $\tilde K$, equal to $\Rr$ on $V_+$ and
$S$ on $V_-$, is well defined, because $\Rr(T(w')-\ii h)=S(w')$ on the band.
The sewing solution is $\tilde K\circ\tilde\Phi^W$, where $\tilde\Phi^W$ is the
lift of its uniformisation; sewing along the real line or along any other
essential curve in $\Sigma_q$ gives the same $\tilde X$.

Define $\tilde\Phi:\mathbb C\to\tilde X$ by $\varphi_+-k\ii h$ on
$\{\operatorname{Im}z>-2Y\}$ and $\varphi_-$ on $\{\operatorname{Im}z<-Y\}$.
These agree on $O$ by (i).  They land in $V_\mp$ as follows.
$\operatorname{Im}Q$ is bounded and harmonic on a half-cylinder, hence extends
harmonically over the puncture.  So on $\{\operatorname{Im}z\le-Y-\delta\}$ it is
bounded by its values on the line $\operatorname{Im}z=-Y-\delta\subset O$.
This gives $\operatorname{Im}\varphi_-\le-Y+h/2+1<d-D$.  Likewise, on a line
in $O$ near $\operatorname{Im}z=-2Y$, (i) gives
$\operatorname{Im}(\varphi_+-k\ii h)\ge\operatorname{Im}z-1-\rho/2\pi$.  The minimum
principle then gives $\operatorname{Im}(\varphi_+-k\ii h)\ge-2Y-1-\rho/2\pi$
above that line.  This exceeds the height
$-h/2-d+D+\rho/2\pi$ of the lower boundary of $V_+$, because
$2Y<h-Y_0-D-3$.

(iii) \emph{$\tilde\Phi$ is a translate of $\tilde\Phi^W$.}  $\tilde\Phi$
commutes with $z\mapsto z+1$ and descends to a holomorphic
$\Phi:\mathbb C/\mathbb Z\to X\cong\mathbb C^*$.  $\Phi$ is proper, since
$\varphi_\pm-\mathrm{id}$ are bounded, and has local degree one at each end,
since $P\to p_0$ and $Q\to q_0$.  A proper holomorphic map of $\mathbb C^*$ that
fixes the ends with local degree one is linear.  Hence
$\tilde\Phi=\tilde\Phi^W(\cdot+c)$ for some $c\in\mathbb C$, and
$F=\tilde K\circ\tilde\Phi=\Kk^W_b(\cdot+c)$ in the coordinate of
Theorem~\ref{thm:W2}(1)--(2).  Univalence of $\varphi_\pm$ is a conclusion, not
an assumption.

(iv) \emph{The root.}  Take a line $\ell\subset O$ with
$\operatorname{Im}z=-Y''$ and $Y''$ close to $2Y$.  There
$\operatorname{Im}\varphi_-\le-2Y+h/2+1+o(1)<H-\rho$, so $\varphi_-(\ell)$ lies
in $\varphi_-^W(\{\operatorname{Im}<H\})$.  The inclusion
$\varphi_-^W(\{\operatorname{Im}<H\})\supset\{\operatorname{Im}<H-\rho\}$ holds because
$|Q^*|\le\rho$: for $w$ in the smaller set, $\varphi^W_--w$ has winding
number $1$ around $0$ on the boundary of a large fundamental rectangle.  By injectivity of
$\tilde\Phi^W$, the points $z+c$ with $z\in\ell$ lie in
$\{\operatorname{Im}<H\}$.  There $Q(z)=c+Q^*(z+c)$, so
$|\operatorname{Im}c-h/2|\le1+\rho$.  On the upper chart,
$\varphi_+(z)-k\ii h=\varphi^W_+(z+c)-\ii h$, and $P(0)=0$ gives
$\varphi^W_+(c)=\ii h(1-k)$.  Since $|P'|\le\rho$ where $\varphi^W_+$ is defined,
$|c-(\ii h(1-k)-t_0)|\le\rho$.  Comparing imaginary parts with the bound on
$\operatorname{Im}c$ gives $|k|h\le1+2\rho$, so $k=0$.  Then $c$ lies in the
Rouch\'e disc of Theorem~\ref{thm:W2}(3), so $c=z_0$ and $F=\Kk^W_b$.  The last
sentence of the theorem is Corollary~\ref{cor:Kw}.
\end{proof}

\begin{theorem}[Boundary values]\label{thm:bv}
Let $b\in(b_1,\eta)$ and $Y$ be as in Theorem~\ref{thm:char}.  Suppose that for
every sequence $y_k\downarrow0$ the Kneser-type solutions $\kappa_k=\kappa_{b+\ii y_k}$
admit representations (a) and (b) with $\Rr_{b+\ii y_k}$, $S_{b+\ii y_k}$ in
place of $\Rr_b,S_b$.  Suppose also that $P_k,Q_k$ are bounded uniformly in
$k$, and that $|\operatorname{Im}Q_k-h(b_k)/2|\le\tfrac12$ on $O$, with $b_k=b+\ii y_k$.  Then
$\kappa_k\to\Kk^W_b$ locally uniformly off the (moving) singular sets, the boundary
value~\eqref{eq:K} exists, and its separation coefficients
satisfy $|\hat c_n(b)-\tau_n(b)|\le C_n\Lambda$.  Consequently, if the
hypothesis holds for all $b$ near $\eta$,
$\hat c_n\to B_n\eee^{2\pi\ii na}$ for all $n$ and
Conjectures~\ref{conj:horn} and~\ref{conj:horn2} hold for $\Kk_b$.
\end{theorem}

\begin{proof}
$\Rr_{b'}\to\Rr_b$ and $S_{b'}\to S_b$ locally uniformly (Koenigs maps depend
holomorphically on the base), and so do the constants $h,\lambda$.  By
Montel's theorem a subsequence of $(P_k,Q_k)$ converges locally uniformly to
$(P,Q)$, bounded and periodic, with $P(0)=0$ and
$|\operatorname{Im}Q-h/2|\le\tfrac12$.  The limit
$F=\Rr_b\circ(\mathrm{id}+P)=S_b\circ(\mathrm{id}+Q)$ solves~\eqref{eq:abel}
with $F(0)=1$, so $F\in\mathcal C_b(Y)$ and $F=\Kk^W_b$ by
Theorem~\ref{thm:char}.  Every subsequence has the same limit.  The Fourier
coefficients of $P_k$ converge to those of $P$, and Theorem~\ref{thm:W2}(4)
applies.
\end{proof}

\noindent
The hypothesis of Theorem~\ref{thm:bv} is the form in which the complex-base
construction computes $\kappa$: \texttt{\_cbuild.py} solves precisely for
bounded $\theta_{\mathrm{up}},\theta_{\mathrm{dn}}$ with these two
representations.  Numerically it holds along every ladder we ran.  The ladder
values of $c_1$ are of size $\Lambda$, not $O(1)$ or $\Lambda^2$.  This is
consistent with the height condition, but it pins the offset only in
combination with the value of $|\tau_1|$.  What is not proved is that the hypothesis holds
for the family continued from $b>\eta$ through the Shell--Thron boundary.
That is a statement about Paulsen's family, not about the separation.

\subsection{Uniform bounds come for free}\label{sec:unif}

The hypothesis of Theorem~\ref{thm:bv} asks for bounds that are uniform along
a ladder.  We show that uniformity is automatic.  It suffices that the
continued family lie in the characterisation class on some open set of bases,
with no quantitative condition.

Let $N=\bigcup_{b_0\in(b_1,\eta)}\Delta(b_0)$ be the union of the discs of
Theorem~\ref{thm:W1}, and $N^+=N\cap\{\operatorname{Im}b'>0\}$.  Both are
connected: the discs are centred on an interval, and consecutive upper
half-discs overlap.  For complex $b'$, Definition~\ref{def:class} is read with
$\Rr_{b'},S_{b'}$, the period $\varpi(b')=2\pi\ii/\log\lambda(b')$ in place of
$-\ii h$, and $h(b')=|\operatorname{Im}\varpi(b')|$.

\begin{theorem}[Uniformity]\label{thm:unif}
Shrinking the discs if necessary:
\begin{enumerate}
\item The sewing solutions $\Kk^W_{b'}$ of the discs agree on overlaps and form
  one family, holomorphic in $(b',z)$ on $N\times\mathbb C$ minus the singular
  set.  Its representation functions $P_{b'},Q_{b'}$ are bounded uniformly on
  $\{\operatorname{Im}z>-2Y\}$ and $\{\operatorname{Im}z<-Y\}$ respectively,
  for $b'$ in any compact subset of $N$, including compact pieces of the real
  segment.
\item For every $b'\in N$, $\mathcal C_{b'}(Y)=\{\Kk^W_{b'}\}$.
\item Let $\kappa$ be any family of solutions, jointly holomorphic in $(b',z)$
  off a closed singular set, whose regular set in $z$ is connected for each
  $b'$ and which contains $N^+\times W$ for some open
  $W\subset\{\operatorname{Im}z<-Y\}$, with $\kappa_{b'}\in\mathcal C_{b'}(Y)$ for all $b'$ in a nonempty open
  subset of $N^+$.  Then $\kappa=\Kk^W$ on $N^+$.  Its
  $\theta$-representations are bounded uniformly along every ladder
  $b+\ii y$, $y\downarrow0$, $b\in(b_1,\eta)$.  So the hypothesis of
  Theorem~\ref{thm:bv} holds, and with it the conclusions there, including
  Conjectures~\ref{conj:horn} and~\ref{conj:horn2} for $\kappa$.
\end{enumerate}
\end{theorem}

\begin{proof}
(1) On $\Delta(b_0)\cap\Delta(b_0')$ let $H\le H'$ be the two strip widths.
Since $\|\cdot\|_H\le\|\cdot\|_{H'}$, the fixed point of the $H'$-contraction
lies in the $H$-ball.  It is a fixed point there, so it equals the
$H$-fixed point by uniqueness.  The normalising roots coincide by Rouch\'e.
The singular set $\{(b',z):\ \varphi_+(z)+\varpi(b')\in\text{cuts of }\Rr_{b'}\}$
is closed in $N\times\mathbb C$, because the cuts move continuously with $b'$.  Uniform bounds on compacta follow
from continuity of the fixed point in the base and from the explicit bounds
$\|Q^*\|_H\le\rho$, $\|P'\|_H\le\rho$ of Theorem~\ref{thm:W2}, which hold
throughout each disc by Theorem~\ref{thm:W1}(1).  The continuation of the
representations to the half-planes of Definition~\ref{def:class} is as in the
proof of Theorem~\ref{thm:char}: through the band, which persists for $b'$
near the real axis because $T_{b'}$ depends holomorphically on $b'$.

(2) The proof of Theorem~\ref{thm:char} runs verbatim with the following
per-base data.  The admissible range is
$Y_0+D+3<Y<(h(b')-Y_0-D-3)/2$ with $H-\tfrac12$ in place of $H$; it contains a
fixed $Y$ for all $b'$ in each small disc.  The strip $\Sigma_q$, on which
$|\hat T_{b'}-\mathrm{id}-\hat t_0|\le\rho/2\pi$ and $|\hat T_{b'}'-1|\le\rho$, persists
by continuity in $b'$.  So does the identity $\Rr_{b'}\circ T_{b'}=S_{b'}$ on
it, by analytic continuation in $b'$ of an identity valid at real $b'$.  The identity $\sigma\circ\Rr(w)=\sigma(1)\lambda^w$
holds for all $b'$.  The sewn surface is again $\mathbb C^*$.  Proper maps of
$\mathbb C^*$ fixing the ends with local degree one are linear.  Finally, the
normalising roots $c_j=c_0-j\varpi+O(|\Lambda|)$ have imaginary parts
separated by $h(b')$.  None of these steps uses reality of the base.  The one place it entered,
$\operatorname{Im}t_0=h/2$, is replaced by $\operatorname{Im}\hat t_0\approx0$ for the
shifted map $\hat T_{b'}$, which holds after shrinking the discs.

(3) By (2), $\kappa_{b'}=\Kk^W_{b'}$ on a nonempty open subset of $N^+$.  For
$z\in W$, where $\Kk^W_{b'}=S_{b'}\circ\varphi_-$ has no cuts, both sides are
holomorphic in $b'$ on the connected set $N^+$.  So they agree on
$N^+\times W$, and then for all $z$ off the singular set by continuation
in $z$ (the complement of the cuts is connected).  Uniform boundedness is then (1).
\end{proof}

\begin{corollary}[A bounded complex cusp corridor]\label{cor:cusp-corridor}
There are $b_3<\eta$, $r>0$, and a connected open set
$N_\eta\subset\mathbb C$ containing $(b_3,\eta)$, with
$\eta\in\partial N_\eta$, such that $\Kk^W$ is jointly holomorphic on
$N_\eta\times B(0,r)$ and
$\Kk_b^W\to F_\eta^{\mathrm{att}}$ locally uniformly there as
$b\to\eta$ within $N_\eta$, together with all $z$-derivatives.
Throughout $N_\eta$, the only preimage of $1$ in $B(0,r)$ is $0$,
and $\Kk_b^W$ and $1/\Kk_b^W$ are uniformly bounded on that disk.
No lower bound on the width of $N_\eta$ is asserted.
\end{corollary}

\begin{proof}
Choose $r$ inside the disk of Corollary~\ref{cor:sewn-cusp-local}
so that $F_\eta^{\mathrm{att}}$ has no zero on $\overline{B(0,r)}$
and $F_\eta^{\mathrm{att}}-1$ has just its simple zero at $0$.
The minimum of $|F_\eta^{\mathrm{att}}|$ on the closed disk and of
$|F_\eta^{\mathrm{att}}-1|$ on its boundary is positive; let one third
of the smaller minimum be $m$.  By that corollary, for every real
$x\in(b_3,\eta)$, with $b_3$ sufficiently close to $\eta$,
$\|\Kk_x^W-F_\eta^{\mathrm{att}}\|_{\overline{B(0,r)}}<m$.
Theorem~\ref{thm:W1} gives a base disk $U_x$ centred at $x$ on which
$\Kk_b^W$ is jointly holomorphic near the closed $z$-disk.  Shrink it
so that $|b-x|<(\eta-x)/4$ and
\[
 \|\Kk_b^W-\Kk_x^W\|_{\overline{B(0,r)}}
 <\min(m,\eta-x)\qquad(b\in U_x).
\]
Two overlapping disks centred on the real line have an overlapping real
interval.  Both families equal the uniquely normalised real solution of
Theorem~\ref{thm:W2} there, so the identity theorem patches them.
Their union $N_\eta=\bigcup_x U_x$ is connected and excludes $\eta$.
The preceding bounds give zero-freeness and, by Rouch\'e's theorem,
exactly one preimage of $1$ in $B(0,r)$, necessarily $0$.
They also bound the functions and their reciprocals uniformly.
If $b\to\eta$ inside $N_\eta$ and $b\in U_x$, then
$|b-x|<(\eta-x)/4$ forces $x\to\eta$; the last displayed bound and
Corollary~\ref{cor:sewn-cusp-local} give the stated convergence.
Cauchy's formula gives convergence of derivatives on smaller disks.
\end{proof}

\begin{remark}\label{rem:open}
The open-set hypothesis in (3) is implied by membership at a single base
$b'_*$ with strict inequality $|\operatorname{Im}Q-h/2|<1$, provided the limits $\kappa_{b'}(z)\to L(b')$, $L_2(b')$ as
$\operatorname{Im}z\to\pm\infty$ are locally uniform in $b'$.  Indeed, at a
complex base these limits give $\Rr$- and $S$-representations near
$\pm\ii\infty$ with bounded periodic corrections, provided there is no winding.
In detail, $H(q)=\lambda^{P}$, with $q=\eee^{2\pi\ii z}$, satisfies
$|H|=o(|q|^{-\arg\lambda/2\pi})$ with $\arg\lambda<2\pi$, so it extends
holomorphically to $q=0$.  A zero of $H$ at $q=0$ is exactly winding.  We take
$\operatorname{Im}z$ large enough to avoid the other zeros of $H$ (points where
$\kappa=L$).  The winding number is an integer that varies continuously, so it
persists.  The $S$-side is the same; $\tau^{-1}$ has no critical points.
On the remaining strip, which is compact modulo $1$, the representations of
$\kappa_{b'_*}$ continue to nearby $b'$ because $\kappa_{b'}\to\kappa_{b'_*}$
uniformly there, and the height condition is open.  Local uniformity of the
defining limits in the base is exactly the kind of statement a construction
of the family, such as Paulsen's~\cite{Paulsen2019}, provides.
\end{remark}

\noindent
Consequently, nothing quantitative about $\kappa$ remains to be checked.  It
is enough that the base-continued family has its two regular representations,
with the seam half a period below the real axis, on some open set of bases
near $(b_1,\eta)$.  That is how \texttt{\_cbuild.py} computes it, and the
ladders are consistent with it.  A proof requires the precise definition and
uniqueness theorem of~\cite{Paulsen2019}, which we have not been able to
consult.

\subsection{Numerics up to $\eta-b=3\cdot10^{-5}$}

Because $T$ needs no ladder, the invariants can be computed essentially at
the parabolic value.  The modes are measured on the line
$\operatorname{Im}z=-h_2/2+Y_0$, $h_2=2\pi/\log\lambda_2$, close to the gate
carrying the positive modes, so that the precision needed is independent of
$h$; results at $Y_0=1,1.5,2,2.5,3$ agree to all printed digits, and at
$b=1.3$ they agree with sampling on $\mathbb R$ at $40$ digits.

\begin{center}
\small\setlength{\tabcolsep}{4pt}
\begin{tabular}{lllllll}
\toprule
$\lambda$ & $\eta-b$ & $|\tau_1|/|B_1|$ & $\delta_1/\varepsilon^2$ & $\arg\tau_1$ & $|t_2/t_1^2|$ & $|t_3/t_1^3-B_3/B_1^3|$\\
\midrule
$0.5$ & $\num{9.04e-2}$ & $0.99601457$ & $0.008295$ & $1.581415$ & $1.555683$ & $0.148$\\
$0.8$ & $\num{1.21e-2}$ & $0.99952725$ & $0.009494$ & $1.585920$ & $1.572187$ & $0.0177$\\
$0.9$ & $\num{2.84e-3}$ & $0.99989062$ & $0.009853$ & $1.586391$ & $1.573905$ & $0.00410$\\
$0.95$ & $\num{6.87e-4}$ & $0.99997362$ & $0.010028$ & $1.586499$ & $1.574298$ & $0.000989$\\
$0.98$ & $\num{1.08e-4}$ & $0.99999587$ & $0.010131$ & $1.586528$ & $1.574403$ & $0.000155$\\
$0.99$ & $\num{2.67e-5}$ & $0.99999897$ & $0.010165$ & $1.586532$ & $1.574418$ & $0.0000385$\\
\midrule
limit & $0$ & $1$ & $\approx0.01020$ & $1.586533$ & $1.574423$ & $0$\\
\bottomrule
\end{tabular}
\end{center}

All three invariants approach their parabolic targets.  The third mode
converges to $B_3/B_1^3=-2.254799-0.024409\ii$, so it is a meaningful and
verified comparison.

\paragraph{The deficit law.}
The natural variable for the approach is not $\varepsilon$ but the product
of the two logarithmic multipliers,
\[
 p:=|\log\lambda|\,\log\lambda_2=\varepsilon^2\bigl(1-\tfrac{\varepsilon}3+O(\varepsilon^2)\bigr).
\]
This is the analytic parameter $\hat\epsilon$ of the
Marde\v{s}i\'c--Roussarie--Rousseau normal form, up to a factor analytic and
nonzero at $0$: the multipliers there are $\exp(\mu_{0,\infty})$ with
$\mu_0\mu_\infty=-4\hat\epsilon/(1-a^2\hat\epsilon)$.  Twelve bases with
$1-\lambda=k/200$, $k=1,\ldots,12$, fitted by polynomials of degree
$2,3,4$ in $p$, give coefficients that agree to eleven digits:
\begin{equation}\label{eq:kappa1}
\begin{aligned}
 \log\frac{\tau_1(b)}{B_1\eee^{2\pi\ii a}}&=\kappa^{(1)}p+O(p^2),&
 \kappa^{(1)}&=-0.0101990063459-0.0132509561825\ii,\\
 \log\frac{t_2/t_1^2}{B_2/B_1^2}&=\kappa^{(2)}p+O(p^2),&
 \kappa^{(2)}&=-0.0306585435-0.0358138682\ii .
\end{aligned}
\end{equation}
The coefficients of $p^2$ are about $4\cdot10^{-4}$ in modulus.  The constant
terms vanish to $10^{-16}$, and in $\varepsilon$ there is no linear term.
Consequently
\[
 \delta_1=C\,p+O(p^2),\qquad C=-\operatorname{Re}\kappa^{(1)}=0.0101990063,
\]
and $\delta_\kappa/\delta_1\to\operatorname{Re}\kappa^{(2)}/\operatorname{Re}\kappa^{(1)}=3.00603$.
This settles Section~\ref{sec:deficit}.  The deficit is analytic of first order
in the unfolding parameter.  The exponents $1.74$ and $2$ and the constant
$0.01013$ reported there came from fitting a finite range in variables that
differ from $p$ at relative order $\varepsilon$.  Candidate closed forms such as
$1/(10\pi^2)$ are excluded.  By Corollary~\ref{cor:cn}, the Kneser-side
$\hat c_n$ inherits~\eqref{eq:kappa1} up to $O(\Lambda)$, which is invisible at
every order in $p$.  The first-order coefficients of $\log(\tau_n/(B_n\eee^{2\pi\ii na}))$ for
$n=1,2,3$ are
\[
 -0.0101990063-0.0132509562\ii,\quad
 -0.0510565562-0.0623157805\ii,\quad
 -0.1400443481-0.1692244639\ii
\]
(script \texttt{deficit\_modes.py}).  They are not linear in $n$, so the
first-order effect of the unfolding is not a translation of Fatou coordinates.
A derivation of $\kappa^{(1)}$ is left open, and it is not a routine
perturbation computation.  Write the attracting coordinate as
$\Psi_p+\sum_kF_p\circ f_p^k$.  The model time is
$\Psi_p=\nu_p\log u-\sum_{m\ge1}s_m(p)/(mu^m)$, up to constants, with
$s_m=A_1u_1^m+A_2u_2^m$ analytic in $p$ and $s_3\propto p$.  Differentiating
term by term in $p$ at $p=0$ diverges: along $u_k\approx-2/k$ the parameter
sensitivity of the orbit grows like $k$ while $F'_0(u_k)$ decays like $1/k$,
and the $s_3/u^3$ term grows like $k^3$.  So $\kappa^{(1)}$ is a regularised,
resurgent quantity of the unfolding (compare the analytic dependence of the
unfolded modulus on the canonical parameter in~\cite{MRR2004,ChristopherRousseau}).
Its closed-form evaluation is a natural question for that theory.


\begin{remark}[On the definition and uniqueness in~\cite{Paulsen2019}]\label{rem:paulsen}
In~\cite{Paulsen2019}, $\kappa_b$ is defined by analytic continuation in $b$ of
Kneser's solution $\kappa_b(z_0)$, $\operatorname{Re}z_0>-2$, from $(\eta,\infty)$
(Prop.~3 there).  That proposition asserts real-analytic dependence of
Kneser's conformal uniformisation on the real base without a separate
parameter-dependence proof; Proposition~\ref{prop:exterior-real-parameter}
supplies such a proof locally along the real exterior ray, and
Corollary~\ref{cor:exterior-uniform-tails} gives uniform vertical tails
there.  Transporting those estimates through a complex corridor to the
Shell--Thron boundary remains open.  Off that boundary, Lemma~1 of
that paper asserts two
regular representations $\psi_1^{-1}(\rho_1)$ and
$\psi_2^{-1}(\rho_2)$ from the limits at $\pm\ii\infty$.
We impose their degree-one form explicitly in
Definition~\ref{def:class}; Lemma~\ref{lem:end-spectrum} records
the integer end label that a fixed-point limit alone leaves open.  The continuation across the Shell--Thron boundary is
argued in~\cite[\S5]{Paulsen2019} by the continuity of a numerical scheme; it is
not proved there.

The uniqueness statement~\cite[Prop.~2]{Paulsen2019} requires only
analyticity on $\operatorname{Re}z>-2$, $F(0)=1$ and the two limits.  As stated,
it is false, already at $b=\eee$, by the following exact existence argument.
Let $K=\operatorname{sexp}_{\eee}$ be Kneser's real-normalised solution and
$D_0=\{|z|<\tfrac12\}$.  Since $K(0)=1$ and $K$ is nonconstant and
holomorphic, the open-mapping theorem gives $\delta>0$ such that
$B(1,\delta)\subset K(D_0)$.  Decrease $\delta$ so that
$\eee(1-\eee^{-\delta})>\delta$.  For every real $t$,
\[
 \exp(B(t,\delta))\supset
 B\bigl(\eee^t,\eee^t(1-\eee^{-\delta})\bigr),
\]
as follows by applying the principal logarithm and
$|\log(1+v)|\le-\log(1-|v|)$.  Put $u_n=K(n)$; then $u_0=1$ and
$u_{n+1}=\eee^{u_n}$.  Induction using the displayed inclusion gives
$B(u_n,\delta)\subset K(D_0+n)$ for every $n$, and indeed
\[
 B\bigl(u_{n+1},u_{n+1}(1-\eee^{-\delta})\bigr)
 \subset K(D_0+n+1).
\]
Choose $n$ so large that the last radius exceeds $1+\pi/2$, and let
$k=\lceil\eee^{u_{n+1}}/(2\pi)\rceil$.  For large $n$,
$0\le\log(2\pi k)-u_{n+1}<1$.  Hence
$w=\log(2\pi k)+\ii\pi/2$ lies in $K(D_0+n+1)$.  Choose
$z_*\in D_0+n+1$ with $K(z_*)=w$ and set $c=z_*+2$.
Then $\operatorname{Re}c>0$ and, exactly,
\[
 K(c-1)=\eee^w=2\pi\ii k,\qquad
 K(c)=\eee^{2\pi\ii k}=1.
\]
Thus $F(z)=K(z+c)$ is analytic on $\operatorname{Re}z>-2$, obeys
$F(z+1)=\eee^{F(z)}$, has $F(0)=1$, and has the same two vertical limits
as $K$.  Yet $F(-1)=2\pi\ii k\ne0=K(-1)$.
The construction works for every real $b>\eta$.  Put $a=\log b>1/\eee$
and $u_n=K_b(n)\to\infty$, since $\eee^{ax}>x$ for $x>0$.
Choose $N$ with $a u_{N+1}>1$ and then
$\delta>0$ so that $B(u_N,\delta)\subset K_b(B(N,\tfrac12))$ and
$u_{N+1}(1-\eee^{-a\delta})>\delta$.  The inclusion
\[
 E_b(B(t,\delta))\supset
 B\bigl(\eee^{at},\eee^{at}(1-\eee^{-a\delta})\bigr)
\]
propagates the fixed $\delta$-disks and yields image disks of radii
$u_{n+1}(1-\eee^{-a\delta})\to\infty$.  For large
$n\ge N$ choose $k=\lceil a\exp(a u_{n+1})/(2\pi)\rceil$.
The target $w_k=a^{-1}(\log(2\pi k/a)+\ii\pi/2)$ lies in the
resulting image disk.  Then
$E_b^2(w_k)=1$ and $E_b(w_k)=2\pi\ii k/a\ne K_b(-1)=0$.
Choose $z_k\in B(n+1,\tfrac12)$ with $K_b(z_k)=w_k$, put $c_k=z_k+2$,
and set $F_k(z)=K_b(z+c_k)$.  Since $\operatorname{Re}c_k>0$, each $F_k$
meets the hypotheses of Proposition~2, but $F_k(-1)=2\pi\ii k/a$.
The values $k$ are unbounded, giving infinitely many distinct solutions.
The gap in~\cite[Prop.~2]{Paulsen2019} is the last step of its proof,
which infers $C=0$ from $F_1(C)=1$.  For real $b>\eta$ the uniqueness
theorem of~\cite{PaulsenCowgill2017} is not affected: it also assumes
analyticity on $\mathbb C\setminus(-\infty,-2]$ and $F(\bar z)=\overline{F(z)}$,
both of which this $F$ violates.  Indeed $c$ is nonreal, and $K$ cannot
extend holomorphically through $-2$, since $K(-1)=0$ whereas an exponential
never vanishes; hence $F$ is singular at $-2-c$ off the standard cut.
For complex bases, some root-selection
condition must be added to~\cite[Prop.~2]{Paulsen2019}.  The seam-height
condition of Definition~\ref{def:class} is such a condition, and
Theorem~\ref{thm:char} shows that it suffices near $(1,\eta)$.
\end{remark}

\subsection{Identification with an independent Kneser-type implementation}\label{sec:fatou}
The hypothesis of Theorem~\ref{thm:bv} can be tested against a completely independent
computation of the base-continued family.  Sheldon Levenstein's PARI/GP
program \texttt{fatou.gp} builds the ``merged'' two-fixed-point solution for
complex bases.  At each base it constructs, independently, the solution with a
$\theta$-mapping at each of the two fixed points: the two regular representations
of~\cite[Lemma~1]{Paulsen2019}.  Its author regards this as the continuation of
Kneser's solution in the base.  We use the unmodified program
(sha256 \texttt{70559daf\ldots}): published reference values with $100$
verified digits at $b=0.8+0.4\ii$ and $b=1+\ii$, and our own runs at
\texttt{\textbackslash p 40} at $b=1.2+0.1\ii$ and $b=1.3+0.2\ii$, near the
interval $(1,\eta)$.  In those runs $\mathrm{sexp}(0)=1$ to $10^{-53}$, and the
functional equation holds to $10^{-54}$.

We solved the sewing equation of Theorem~\ref{thm:W2} directly at these complex
bases (script \texttt{weld\_complex.py}).  $\Rr$ is the regular solution at the
attracting fixed point and $S$ the one at the repelling fixed point; the seam
lies in the band of holomorphy of $T=\Rr^{-1}\circ S$, and the root is selected
as in Definition~\ref{def:class}.  Table~\ref{tab:fatou} shows the result.  At
the three bases where the band is not too narrow, the sewing solution agrees
with the independent value \emph{to the full working precision of $40$ digits}.
That is $27$ to $35$ orders of magnitude below the size of the separation
$|\kappa-\Rr|$ itself.  At $1+\ii$ the band is narrowest and the agreement
($2\cdot10^{-25}$) is limited by the number of samples.  The neighbouring roots
$c_{\pm1}$ reproduce $\Rr$ instead, and miss by the full separation.  The
$\theta$-mapping builder used for the ladders does not converge at $0.8+0.4\ii$,
whereas the sewing iteration converges in a few dozen steps.

\begin{table}[htbp]
\centering
\caption{Sewing solution versus the independent \texttt{fatou.gp} values of
$\kappa_b(\tfrac12)$ (working precision $40$ digits).}
\label{tab:fatou}
\small\setlength{\tabcolsep}{5pt}
\begin{tabular}{lllll}
\toprule
$b$ & $|\Lambda|$ & $|\kappa_b(\tfrac12)-\Rr_b(\tfrac12)|$ & $|\Kk^W_b(\tfrac12)-\kappa_b(\tfrac12)|$ & $|\hat c_1|$\\
\midrule
$1.2+0.1\ii$ & \num{1.91e-11} & \num{6.91e-13} & \num{4.6e-40} & $0.0869956$\\
$1.3+0.2\ii$ & \num{2.92e-11} & \num{1.81e-12} & \num{8.3e-41} & $0.0876095$\\
$0.8+0.4\ii$ & \num{1.05e-3} & \num{9.86e-5} & \num{5.5e-39} & $0.0856072$\\
$1+\ii$ & \num{2.13e-2} & \num{4.41e-3} & \num{2.0e-25} & $0.0873850$\\
\bottomrule
\end{tabular}
\end{table}

The agreement identifies the sewing solution with an independent implementation
of the two-sided regular representation, at four bases, two of them near
$(1,\eta)$.  It does \emph{not} test the analytic continuation across the
Shell--Thron boundary: both programs work at one base at a time.  That
continuation is the one remaining hypothesis (Section~\ref{sec:brjuno}).

\subsection{A pointwise test near the Shell--Thron boundary}\label{sec:bndtest}
The analytic continuation of Kneser's solution across the Shell--Thron boundary is
argued in~\cite[\S5]{Paulsen2019} only through the continuity of a numerical
scheme.  In particular, continuity of each finite quadrature approximant
does not imply continuity, let alone holomorphy, of its limit without
locally uniform convergence in the base parameter.  The transfer in
\cite[Prop.~3]{Paulsen2019} of fixed-point limits from real to complex
bases likewise requires such uniformity.  We tested whether the pointwise
two-fixed-point constructions on the two sides are numerically compatible
near the irrational boundary point
$\lambda_c=\eee^{\ii\phi}$, $\phi=2\pi(\sqrt2-1)/5$, where
$b_c=\exp(\lambda_c\eee^{-\lambda_c})\approx1.5217+0.0148\ii$.  Along the normal
$\lambda=r\eee^{\ii\phi}$ we computed the pointwise output
$F_b^{\mathrm{GP}}(\tfrac12)$ of \texttt{fatou.gp}
(\texttt{\textbackslash p 24}).  We used three bases inside
($r=0.98,0.99,0.995$, one attracting fixed point) and three outside
($r=1.005,1.01,1.02$, both fixed points repelling).  So the two sides are
computed with structurally different representations.  A degree-$4$ polynomial
in $b$ fitted to five of the six values predicts the sixth to
$10^{-12}$--$10^{-13}$, uniformly on both sides.  The quadratic through the three
inside values extrapolates to the outside values with errors
$1.3\cdot10^{-8}$, $3.2\cdot10^{-8}$ and $1.1\cdot10^{-7}$, growing with the
distance as the cubic truncation term of a smooth function should.  The reverse
extrapolation gives the mirror-image errors.  The errors are consistent with
a smooth join of these pointwise constructions at the tested scales.
Finite sampling cannot exclude a smaller
jump, a higher-order nonanalytic term, or failure of continuation elsewhere
on the boundary.  Nor does the pointwise algorithm identify its outputs
with the base-continued family $\kappa$.  This is numerical evidence
compatible with the continuation claimed in~\cite{Paulsen2019}, not a proof.

We repeated the six-base calculation at five noninteger heights, using the
unmodified \texttt{fatou.gp} at $50$-digit precision
(\nolinkurl{boundary_multiheight.py}; raw values in
\nolinkurl{data/boundary-multiheight.json}).  Table~\ref{tab:bnd-multiheight}
reports the largest leave-one-out error when a degree-$4$ polynomial in $b$
through five bases predicts the sixth.  It also reports the difference at
$b_c$ between quadratics fitted separately to the three inside and three
outside values.  Recomputing the two closest bases at $70$-digit precision
changed all ten values by less than $5\cdot10^{-50}$.
The five tested heights form only a finite set, and sampling at one
boundary point cannot verify condition (P2).

\begin{table}[htbp]
\centering
\caption{Finite two-sided boundary comparison at five noninteger heights.
The quadratic gap is an extrapolation diagnostic, not a bound on a jump.}
\label{tab:bnd-multiheight}
\begin{tabular}{ccc}
\toprule
$z$ & max. degree-$4$ leave-one-out error & two-sided quadratic gap at $b_c$\\
\midrule
$1/6$ & $1.21\cdot10^{-12}$ & $4.25\cdot10^{-9}$\\
$1/3$ & $1.62\cdot10^{-12}$ & $6.64\cdot10^{-9}$\\
$1/2$ & $1.45\cdot10^{-12}$ & $7.02\cdot10^{-9}$\\
$2/3$ & $9.68\cdot10^{-13}$ & $5.57\cdot10^{-9}$\\
$5/6$ & $4.28\cdot10^{-13}$ & $2.89\cdot10^{-9}$\\
\bottomrule
\end{tabular}
\end{table}

\Needspace{22\baselineskip}
\begin{proposition}[Local base analyticity on the real exterior ray]
\label{prop:exterior-real-parameter}
For every real $b_0>\eta$, Kneser's real-normalised solution admits a
jointly holomorphic extension in $(b,z)$ to $\Delta\times W$, where
$\Delta$ is a complex disc about $b_0$ and $W$ is a complex
neighbourhood of $[0,1]$.  These local extensions agree on overlap,
so the classical exterior branch has a jointly holomorphic germ along
the entire real ray $(\eta,\infty)$.  In particular, every compact
subinterval of that ray has a common complex base neighbourhood and a
common complex height neighbourhood of $[0,1]$ on which this family is
jointly holomorphic.  The assertion supplies no
corridor reaching a non-cusp Shell--Thron boundary arc.
\end{proposition}

\begin{proof}
Write $a=\log b$ and choose the principal fixed-point pair $L_+(b)$ and
$L_-(b)$, conjugate when $b>\eta$ and characterised there by
$0<a\operatorname{Im}L_+<\pi$.  They are simple fixed
points and hence holomorphic in $b$ near $b_0$.  Set
$u_b=(L_++L_-)/2$, $v_b=(L_+-L_-)/(2\ii)$, and
$\gamma_b(t)=u_b+\ii v_b t$ for $-1\le t\le1$.
At a real base, $u_b,v_b\in\mathbb R$, $v_b>0$, and
\[
 E_b(\gamma_b(t))=r_b\eee^{\ii\theta_b t},\qquad
 r_b=\eee^{a u_b}=|L_+|,\quad
 \theta_b=a v_b\in(0,\pi).
\]
Since $v_b=r_b\sin\theta_b$, the multipliers satisfy
$|E_b'(L_\pm)|=a r_b=\theta_b/\sin\theta_b>1$.
The straight chord $\gamma_b$ and this circular arc meet only at
$L_\pm$ and bound the initial region
$H_b=\{\operatorname{Re}w\ge u_b,\ |w|\le r_b\}\setminus\{L_+,L_-\}$.
This is Kneser's initial region as described in
\cite[\S4, Lem.~4]{TrappmannKouznetsov2011}.
The arcs depend holomorphically on $b$ through their real-analytic
parameter $t$.  They remain disjoint away from their endpoints after
shrinking a disc $\Delta$ about $b_0$; at the endpoints their tangents
remain distinct because the fixed-point multipliers are nonreal.
The bounded Jordan interiors vary continuously, so their union in
$\Delta\times\mathbb C$ is open; the ambient coordinates $(b,w)$ make
its projection a holomorphic family.  Identify thin analytic collars
of the two arcs by
$(b,w)\mapsto(b,E_b(w))$.  The map is injective on the collar after
shrinking $\Delta$, since it is injective on the compact chord at
$b_0$ and its derivative never vanishes.  This gives a Hausdorff
holomorphic family of cylinders with one end at each $L_\pm(b)$.

The ends can be filled holomorphically.  Indeed, parameter-dependent
Koenigs coordinates $\sigma_{\pm,b}$, normalised by
$\sigma_{\pm,b}(L_\pm)=0$ and $\sigma_{\pm,b}'(L_\pm)=1$,
linearise $E_b$ near the two repelling fixed points.  If
$\ell_{\pm,b}=\log E_b'(L_\pm(b))$ is
chosen continuously, then
$\exp(2\pi\ii\log\sigma_{\pm,b}/\ell_{\pm,b})$ (or its reciprocal)
is a puncture coordinate on the corresponding glued end.  Filling
both punctures yields a proper holomorphic submersion of compact
genus-zero surfaces over $\Delta$, just as in the collar construction
of Proposition~\ref{prop:weld-parameter}.

We also have a holomorphic marked section represented by $w=1$.
For real $b>\eta$, $E_b(x)>x$ for every real $x$: the difference is
positive at $x\to\pm\infty$ and has no real zero.  The adjacent real
intervals between the chord and its image are $(u_b,r_b)$, with
$r_b=E_b(u_b)>u_b$.  The forward endpoints $E_b^n(u_b)$ increase to
$+\infty$, since a finite limit would be a real fixed point.  For any
$x>0$ below $u_b$, forward iteration eventually enters $(u_b,r_b)$
or one of its boundary seams; for $x>r_b$, repeated real logarithms
eventually do the same.
Indeed, a monotone orbit staying on either side would have a finite
real fixed-point limit.  Thus these dynamical translates and their
seam collars cover the positive real axis.  In particular, $1$ is
in the interior of one such interval or on a seam.  A finite chain of
holomorphic dynamical charts represents it near $b_0$; at a seam the
analytic collar gives the same section.  The entire compact real
segment from $1$ to $b=E_b(1)$ is covered by finitely many of these
charts.

The two filled ends and this marked point are three disjoint
holomorphic sections.  Fischer--Grauert local triviality
\cite{FischerGrauert1965} and the fibrewise three-point M\"obius
normalisation give a jointly holomorphic coordinate $q_b$ on the
compactified cylinder, taking the upper end to $0$, the lower end to
$\infty$, and $w=1$ to $1$.  The upper puncture coordinate
$\exp(2\pi\ii\log\sigma_{+,b}/\ell_{+,b})$ shows that the logarithmic
lift $z=(2\pi\ii)^{-1}\log q_b$, with $z(1)=0$, gains $+1$ under the edge
identification by $E_b$.  On the real slice, reflection exchanges the
two ends and fixes the marked point, so $q_b(\bar w)=1/\overline{q_b(w)}$
in the real dynamical charts.  Thus $z$ is real there; its local
injectivity and $z(E_b(w))=z(w)+1$ make it strictly increasing along
the positive axis.  Inverting on the finite
chart chain covering the real segment from $1$ to $b$ gives a jointly
holomorphic superfunction on $\Delta\times W$, after shrinking both.

For real $b$, the lift restricted to $H_b$ is injective and its integer
translates cover the universal-cover plane.  Choosing any $d\in H_b$
and subtracting its value gives the initial-region criterion
of~\cite[Thm.~1]{TrappmannKouznetsov2011}; the injectivity and tiling
assertions for the Kneser generalisation are proved in
\cite[Thm.~5]{TrappmannKouznetsov2011}.  We normalise at $d$ here,
since $1$ need not itself belong to $H_b$.
Their Abel coordinates differ only by a constant on $H_b$, and
continuation along the real dynamical charts to $w=1$ makes that
constant zero.  Thus the inverse just
constructed is the classical real-normalised Kneser branch.
The marking fixes the local family uniquely, and the identity theorem
patches the parameter discs on overlaps.  A finite subcover of any
compact real base interval, followed by intersection of its finitely
many height neighbourhoods, gives the final uniformity assertion.
\end{proof}

\begin{proposition}[Persistence of the exterior initial curve near the cusp]
\label{prop:exterior-curve-collar}
Put, for small $\theta\ne0$,
\[
 r(\theta)=\eee^{\theta\cot\theta},\qquad
 a(\theta)=\frac{\theta\eee^{-\theta\cot\theta}}{\sin\theta},
 \qquad b(\theta)=\eee^{a(\theta)},\qquad
 L_\pm(\theta)=r(\theta)\eee^{\pm\ii\theta}.
\]
For every fixed $M>0$ there is $p_0>0$ such that, whenever
$\theta=p+\ii q$, $0<p<p_0$ and $|q|\le Mp$, the chord
$\gamma_\theta(t)=r(\theta)(\cos\theta+\ii t\sin\theta)$ and its
image $E_{b(\theta)}(\gamma_\theta(t))=r(\theta)\eee^{\ii\theta t}$
are individually injective and meet only at their common endpoints
$L_\pm(\theta)$.  Thus the exterior initial Jordan curve persists
geometrically in this complex parameter cone.

There is moreover a real-analytic function
$q_c(p)=p^2/6+O(p^3)$ such that, for $0<p<p_0$ and
$0\le q<q_c(p)$, both fixed-point multipliers have modulus greater
than one.  The upper multiplier has modulus one at $q=q_c(p)$;
this is a non-cusp Shell--Thron boundary arc.  Therefore the
unmarked two-end quotient surfaces of
Proposition~\ref{prop:exterior-real-parameter} exist locally throughout
this exterior wedge, arbitrarily close to that boundary arc.  The map
$\theta\mapsto b(\theta)$ is injective on a sufficiently small such
wedge, so its image is a genuine complex base corridor attached to
the real exterior ray.  This
geometric assertion does not supply a holomorphic section represented
by $w=1$ on a common parameter corridor or a common regular height
domain there.
\end{proposition}

\begin{proof}
The identities $a r\sin\theta=\theta$ and
$\log r=\theta\cot\theta$ show that $L_\pm$ are fixed points and
give the displayed image arc.  The functions $a(\theta)$ and
$b(\theta)$ are even and holomorphic at $0$, with
$b(\theta)=\eta\bigl(1+\theta^2/(2\eee)+O(\theta^4)\bigr)$;
in particular $b'(\theta)\ne0$ for small $\theta\ne0$.  More strongly,
$b(\theta)=B(\theta^2)$ with $B'(0)=\eta/(2\eee)\ne0$.
The map $B$ is locally injective, as is $\theta\mapsto\theta^2$ in
$\operatorname{Re}\theta>0$; hence $b$ is injective on the small wedge.

An intersection at an interior arc point would require the number
\[
 s_\theta(t)=
 \frac{\eee^{\ii\theta t}-\cos\theta}{\ii\sin\theta}
 \qquad(-1<t<1)
\]
to be real.  Taylor expansion with its endpoint zeros factored out
gives, uniformly for $-1\le t\le1$,
\[
 \eee^{\ii\theta t}-\cos\theta-\ii t\sin\theta
  =(1-t^2)\theta^2\bigl(\tfrac12+O(|\theta|)\bigr),
 \qquad
 \operatorname{Im}s_\theta(t)
  =-(1-t^2)\bigl(\tfrac p2+O(|\theta|^2)\bigr).
\]
The error is uniform even as $t\to\pm1$: every term of order at
least two in the power series has the factor $1-t^2$, with the
remaining polynomial bounded on $[-1,1]$.  In the stated cone,
$|\theta|^2\le(1+M^2)p^2$, so the last imaginary part is strictly
negative for $-1<t<1$ and small $p$.  The chord is injective since
$\sin\theta\ne0$, and the arc is injective since
$\eee^{\ii\theta(t_1-t_2)}=1$ with $|t_1-t_2|\le2$ and small
$|\theta|$ forces $t_1=t_2$.  The same expansion gives distinct
endpoint tangents, hence a Jordan curve.

The multipliers are
$\lambda_\pm(\theta)=a(\theta)L_\pm(\theta)
 =\theta\cot\theta\pm\ii\theta$.  For $\theta=p+\ii q$ near zero,
\[
 |\lambda_+|^2=1+\frac{p^2}{3}-2q
       +O(p^3+p|q|+q^2),\qquad
 |\lambda_-|^2=1+\frac{p^2}{3}+2q
       +O(p^3+p|q|+q^2).
\]
At $q=0$ both are repelling, while
$\partial_q|\lambda_+|^2=-2+O(p+|q|)$.  The real implicit function
theorem gives a unique nearby zero $q_c(p)=p^2/6+O(p^3)$, and the
signs above show that both multipliers are repelling for
$0\le q<q_c(p)$.  Since
$a(\theta)=\lambda_+(\theta)\eee^{-\lambda_+(\theta)}$,
the locus $|\lambda_+|=1$ is precisely the stated
Shell--Thron boundary parameterisation.  The collar gluing and
repelling-end filling in Proposition~\ref{prop:exterior-real-parameter}
apply locally at every point of the open wedge.
\end{proof}

\begin{proposition}[A certified marked exterior path]
\label{prop:exterior-marked-path}
Set $\theta(q)=\frac1{10}+\ii q$ and use $b(\theta)$, $L_\pm(\theta)$
and the initial region $H_\theta$ of
Proposition~\ref{prop:exterior-curve-collar}.  There is a unique
$q_c\in(0,1667/10^6)$ with
$|\lambda_+(\theta(q_c))|=1$.  For every $0\le q\le q_c$,
\[
 E_{b(\theta(q))}^{\circ28}(1)\in\operatorname{int}H_{\theta(q)}.
\]
For $0\le q<q_c$ both fixed points are repelling, and the exterior
Kneser germ normalised by $F(0)=1$ analytically continues along the
embedded base path $q\mapsto b(\theta(q))$.  The result concerns the
germ at height $z=0$; it does not give a single regular height domain
containing $[0,1]$ uniformly as $q\uparrow q_c$.
\end{proposition}

\begin{proof}
The following strict enclosures are certified with outward-rounded
Arb complex ball arithmetic by the reproducible script
\texttt{docs/certify\_exterior\_mark.py}.  It uses only the rational
inputs $p=1/10$, $Q=1667/10^6$, $28$ iterates, and $1000$ equal
subintervals of $[0,Q]$; the recorded run used
\texttt{python-flint} 0.9.0 with FLINT 3.6.0.  For $0\le q\le Q$ it checks that the arc
coordinate
$s_\theta(t)=(e^{\ii\theta t}-\cos\theta)/(\ii\sin\theta)$
has $\partial_t\operatorname{Re}s_\theta(t)>0.98$ for $-1\le t\le1$,
with $\operatorname{Re}s_\theta(-1/10)<-0.08$ and
$\operatorname{Re}s_\theta(1/10)>0.08$.
The estimate behind the remaining whole-arc check is elementary:
writing
\[
 e^{\ii\theta t}-\cos\theta-\ii t\sin\theta
 =(1-t^2)\theta^2H(\theta,t),
 \qquad
 |H(\theta,t)-\tfrac12|
 \le\frac{e^R-1-R}{2R},\quad R=0.101,
\]
the termwise interpolation remainder and
$|\sin\theta/\theta-1|\le\sinh(R)/R-1$ imply
$\operatorname{Im}s_\theta(t)<-0.04(1-t^2)$ for $-1<t<1$.
Indeed, after division by $1-t^2$ the order-$n$ interpolation
polynomial is a sum of at most $n/2$ powers of $t$, each of modulus
at most one; this gives the stated bound on $H$.
Thus the image arc is a graph strictly below the chord in the
$s$-plane.  On the horizontal interval $[-0.01,0.04]$ its height is
below $-0.0396$.

Let $w_{28}(\theta)=E_{b(\theta)}^{\circ28}(1)$ and
$\zeta=(w_{28}/r-\cos\theta)/(\ii\sin\theta)$.
The same script encloses $\zeta$ on every parameter subinterval and
verifies $-0.01<\operatorname{Re}\zeta<0.04$ and
$-0.03<\operatorname{Im}\zeta<-0.01$.  Hence $w_{28}$ lies strictly
between the two arcs.  It also verifies that
$|\lambda_+|^2-1$ is positive at $q=0$, negative at $q=Q$, and has
derivative less than $-1$ throughout $[0,Q]$; that
$|\lambda_-|^2>1.003$ there; and that
$\operatorname{Im}(d b(\theta(q))/dq)>0.05$.
These strict inequalities prove the root, repelling and embedded-path
claims.

For $q<q_c$, the two repelling ends and the analytic collars give a
holomorphic family of quotient spheres locally in $q$.
The point $w_{28}(\theta)$ is a holomorphic section in the initial
region, so it marks the physical point $w=1$ through the finite
iterate $E_b^{\circ28}$; that iterate has nonzero derivative at $1$.
After three-point normalisation choose the logarithmic Abel lift with
value $28$ at $w_{28}$, and pull it back by the local inverse of
$E_b^{\circ28}$ to obtain value $0$ at $w=1$.  Local inversion of
this Abel coordinate gives a jointly holomorphic normalised
superfunction germ at $z=0$.
At $q=0$ it is the classical Kneser germ by
Proposition~\ref{prop:exterior-real-parameter}; the embedded path
continues that germ uniquely.
\end{proof}

\begin{proposition}[A prescribed Brjuno endpoint for the marked path]
\label{prop:exterior-brjuno-path}
Let
\[
 \alpha=\frac{\sqrt2-1}{26},\qquad
 \lambda_c=\eee^{2\pi\ii\alpha}.
\]
There is a unique solution $\theta_c=p_c+\ii q_c$ of
$\lambda_+(\theta_c)=\lambda_c$ in the coordinate square of radius
$10^{-8}$ centred at
\[
 \theta_0=0.1000435842+0.0016682119\,\ii.
\]
The number $\alpha$ is Brjuno.  Along the embedded path
$\theta(q)=p_c+\ii q$, $0\le q<q_c$, the two fixed points remain
repelling and $E_{b(\theta(q))}^{\circ28}(1)$ remains strictly inside
$H_{\theta(q)}$, also at $q=q_c$ in the latter geometric assertion.
Consequently the classical real Kneser germ has a marked analytic
continuation along an exterior path ending at a specified non-cusp
Brjuno boundary point.  No common regular $z$-domain containing
$[0,1]$ is asserted.
\end{proposition}

\begin{proof}
The number $\alpha$ is a root of $676X^2+52X-1$.  For a rational
$m/n$ near it, that integer polynomial has nonzero value of modulus
at least $1/n^2$ at $m/n$.  Its derivative is bounded nearby, hence
$|\alpha-m/n|\ge c/n^2$ for some $c>0$.  The denominators of its
continued-fraction convergents therefore satisfy
$n_{k+1}\le c^{-1}n_k$, while they grow at least as fast as Fibonacci
numbers.  Thus $\sum_k(\log n_{k+1})/n_k<\infty$, the Brjuno condition.

The remaining claims are certified with outward-rounded Arb complex
balls by \path{docs/certify_brjuno_exterior_mark.py}.  Put
$f(\theta)=\theta\cot\theta+\ii\theta-\lambda_c$ and
$T(\theta)=\theta+\ii f(\theta)$.  On the stated closed square the
script verifies
\[
 |f(\theta_0)|<3\cdot10^{-11},\qquad
 \sup|T'(\theta)|<0.07,\qquad
 3\cdot10^{-11}+0.07\sqrt2\cdot10^{-8}<10^{-8}.
\]
The contraction theorem gives the unique root $\theta_c$ in the
square.  In particular $0<q_c<0.00166823$.

The same script subdivides $[0,0.00166823]$ into $1000$ rational
intervals and encloses the entire coordinate box
$|p-0.1000435842|\le10^{-8}$.  Uniformly there, it verifies the
strict arc-graph and multiplier inequalities used in
Proposition~\ref{prop:exterior-marked-path}, with the whole-arc bound
$\operatorname{Im}s_\theta(t)<-0.04(1-t^2)$ obtained now using
$|\theta|<0.102$.  It verifies for
$\zeta=(E_{b(\theta)}^{\circ28}(1)/r-\cos\theta)/(\ii\sin\theta)$
that
\[
 -0.01<\operatorname{Re}\zeta<0.04,\qquad
 -0.03<\operatorname{Im}\zeta<-0.01.
\]
The arc is below $-0.0396$ on this horizontal range, so the marked
iterate lies strictly inside $H_\theta$.  Finally the script checks
$\partial_q|\lambda_+|^2<-1$, $|\lambda_-|^2>1.003$, and
$\operatorname{Im}(d b(p+\ii q)/dq)>0.05$ throughout the enclosing
box.  Since $|\lambda_+(\theta_c)|=1$, both multipliers are
repelling for $q<q_c$ and the base path is embedded.  The marked
quotient construction and local inversion from
Proposition~\ref{prop:exterior-marked-path} now apply along this path.
\end{proof}

\Needspace{13\baselineskip}
\begin{proposition}[A marked pullback cell along the Brjuno path]
\label{prop:exterior-pullback-cell}
For every $0\le q\le q_c$ on the path of
Proposition~\ref{prop:exterior-brjuno-path}, the inverse branch of
$E_{b(q)}^{\circ28}$ taking $w_{28}(q)$ to $1$ extends conformally
to a neighbourhood of $\overline{H_q}$.  Denote it by $\Phi_q$.
It depends continuously on $q$ up to $q_c$ and holomorphically in
the base parameter locally along the path.  Moreover
\[
 \Phi_q(E_{b(q)}(w))=E_{b(q)}(\Phi_q(w))
 \qquad(w\in\gamma_q([-1,1])).
\]
Consequently $D_q=\Phi_q(H_q)$ is a Jordan cell containing $w=1$;
its two boundary arcs are $\Phi_q(\gamma_q)$ and
$E_{b(q)}(\Phi_q(\gamma_q))$, and their endpoints are the original
fixed points $L_\pm(q)$.  This supplies a global physical pullback
of the marked initial region, including at the Brjuno endpoint.
\end{proposition}

\begin{proof}
The certificate \path{docs/certify_brjuno_exterior_mark.py} now also
encloses the singular orbit $E_{b(q)}^{\circ k}(0)$ for every
$0\le k\le28$ on all $1000$ parameter slabs.  In the chord
coordinate $\zeta=(w/r-\cos\theta)/(\ii\sin\theta)$ it verifies
\[
 -0.1<\operatorname{Re}\zeta<0.2,\qquad
 \operatorname{Im}\zeta>0.02
 \quad\bigl(w=E_{b(q)}^{\circ k}(0),\ 0\le k\le28\bigr).
\]
The closed cell $\overline{H_q}$ lies below the chord
$\operatorname{Im}\zeta=0$, so it avoids every one of these orbit
points with a uniform margin.  Choose a simply connected thin
neighbourhood $U_q$ of $\overline{H_q}$ which still avoids them.

Set $g_0(w)=w$.  Inductively choose a logarithm on $U_q$ so that
\[
 E_{b(q)}(g_j(w))=g_{j-1}(w),\qquad
 g_j(w_{28}(q))=E_{b(q)}^{\circ(28-j)}(1)
 \quad(1\le j\le28).
\]
For $j=1$, $g_0$ never vanishes on $U_q$.  If $g_{j-1}$ vanished at
$w\in U_q$, then $w=E_{b(q)}^{\circ(j-1)}(0)$, contrary to the
certificate.  Thus each logarithm exists on the simply connected
$U_q$.  It is univalent, and the induction gives
$E_{b(q)}^{\circ j}\circ g_j=\operatorname{id}$.  In particular
$\Phi_q=g_{28}$ is the asserted inverse branch.  The uniform
clearance and the analytic motion of the two boundary arcs allow the
$U_q$ and logarithm branches to be chosen locally uniformly in $q$,
including at $q_c$.  This proves the stated parameter dependence.

It remains to check the seam identity; it does not follow merely from
the equation $E_b^{\circ28}\circ\Phi_q=\operatorname{id}$.  At the
real base $q=0$, put $u=\gamma_0(0)$ and $r=E_{b(0)}(u)$.
The branch $\Phi_0$ is the real inverse of the strictly increasing
map $E_{b(0)}^{\circ28}$ on the real interval through $w_{28}(0)$.
Hence $\Phi_0(r)=E_{b(0)}(\Phi_0(u))$.
For arbitrary $(q,t)$ in the connected rectangle
$[0,q_c]\times[-1,1]$, both
$\Phi_q(E_{b(q)}(\gamma_q(t)))$ and
$E_{b(q)}(\Phi_q(\gamma_q(t)))$ are inverse images of
$E_{b(q)}(\gamma_q(t))$ under $E_{b(q)}^{\circ28}$.
They vary continuously, and $E_b^{\circ28}$ has nonzero derivative
everywhere.  Their equality set is therefore both open and closed
in the rectangle.  It contains $(0,0)$, so it is the whole rectangle.

The conformal image $D_q=\Phi_q(H_q)$ is Jordan and contains
$\Phi_q(w_{28}(q))=1$.  The seam identity identifies its second
boundary arc with the image under $E_b$ of its first.  At either
endpoint $L_\pm$, it says that $\Phi_q(L_\pm)$ is fixed by $E_b$;
since $E_b^{\circ28}(\Phi_q(L_\pm))=L_\pm$, it follows that
$\Phi_q(L_\pm)=L_\pm$.  These are the claimed endpoints.
\end{proof}

\Needspace{15\baselineskip}
\begin{proposition}[A quantitative flat-cylinder model]
\label{prop:exterior-flat-cylinder}
For every parameter on the marked Brjuno path, including its boundary
endpoint, the sewn two-end cylinder has a quasiconformal
parameterisation from the standard cylinder $\mathbb C/\mathbb Z$
with maximal dilatation less than $11/9$.  In the chord coordinate of
Proposition~\ref{prop:exterior-curve-collar}, it is induced by
\[
 z=x+20\ii y,\quad t=\tanh y,\quad
 \zeta_q(z)=t+x\bigl(s_{\theta(q)}(t)-t\bigr),
 \qquad 0\le x\le1,\quad y\in\mathbb R,
\]
where the sides $x=0$ and $x=1$ are identified by $E_{b(q)}$.
The bound is uniform as $q\uparrow q_c$ and controls the entire
punctured cylinder, not only a compact middle part.  It does not by
itself assert that the image of the flat seam meets a normalised
height circle only once.
\end{proposition}

\begin{proof}
The chord and image arc are disjoint Jordan edges, and the latter is
an increasing graph in the chord coordinate by the certified
inequalities of Proposition~\ref{prop:exterior-brjuno-path}.  Write
\[
 s_\theta(t)-t=(1-t^2)D_\theta(t),\qquad
 D_\theta(t)=\frac{\theta^2H(\theta,t)}{\ii\sin\theta},
\]
where the entire interpolation remainder $H$ is defined by
\[
 e^{\ii\theta t}-\cos\theta-\ii t\sin\theta
       =(1-t^2)\theta^2H(\theta,t).
\]
Set $A_\theta(t,x)=1+x(s_\theta'(t)-1)$.
Since $dt/dy=1-t^2$, differentiation in $z=x+20\ii y$
gives the Beltrami coefficient
\[
 \mu_\theta(t,x)=
 \frac{D_\theta(t)+\ii A_\theta(t,x)/20}
      {D_\theta(t)-\ii A_\theta(t,x)/20}.
\]
The affine change from $\zeta$ to the physical $w$-coordinate is
conformal and leaves this coefficient unchanged.

The reproducible outward-rounded Arb calculation in
\texttt{certify\_exterior\_flat\_cylinder.py} encloses this formula
on the whole box
$|\operatorname{Re}\theta-0.1000435842|\le10^{-8}$,
$0\le\operatorname{Im}\theta\le0.00166823$,
$-1\le t\le1$, $0\le x\le1$, and proves
$|\mu_\theta|<1/10$.  To avoid interval division by $1-t^2$ at the
ends, it evaluates the Taylor terms of orders $2$ through $11$ and bounds
the remaining modulus by
\[
 \frac{R^{10}}{2\,11!\,(1-R/12)},\qquad R=0.102.
\]
Thus the interpolation has positive Jacobian in its interior.  Its
boundary maps once around the certified Jordan curve; the degree
formula then makes it a homeomorphism onto $H_q$ and a local
diffeomorphism inside.  The two vertical sides match under $E_b$,
so it descends to a quasiconformal homeomorphism of the sewn
punctured cylinder.  The standard inequality
$K\le(1+\|\mu\|_\infty)/(1-\|\mu\|_\infty)$ gives $K<11/9$.
The uniform bound also makes the two isolated ends removable for the
pulled-back complex structure.
\end{proof}

\Needspace{14\baselineskip}
\begin{proposition}[Marked exterior germ limits on a boundary subarc]
\label{prop:exterior-brjuno-germ-limit}
Let $I=[0.1000435842-10^{-8},0.1000435842+10^{-8}]$.
For every $p\in I$ there is a unique
$q_c(p)\in(0,0.00166823)$ with
$|\lambda_+(p+\ii q_c(p))|=1$.  On the exterior side
$0\le q<q_c(p)$, the real-normalised Kneser germ transported from
$q=0$ has a common complex height disc $V$ about $z=0$,
independent of $(p,q)$ in this parameter region.  As
$q\uparrow q_c(p)$, these germs converge locally uniformly on $V$
to nonconstant holomorphic boundary germs, uniformly in $p\in I$.
Every limit takes the value $1$ and has nonzero derivative at $0$.
In particular this holds at the prescribed Brjuno multiplier of
Proposition~\ref{prop:exterior-brjuno-path}; the existence of these
local boundary germs does not require a Brjuno hypothesis.  No
common domain containing $[0,1]$ or matching with an inner branch
is asserted.
\end{proposition}

\begin{proof}
The interval certificate of
Proposition~\ref{prop:exterior-brjuno-path} verifies
$|\lambda_+(p)|>1$ throughout $I$ and
$\partial_q|\lambda_+(p+\ii q)|^2<-1$ on the whole enclosing box.
It also subdivides $I$ into $100$ rational intervals and checks
$|\lambda_+(p+0.00166823\,\ii)|<1$ on each one.
The intermediate-value theorem and strict derivative give the
unique $q_c(p)$, which depends real-analytically on $p$.
The same certificate puts the marked iterate in $H_{p+\ii q}$
with a uniform strict margin for the entire closed region
$K=\{(p,q):p\in I,\ 0\le q\le q_c(p)\}$.

Use the standard punctured sphere coordinate
$\xi=\exp(2\pi\ii z)$ on the flat cylinder of
Proposition~\ref{prop:exterior-flat-cylinder}.  Pulling back the
physical complex structure of the sewn cylinder gives a Beltrami
coefficient $\mu_{p,q}$ on this fixed sphere, with
$\|\mu_{p,q}\|_\infty<1/10$.  The explicit formula in that
proposition, including its factored endpoint limits, shows that
$(p,q)\mapsto\mu_{p,q}$ is continuous in the $L^\infty$ norm on
$K$.  The normalised solutions of the Beltrami equation therefore
vary continuously on the sphere by~\cite[Thm.~8]{AhlforsBers1960}.
Their source-coordinate marked points also vary continuously,
because $w_{28}(p,q)$ stays uniformly inside the Jordan cell and the
flat-cylinder parameterisations are homeomorphisms varying
continuously.  Three-point normalisation at the two punctures and
$w_{28}$ consequently gives conformal coordinates $Q_{p,q}$ that
vary locally uniformly in the physical chart near the marked point.
The boundary coordinates are defined by this same uniformisation,
even when the upper multiplier is neutral.  For $q<q_c(p)$ they
coincide with the repelling-end quotient coordinates used to
transport the classical exterior germ, by uniqueness of the
three-point normalisation.

The certified interior margin supplies a fixed physical disc about
$w_{28}(p,q)$ in the moving cells.  On this compact family of discs,
$Q_{p,q}$ and all its $w$-derivatives depend continuously on
$(p,q)$ by Cauchy's formula.  Their first derivatives at the marked
points never vanish, since each $Q_{p,q}$ is conformal.
Compactness of $K$ and the inverse function theorem give a single
small disc $V$ on which the inverses of the normalised Abel charts
\[
 A_{p,q}(w)=\frac{1}{2\pi\ii}
       \log Q_{p,q}(E_{b(p+\ii q)}^{\circ28}(w)),
 \qquad A_{p,q}(1)=0,
\]
exist and depend continuously on $(p,q)$.  For $q<q_c(p)$ their
inverses are precisely the transported Kneser germs.  Uniform
continuity on the compact set $K$ proves the asserted boundary
convergence and nondegeneracy, uniformly in $p$.
\end{proof}

\Needspace{17\baselineskip}
\begin{proposition}[An exterior analytic base corridor for the germ]
\label{prop:exterior-analytic-germ-corridor}
There are a convex complex $\theta$-neighbourhood $U$ of the closed
rectangle
\[
 \widehat K=\{\theta:|\operatorname{Re}\theta-0.1000435842|
       \le10^{-8},\ |\operatorname{Im}\theta|\le10^{-2}\}
\]
and a disc $V$ about $z=0$ such that the real-normalised exterior Kneser germ
extends jointly holomorphically to $U\times V$.  The map
$\theta\mapsto b(\theta)$ is injective on $U$, so its image is a
genuine complex base corridor containing the real exterior seed
segment and crossing a non-cusp Shell--Thron boundary subarc.
The assertion is about the germ on $V$; it does not give a common
regular neighbourhood of $[0,1]$ or identify the continuation with
the inner sewn branch.
\end{proposition}

\begin{proof}
The outward-rounded certificate
\texttt{certify\_exterior\_wide\_corridor.py} checks on $\widehat K$
that the image arc remains an increasing graph below the chord,
$E_b^{\circ28}(1)$ remains inside the Jordan cell, and the singular
orbit $E_b^{\circ k}(0)$ lies above its chord for $0\le k\le28$.
It also gives $|\mu_\theta|<1/10$ on the entire flat cylinder and
$\operatorname{Re}b'(\theta)>0.05$.  All inequalities are strict,
so they persist on a slightly larger convex open rectangle $U$.
In the fixed flat-cylinder coordinate
$\xi=e^{2\pi\ii z}$ of
Proposition~\ref{prop:exterior-flat-cylinder}, the pulled-back
Beltrami coefficient $\mu_\theta$ depends holomorphically on
$\theta$ as an $L^\infty$-valued function: its explicit formula is
holomorphic in $\theta$ pointwise, uniformly bounded away from
modulus one, and has uniform Taylor bounds on a slightly larger
parameter rectangle.  The Ahlfors--Bers holomorphic-parameter
theorem~\cite[Thm.~11]{AhlforsBers1960} supplies a solution
$h_\theta(\xi)$ of the Beltrami equation, normalised at the fixed
source points $0,1,\infty$, that is holomorphic in $\theta$ for
each fixed $\xi$.

One must account for the fact that the inverse of the explicit
quasiconformal map $\Psi_\theta$ from the flat cylinder to the
physical Jordan cell is not holomorphic in the parameter at fixed
physical $w$.  Put $\xi_\theta(w)=\Psi_\theta^{-1}(w)$ near the
marked point.  The maps $h_\theta$ and $\Psi_\theta$ have the same
Beltrami coefficient in $\xi$, while $\Psi_\theta(\xi)$ is
holomorphic in $\theta$ for fixed $\xi$.  Differentiating
$\Psi_\theta(\xi_\theta(w))=w$ at fixed $w$ gives
\[
 \partial_{\bar\theta}\xi_\theta(w)
 +\mu_\theta(\xi_\theta(w))
       \partial_{\bar\theta}\overline{\xi_\theta(w)}=0.
\]
The same identity and $\partial_{\bar\xi}h_\theta
=\mu_\theta\partial_\xi h_\theta$ imply
$\partial_{\bar\theta}
  [h_\theta(\xi_\theta(w))]=0$.
Differentiating instead at fixed $\theta$ with respect to
$\bar w$ proves that this composition is holomorphic in $w$.
These calculations take place in a smooth interior chart, where
all derivatives are classical; hence the composition is jointly
holomorphic.  They apply also at the holomorphic marked section
$w_{28}(\theta)=E_{b(\theta)}^{\circ28}(1)$.

Divide this composition by its nonzero value at $w_{28}(\theta)$.
The resulting coordinate $Q_\theta(w)$ takes the two punctures
to $0,\infty$ and the marked point to $1$, and is jointly
holomorphic in $(\theta,w)$ near the marked section.  Composing
with $E_{b(\theta)}^{\circ28}$, taking the logarithm with value
zero at $w=1$, and applying the parameter-dependent inverse
function theorem yields a jointly holomorphic normalised
superfunction on $U\times V$, after shrinking $U,V$ if necessary.
At $q=0$ it is the classical Kneser germ, so uniqueness of analytic
continuation identifies it throughout the connected corridor.

Finally the wide certificate checks
$\operatorname{Im}(db(p+\ii q)/dq)>0.05$ on $\widehat K$.
Since $db/dq=\ii b'(\theta)$, this is
$\operatorname{Re}b'(\theta)>0.05$ there.  Shrink the convex $U$ so
that $\operatorname{Re}b'>0$ on it.  For distinct
$\theta_1,\theta_2\in U$, the integral of $b'$ along their joining
segment has positive real part, so
$(b(\theta_2)-b(\theta_1))/(\theta_2-\theta_1)\ne0$.
Thus $b$ is injective on $U$ as claimed.
\end{proof}

\Needspace{17\baselineskip}
\begin{proposition}[The real seam anchor and an elliptic crossing test]
\label{prop:exterior-real-seam-anchor}
For every real $p\in I$ of
Proposition~\ref{prop:exterior-brjuno-germ-limit}, the normalised seam
function
\[
 g_{p,q}(t)=\log|Q_{p,q}(\gamma_{p+\ii q}(t))|,
 \qquad -1<t<1,
\]
has exactly one zero at $q=0$, namely $t=0$, and that zero is simple
with $\partial_tg_{p,0}(0)<0$.  For fixed $p$, a change in the number
of seam crossings on $0\le q\le q_c(p)$ requires a pair $(q,t)$
with
\[
 g_{p,q}(t)=\partial_tg_{p,q}(t)=0.
\]
On the symmetric wide parameter box, the seam functions also obey
$g_{p,-q}(-t)=-g_{p,q}(t)$.
In particular, for $p=p_c$, exclusion of such pairs on this compact
parameter path proves the hypothesis of
Proposition~\ref{prop:exterior-single-crossing} at its prescribed
endpoint.

There is an equivalent scalar elliptic test on the fixed flat
cylinder.  Write $Y=20\operatorname{artanh}t$, let $\Psi_{p,q}$
denote the parameterisation of
Proposition~\ref{prop:exterior-flat-cylinder} in the physical
$w$-coordinate, and set
\[
 v_{p,q}(x,Y)=-\frac{1}{2\pi}
  \log|Q_{p,q}(\Psi_{p,q}(x+\ii Y))|.
\]
With $\mu=\mu_{p+\ii q}$ from that proposition, $v$ is periodic in
$x$ and satisfies, weakly across the seam,
\[
 \nabla\!\cdot(\mathsf A_\mu\nabla v)=0,\qquad
 \mathsf A_\mu=
 \frac1{1-|\mu|^2}
 \begin{pmatrix}
 |1-\mu|^2&-2\operatorname{Im}\mu\\
 -2\operatorname{Im}\mu&|1+\mu|^2
 \end{pmatrix}.
\]
The two equations in the displayed crossing condition are exactly
$v_{p,q}(0,Y)=\partial_Yv_{p,q}(0,Y)=0$.
\end{proposition}

\begin{proof}
At a real parameter, complex conjugation preserves the initial
region and commutes with $E_b$.  It descends to an antiholomorphic
involution of the sewn sphere that exchanges its two punctures and
fixes the real marked point.  In the three-point coordinate it is
$Q\mapsto1/\overline Q$; its fixed locus is $|Q|=1$.
Conjugation sends the seam point labelled $t$ to the one labelled
$-t$.  The quotient seam is embedded, so a point on it is fixed
only when $t=0$.  Near this point the real fixed locus has the
horizontal physical tangent, while the chord seam has a nonzero
vertical tangent.  Conformality of the quotient chart makes their
intersection transverse.  Since $g$ tends from $+\infty$ to
$-\infty$ as $t$ increases, its derivative at the unique zero is
negative.
For complex-conjugate parameters, the same involution interchanges
the two punctures and preserves the marked normalisation, giving the
displayed symmetry for $g$.

The uniform quasiconformal bound and compactness of the parameter
region give uniform puncture tails for the normalised $Q_{p,q}$.
Thus all zeros of $g_{p,q}$ lie in one compact subinterval of
$(-1,1)$.  On that subinterval, the physical seam charts and the
normalised uniformisation vary continuously in $C^1$ with $q$.
If no double equation in the statement occurs, the implicit
function theorem continues every simple zero locally in $q$;
compactness and the uniform tails keep their number constant.

For the elliptic formula, choose a local lift
$h=(2\pi\ii)^{-1}\log(Q\circ\Psi)$ on the flat cylinder.
Its imaginary part is $v$, and it satisfies
$h_{\bar z}=\mu h_z$.  Writing $h=u+\ii v$ and separating real
and imaginary parts gives
$\mathsf A_\mu\nabla v=(-u_Y,u_x)$.
Taking the divergence proves the weak equation, including the
periodic seam matching.  Finally,
$g_{p,q}(t)=-2\pi v_{p,q}(0,20\operatorname{artanh}t)$ and
$dY/dt=20/(1-t^2)$, which gives the stated equivalence.
\end{proof}

\Needspace{16\baselineskip}
\begin{lemma}[Uniform positive slopes at the two punctures]
\label{lem:exterior-seam-end-slope}
On the wide rectangle $\widehat K$ of
Proposition~\ref{prop:exterior-analytic-germ-corridor}, let
$v_\theta$ be the elliptic height function of
Proposition~\ref{prop:exterior-real-seam-anchor}.  Put
$\lambda_\pm=\theta\cot\theta\pm\ii\theta$ and take the
principal logarithms $\ell_\pm=\operatorname{Log}\lambda_\pm$.
Uniformly for $\theta\in\widehat K$,
\[
 \begin{aligned}
 \partial_Yv_\theta(0,Y)&\longrightarrow
       \operatorname{Im}\!\left(-\frac1{10\ell_+}\right)>0.98
       &&(Y\to+\infty),\\
 \partial_Yv_\theta(0,Y)&\longrightarrow
       \operatorname{Im}\!\left(\frac1{10\ell_-}\right)>0.98
       &&(Y\to-\infty).
 \end{aligned}
\]
Thus some finite $Y_0$ works uniformly on $\widehat K$:
$\partial_Yv_\theta(0,Y)>0$ for $|Y|>Y_0$.
Every tangency in Proposition~\ref{prop:exterior-real-seam-anchor}
therefore lies in $|Y|\le Y_0$.
\end{lemma}

\begin{proof}
The wide interval certificate also verifies
$\operatorname{Re}\lambda_\pm>0.9$,
$\operatorname{Im}\lambda_+>0.09$, and
$\operatorname{Im}\lambda_-<-0.09$ throughout $\widehat K$,
as well as the two $0.98$ slope bounds above.
In particular the stated logarithms vary continuously, have
imaginary parts of the indicated signs, and never vanish.
The factored formula for $\mu_\theta(x,Y)$ is analytic in
$t=\tanh(Y/20)$ at $t=\pm1$.  It therefore tends, uniformly in
$\theta$ and piecewise smoothly in $x$, to periodic coefficients
$\mu_{\theta,\pm}(x)$ as $Y\to\pm\infty$.

Choose the uniformising lift $h_\theta$ on the universal cover of
the cylinder so that $h_\theta(z+1)=h_\theta(z)+1$ and
$h_\theta(0)=0$.  It is a quasiconformal homeomorphism of the
plane with coefficient $\mu_\theta$ and dilatation below $11/9$;
its imaginary part differs from $v_\theta$ by a constant.
For any $\theta_j\to\theta_\infty$ in $\widehat K$ and
$Y_j\to+\infty$, the maps
\[
 h_{\theta_j}(z+\ii Y_j)-h_{\theta_j}(\ii Y_j)
\]
fix $0,1,\infty$ and have a locally uniformly convergent
subsequence by quasiconformal compactness.  The limit is the unique
normalised solution for $\mu_{\theta_\infty,+}$.

This limit is explicit.  Since $1-\tanh(Y/20)$ is asymptotic to
$2e^{-Y/10}$, the physical interpolation near $L_+$, after a
nonzero conformal rescaling, tends to
$e^{-Y/10}[1+x(\lambda_+-1)]$.  The normalised straight-sector
solution for its coefficient is
\[
 h_{\theta,+}(x+\ii Y)=
 \frac{\log[1+x(\lambda_+-1)]-Y/10}{\ell_+},
 \qquad 0\le x\le1,
\]
extended by $h_{\theta,+}(z+1)=h_{\theta,+}(z)+1$.
The logarithm along the segment from $1$ to $\lambda_+$ is chosen
to run from $0$ to $\ell_+$.  The segment remains in the right
half-plane, and its argument increases strictly; this formula is
the normalised quasiconformal homeomorphism with coefficient
$\mu_{\theta,+}$.  Uniqueness identifies it with every translated
limit.  At the lower end, $1+\tanh(Y/20)\sim2e^{Y/10}$ and the
same argument gives
\[
 h_{\theta,-}(x+\ii Y)=
 \frac{\log[1+x(\lambda_--1)]+Y/10}{\ell_-}.
\]

The preceding convergence also holds for the tangential
$Y$-derivatives at the seam.  Indeed, after dividing the physical
coordinate $w-L_\pm$ by the corresponding factor
$e^{\mp Y_j/10}$, the analytic gluing maps converge on fixed
seam disks to multiplication by $\lambda_\pm$, uniformly in
$\theta_j$.  In these rescaled physical charts the lifted maps
are holomorphic.  Local uniform convergence followed by Cauchy's
formula gives convergence of their physical derivatives, and the
explicit rescaled interpolation gives convergence of their
$Y$-derivatives.  Hence the limits of $\partial_Yv$ are the
imaginary parts of $\partial_Yh_{\theta,\pm}$ displayed above.
Compactness of $\widehat K$ makes the convergence uniform.
The two positive continuous limiting slopes have a positive
minimum there, proving the last assertion.
\end{proof}

\Needspace{17\baselineskip}
\begin{proposition}[A certified frozen Abel coordinate]
\label{prop:exterior-frozen-abel}
Use the chord data $D_\theta(t)$ and $A_\theta(t,x)$ of
Proposition~\ref{prop:exterior-flat-cylinder}, with
$t=\tanh(Y/20)$.  Set
\[
 L_\theta(Y)=
   \operatorname{Log}\!\left(\frac{\theta}{\sin\theta}\right)
       +\ii\theta t,\quad
 B_\theta=e^{L_\theta}=s'_\theta(t),\quad
 P_\theta(x,Y)=\frac{A_\theta(t,x)}{20D_\theta(t)}.
\]
On the full marked path, including its boundary endpoint,
$|P_\theta-\ii|<0.09$.  On the wider rectangle $\widehat K$,
the explicit frozen vertical slope
\[
 c_\theta(Y)=
 \frac{e^{L_\theta(Y)}-1}
      {20D_\theta(t)L_\theta(Y)}
\]
has $\operatorname{Im}c_\theta(Y)>0.98$ for every real $Y$.
The quotient at $L_\theta=0$ is interpreted by analytic
continuation.  Define
\[
 F_\theta(x,Y)=
 \frac{\operatorname{Log}
    [1+x(e^{L_\theta(Y)}-1)]}{L_\theta(Y)},\qquad
 h^0_\theta(x+\ii Y)=F_\theta(x,Y)+\int_0^Yc_\theta(s)\,ds.
\]
Extending $h^0_\theta$ by $h^0_\theta(z+1)=h^0_\theta(z)+1$
gives an explicit approximate Abel lift.  Its exact residual in
the pulled-back Beltrami equation is
\[
 \partial_Yh^0_\theta-P_\theta\partial_xh^0_\theta
       =\partial_YF_\theta,\qquad
 |\partial_YF_\theta(x,Y)|
       <10^{-3}\operatorname{sech}^2(Y/20).
\]
Its residual has $L^2$ norm below $0.005$ on the marked path.
At the seam the approximate height derivative is greater than
$0.98$.
The residual and norm bounds do not yet give pointwise control of
the exact uniformising derivative.
\end{proposition}

\begin{proof}
The Beltrami formula of
Proposition~\ref{prop:exterior-flat-cylinder} rearranges exactly to
$h_Y=P_\theta h_x$.  Since
$A_\theta=1+x(B_\theta-1)$, differentiation gives
\[
 \partial_xF_\theta=
 \frac{B_\theta-1}
      {L_\theta A_\theta},\qquad
 P_\theta\partial_xF_\theta=c_\theta.
\]
The displayed residual follows.  The factor $F_\theta$ equals
$0$ at $x=0$ and $1$ at $x=1$, so the periodic extension is
continuous and its seam height derivative is $c_\theta$.

For a bound without cancellation at $L=0$, write
\[
 F(L,x)=\int_0^1
    \frac{x e^{sL}}{1+x(e^{sL}-1)}\,ds.
\]
Its $L$-derivative has integrand
\[
 \frac{s x(1-x)e^{sL}}{[1+x(e^{sL}-1)]^2}.
\]
The accompanying Arb script
\texttt{certify\_exterior\_frozen\_abel.py} verifies
$|L|<0.12$ on $\widehat K$.  Consequently,
\[
 |F_L|\le
 \frac{e^{0.12}}{8(2-e^{0.12})^2},\qquad
 |L_Y|\le\frac{0.102}{20}\operatorname{sech}^2(Y/20).
\]
Their product is less than
$0.000945\operatorname{sech}^2(Y/20)$.
Squaring and integrating uses
\[
 \int_{\mathbb R}\operatorname{sech}^4(Y/20)\,dY=80/3
\]
and gives the stated $L^2$ bound.
The same outward-rounded script encloses
$P_\theta$ on the marked path and $c_\theta$ on the wide box;
it evaluates $(e^L-1)/L$ by a Taylor polynomial with a rigorous
tail, so intervals containing $L=0$ cause no division problem.
\end{proof}

\begin{proposition}[An $L^4$ correction to the frozen coordinate]
\label{prop:exterior-frozen-correction}
On the full marked path let $h_\theta$ be the exact uniformising
lift with $h_\theta(0)=0$ and
$h_\theta(z+1)=h_\theta(z)+1$.  The coordinate $h^0_\theta$
of Proposition~\ref{prop:exterior-frozen-abel} satisfies, on the
infinite flat cylinder,
\[
 \|\partial_x(h_\theta-h^0_\theta)\|_{L^4}<0.032,
 \qquad
 \|\partial_Y(h_\theta-h^0_\theta)
       -\ii\partial_x(h_\theta-h^0_\theta)\|_{L^4}<0.005.
\]
In particular
$\|\partial_Y(h_\theta-h^0_\theta)\|_{L^4}<0.037$.
These are global integral bounds; they do not assert a pointwise
sign for the exact seam derivative.
\end{proposition}

\begin{proof}
Write $\mathcal L=\partial_Y-\ii\partial_x$ and
$f=\partial_YF_\theta$.  The preceding proposition gives
\[
 \|P_\theta-\ii\|_\infty<0.09,
 \qquad
 \|f\|_4<0.000945
       \left(20\int_{\mathbb R}\!\operatorname{sech}^8s\,ds\right)^{1/4}
       =0.000945(640/35)^{1/4}<0.002.
\]
An exact correction $e$ must solve
$\mathcal Le=(P_\theta-\ii)e_x-f$.
On the cylinder let $B_{\rm per}$ be the periodised Beurling
transform.  The planar bound of~\cite[Thm.~1.2]{BanuelosJanakiraman}
at exponent $4$ gives $\|B\|_4\le4.725<12$.
The same bound $\|B_{\rm per}\|_4<12$ follows by cutting a smooth
periodic input of compact vertical support to $M$ horizontal periods,
applying the planar estimate, and averaging the output over the
interior $(1-\varepsilon)M$ periods.  On each fixed-height output
strip the omitted kernel tails tend uniformly to zero, since the
kernel is $O(|x|^{-2})$ away from the input; first let $M\to\infty$,
then $\varepsilon\downarrow0$, then release the output-height cutoff.
Density gives the bound on cylinder $L^4$.

Since $\mathcal L=-2\ii\partial_{\bar z}$, the cylinder
Cauchy operator has
\[
 T:=\partial_x\mathcal L^{-1}
   =-\frac{I+B_{\rm per}}{2\ii},
 \qquad \|T\|_4<6.5.
\]
Solve for $u$ as the unique fixed point in cylinder $L^4$ of
$u=T[(P_\theta-\ii)u-f]$, because $6.5(0.09)=0.585<1$.
It obeys
\[
 \|u\|_4<\frac{6.5(0.002)}{1-0.585}<0.032,
 \qquad
 \|g\|_4
 \le0.09\|u\|_4+\|f\|_4<0.005.
\]
For completeness, construct $e$ from this fixed point by applying
the periodic Cauchy potential to
$g=(P_\theta-\ii)u-f$.  For compactly supported smooth inputs,
its gradient is the usual Cauchy and Beurling pair.  Approximation
in $L^4$, the bounds above, and local Morrey compactness give a
periodic $e\in W^{1,4}_{\rm loc}$ with $\nabla e\in L^4$ and
$e(0)=0$.  The identity $e_x=u$ passes to the limit; uniqueness
modulo constants follows because a periodic entire holomorphic
function with derivative in $L^4$ is constant.

Set $H=h^0_\theta+e$.  The exact equation
$H_Y=P_\theta H_x$ says that $H$ has the same Beltrami
coefficient as $h_\theta$, of modulus below $1/10$.
The positive bound $\operatorname{Im}c_\theta>0.98$ and boundedness
of $F_\theta$ show that $\operatorname{Im}h^0_\theta$ tends to
$+\infty$ and $-\infty$ at the corresponding cylinder ends.
The $L^4$ gradient bound makes $e=o(|Y|)$ uniformly in
$0\le x\le1$: its horizontal mean grows at most
$O(|Y|^{3/4})$ by H\"older's inequality, while the oscillation
on each unit-height strip is uniformly bounded by Morrey's
inequality.  Hence $H$ is proper.  The homotopy
$(1-s)z+sH(z)$ is proper for $0\le s\le1$ and has degree one.
The weak quasiregular factorisation (or the quasiconformal chain
rule followed by Weyl's lemma) makes $H\circ h_\theta^{-1}$
an entire holomorphic map.  It is proper of degree one, so affine;
the period relation forces its slope to be $1$, and the value at
$0$ removes the constant.  Therefore $H=h_\theta$, proving the
claimed bounds for the exact lift.
\end{proof}

\begin{proposition}[Certified pointwise seam monotonicity]
\label{prop:exterior-seam-monotonicity}
On the closed parameter rectangle
$|\operatorname{Re}\theta-0.1000435842|\le10^{-8}$,
$|\operatorname{Im}\theta|\le0.01$, including both sides of its
non-cusp boundary subarc, the exact uniformising lift satisfies
\[
 \sup_{z\in\mathbb C}|\partial_Y(h_\theta-h^0_\theta)(z)|<0.041,
 \qquad
 \partial_Y\operatorname{Im}h_\theta(\ii Y)>0.93
 \quad(Y\in\mathbb R).
\]
Consequently the image of the chord seam under the
three-point-normalised coordinate $Q_\theta$ crosses $|Q|=1$
exactly once and transversely, for every parameter in this
rectangle.
\end{proposition}

\begin{proof}
Let $e=h_\theta-h^0_\theta$ and $f=\partial_YF_\theta$.
The same outward-rounded interval computation as in
Proposition~\ref{prop:exterior-frozen-abel}, now differentiating
the factored interpolation remainder in $t$, gives throughout the
wider rectangle
$\|P_\theta-\ii\|_\infty<0.18$ and
$\|\partial_YP_\theta\|_\infty<0.004$.
Here the derivative of its order-$n$ polynomial has modulus at most
$n^2/4$ on $[-1,1]$; starting at $n=12$, its Taylor tail is bounded by
$36(0.102)^{10}/[12!(1-0.102/10)]$.
For the residual, put $E=e^{0.12}$ and $a=2-E$.  Differentiating
the integral formula for $F_L$ in the preceding proof gives
\[
 |F_L|\le\frac{E}{8a^2},\qquad
 |F_{LL}|\le\frac{E}{12a^2}+\frac{E^2}{6a^3}.
\]
Since $|L_Y|\le(0.102/20)\operatorname{sech}^2(Y/20)$ and
$|L_{YY}|\le(0.102/200)\operatorname{sech}^2(Y/20)$,
the interval evaluation yields
\[
 |f_Y|<0.00012\operatorname{sech}^2(Y/20),
 \qquad \|f_Y\|_4<0.00025.
\]

The planar Beurling bound $\|B\|_4\le4.725<6$ quoted in
Proposition~\ref{prop:exterior-frozen-correction} transfers to the
periodic cylinder by the same cutoff-and-averaging argument used
there.  Hence the cylinder operator
$T=\partial_x\mathcal L^{-1}$ has norm below $3.5$.
The fixed-point construction of that proposition now applies on
the wider rectangle because $3.5(0.18)=0.63<1$.  It gives
\[
 \|e_x\|_4<\frac{3.5(0.002)}{1-0.63}<0.019,
 \quad \|\mathcal Le\|_4<0.18(0.019)+0.002<0.0055,
 \quad \|e_Y\|_4<0.025.
\]
The same properness and degree-one argument identifies this
correction with the exact uniformising lift.

Take a vertical difference quotient $e^\delta(x,Y)
=[e(x,Y+\delta)-e(x,Y)]/\delta$.  Its $x$ derivative belongs
to cylinder $L^4$ for each $\delta\ne0$, and it satisfies
\[
 \mathcal Le^\delta=(P_\theta-\ii)(e^\delta)_x+g^\delta,
 \quad
 g^\delta=P_\theta^\delta e_x(x,Y+\delta)-f^\delta.
\]
The bounds above give
$\|g^\delta\|_4<0.004(0.019)+0.00025<0.00033$.
The same cylinder operator $T$ as before therefore gives
\[
 \|(e^\delta)_x\|_4
 <\frac{3.5(0.00033)}{1-0.63}<0.0032,
 \qquad
 \|(e^\delta)_Y\|_4
 <1.18(0.0032)+0.00033<0.0042.
\]
Passing to weak limits of the difference quotients proves that
$v=e_Y\in W^{1,4}(\mathbb C/\mathbb Z)$ and
$\|v_x\|_4<0.0032$, $\|v_Y\|_4<0.0042$;
the preceding bound gives $\|v\|_4<0.025$.

For clarity, a disk of radius $r<1/2$ on the flat cylinder obeys
the following pointwise Sobolev estimate, obtained by averaging
along radial segments and applying H\"older's inequality:
\[
 |v(z_0)|\le(\pi r^2)^{-1/4}\|v\|_{L^4(B_r(z_0))}
 +(3/2)^{3/4}(2\pi)^{-1/4}r^{1/2}
        \|\nabla v\|_{L^4(B_r(z_0))}.
\]
At $r=1/4$ its two coefficients are below $1.51$ and $0.43$.
Hence $|e_Y(z_0)|<1.51(0.025)+0.43(0.0032+0.0042)<0.041$
for every $z_0$.  The $W^{1,4}$ representative is continuous;
on the seam it equals the classical tangential derivative supplied
by the analytic physical chart.  Since
$h^0_Y(\ii Y)=c_\theta(Y)$ and
$\operatorname{Im}c_\theta(Y)>0.98$, the stated strict derivative
bound follows.

The seam height $\operatorname{Im}h_\theta(\ii Y)$ tends to the
opposite infinite limits at the two ends by the proof of
Proposition~\ref{prop:exterior-frozen-correction}.
It is therefore strictly increasing through every real level.
The three-point-normalised $Q_\theta$ differs from
$\exp(2\pi\ii h_\theta)$ by a nonzero constant factor, so
$|Q_\theta|=1$ is one such height level.  It is crossed exactly
once, with nonzero derivative.
\end{proof}

\Needspace{12\baselineskip}
\begin{proposition}[An attracting overlap of the exterior chord]
\label{prop:interior-chord-overlap}
At the exact interior parameter
$\theta_0=0.1000435842+0.01\ii$, the upper fixed point
$L_+(\theta_0)$ is attracting and the lower fixed point is
repelling.  The entire physical chord subarc
$\gamma_{\theta_0}([-13/100,1])$ lies in the immediate attracting
basin of $L_+(\theta_0)$.  This supplies a common physical chart
region for the exterior quotient and the upper attracting dynamics;
it does not identify the exterior solution with the inner sewn
solution.  Moreover, the unique point of the chord mapped by the
marked quotient coordinate to $|Q_{\theta_0}|=1$ lies on this
certified attracting subarc, at a parameter $t_*>0$.
Consequently there is a noninteger real height $z_*\in(0,1)$
and product neighbourhoods of $(\theta_0,z_*)$ on which the
continued exterior superfunction takes values in the attracting
basin and admits a jointly holomorphic upper Koenigs coordinate.
\end{proposition}

\begin{proof}
The outward-rounded Arb certificate
\texttt{certify\_interior\_chord\_overlap.py} uses exact rational
coordinates for $\theta_0$ and partitions the indicated real
$t$-interval into $113$ cells of width $1/100$.  Put
$a=\log b(\theta_0)$, $L=L_+(\theta_0)$, and $\delta=1/100$.
It verifies $|aL|e^{|a|\delta}<1$ and $|aL_-|>1$.
For a cell centred at $t_j$, let
$c_0=\gamma_{\theta_0}(t_j)$ and
$\rho_0=|\ii r\sin\theta_0|/200$.
The rigorous recurrence
\[
 c_{n+1}=e^{a c_n},\qquad
 \rho_{n+1}=|c_{n+1}|\bigl(e^{|a|\rho_n}-1\bigr)
\]
encloses the whole cell image under $E_b^{\circ n}$ in
$\overline{B(c_n,\rho_n)}$.  For each of the $113$ cells the
certificate finds an $n\le505$ with
$|c_n-L|+\rho_n<\delta$.
The disk $B(L,\delta)$ is strictly forward invariant by the
first verified inequality, because
$|E_b(L+u)-L|\le|aL|e^{|a|\delta}|u|$.
Thus the whole subarc belongs to the attracting basin.
It is connected and contains $L$, so it lies in the immediate
component.

For the final assertion, write the marked iterate in the flat
interpolation coordinate as
$\zeta_{28}=t_m+x_m(s_{\theta_0}(t_m)-t_m)$,
$0<x_m<1$, and put $Y_m=20\operatorname{artanh}t_m$.
The same Arb certificate checks
$\operatorname{Re}\zeta_{28}>0.15$ and
$\operatorname{Re}(s_{\theta_0}(t)-t)<0.01$ for every
$-1\le t\le1$; hence $t_m>0.14$ and $Y_m>2.8$.
The formula for $F_x$ in
Proposition~\ref{prop:exterior-frozen-abel}, together with
$|L_{\theta_0}(Y)|<0.12$, gives
\[
 |F_x-1|\le
 \frac{(e^{0.12}-1-0.12)/0.12+(e^{0.12}-1)}{2-e^{0.12}}
 <0.23.
\]
The exact equation $e_Y=P_\theta e_x-f$, the wide-box bounds
$|P_\theta-\ii|<0.18$, $|f|<0.001$, and
$|e_Y|<0.041$ give $|e_x|<0.052$.  Thus
$|h_{\theta_0,x}-1|<0.3$ on the flat cylinder.
Since $h_{\theta_0}(0)=0$ and its seam-height slope exceeds
$0.93$, we have
\[
 \operatorname{Im}h_{\theta_0}(x_m+\ii Y_m)
 >0.93Y_m-0.3>0.
\]
Three-point normalisation at the mark divides
$\exp(2\pi\ii h_{\theta_0})$ by its value there.
Its unit-circle level on the seam therefore has positive height.
Strict seam monotonicity makes its unique crossing satisfy
$Y_*>0$, hence $t_*=\tanh(Y_*/20)>0$.
The seam point represents a noninteger height $z_*\in(0,1)$,
because the marked point is in the interior of the Jordan cell.
The developed value at $z_*$ is a preimage under
$E_b^{\circ28}$ of either $\gamma_{\theta_0}(t_*)$ or its
forward image.  Both lie in the immediate attracting basin, so
that value lies in the full attracting basin.  Finite entry into
the invariant disk persists on product neighbourhoods of
$(\theta_0,z_*)$; Proposition~\ref{prop:upper-orbit} then gives
the stated joint Koenigs coordinate.
\end{proof}

\Needspace{12\baselineskip}
\begin{proposition}[A certified interior sewing circle]
\label{prop:interior-sewing-circle}
At $\theta_0=0.1000435842+0.01\ii$, the entire lower regular
curve $S_{\theta_0}([0,1])$ lies in the full attracting basin of
$L_+(\theta_0)$ and avoids every backward preimage of $L_+$.
Let $\sigma_+$ be the upper Koenigs coordinate, normalised by
$\sigma_+'(L_+)=1$, and let
$\ell_+=\operatorname{Log}\lambda_+$ be the branch continued from the inner real
axis.  The upper time
\[
 T(x)=\frac{1}{\ell_+}
 \log\frac{\sigma_+(S_{\theta_0}(x))}{\sigma_+(1)}
\]
has a branch along $[0,1]$ for which
\[
 0.94<\operatorname{Re}T'(x)<1.06,
 \qquad |\operatorname{Im}T'(x)|<0.056,
 \qquad T(1)-T(0)=1.
\]
Consequently $T([0,1])/\mathbb Z$ is an embedded analytic Jordan
circle in the upper cylinder.  With
$\varpi=2\pi\ii/\ell_+$, its shifted image
$(T+\varpi)([0,1])$ lies in $\operatorname{Im}w<-2.5$.
It therefore defines an analytic two-chart sewing near this
strictly interior parameter, with the normalising upper point $w=0$
above the seam.  The resulting locally normalised sewn solution is
jointly holomorphic in the base and height on its regular domains.
This local family has not been identified with the continuation of
the real inner sewn family or with the exterior family.
\end{proposition}

\begin{proof}
The accompanying outward-rounded Arb certificates use the exact
rational coordinates of $\theta_0$.  Write $a=\log b$, $L_-=L_-(\theta_0)$,
$\mu=aL_-$, $M=|\mu|$, $A=|a|$, and
$f(u)=L_-(e^{au}-1)$.  They verify
$1.01<M<1.02$, $A<0.38$, and, with $\rho=1/3000$,
$\exp(4A\rho/(M-1))<2$.  The normalised lower Poincar\'e
function is the limit of
\[
 P_N(t)=L_-+f^{\circ N}(t\mu^{-N}).
\]
For $|t|\le\rho$, elementary exponential bounds give
$|f^{\circ j}(v)|\le4\rho M^{j-N}$ and
$|(f^{\circ N})'(v)|\le2M^N$ on the line segment joining
$t\mu^{-N}$ to $f(t\mu^{-(N+1)})$.
Indeed the latter endpoint differs from the former by at most
$\frac12 A\rho^2M^{-2N-1}e^{A\rho}$, and both have modulus below
$2\rho M^{-N}$.  Summing the resulting telescoping differences gives
\[
 \sup_{|t|\le\rho}|P(t)-P_N(t)|
 \le \frac{A\rho^2e^{A\rho}}{M-1}M^{-N}
 <39\rho^2(100/101)^N.
\]
At $N=1000$ this last bound is below $2.1\cdot10^{-10}$.

Put $m=700$ and $t(z)=-\exp((z-m)\operatorname{Log}\mu)$.
The certificates check $|t(z)|<\rho$ on $0\le z\le1$ and use
$S(z)=E_b^{\circ m}(P(t(z)))$.  They partition $[0,1]$ into
$32$ rational cells.  On each cell the initial $t$-disk has radius
$|t(z_j)|(e^{|\operatorname{Log}\mu|/64}-1)$.  For a disk of radius $r$ centred
at $u$ under $f$, the new radius is at most
$|L_-e^{au}|(e^{Ar}-1)$; under $E_b$ the analogous radius is
$|E_b(u)|(e^{Ar}-1)$.  Adding the above uniform Poincar\'e tail
after $1000$ steps, the certificates find that every cell enters
$B(L_+,1/100)$ within a further $427$ iterates after the $700$
shifts.  The disk is strictly invariant since
$|aL_+|e^{A/100}<1$.  At entry each enclosure also excludes
$L_+$; since $|a|/100<2\pi$, no later point of the invariant
disk can map exactly to $L_+$.  Thus $\sigma_+(S(x))\ne0$
throughout the interval.

For the slope, the second certificate propagates derivative disks
on a finer $64$-cell partition, using Cauchy's estimate for the
Poincar\'e tail on $|t|<\rho$.  Alongside a value disk
$B(c,r)$ and a derivative disk $B(d,s)$, one exponential step uses
\[
 c_{\rm new}=e^{ac},\qquad d_{\rm new}=ac_{\rm new}d,\qquad
 s_{\rm new}=|ac_{\rm new}|
 \bigl((e^{Ar}-1)|d|+e^{Ar}s\bigr).
\]
The lower Poincar\'e steps have the analogous formula with
$c_{\rm new}=f(c)$.  If
$u_n=E_b^{\circ n}(S(z))-L_+$ and
$d_n=(E_b^{\circ n}\circ S)'(z)$, then
\[
 D_n=\frac{d_n}{u_n\ell_+}\longrightarrow T'(z),\qquad
 \frac{D_{n+1}}{D_n}
 =\frac{au_ne^{au_n}}{e^{au_n}-1}.
\]
On $|u|<1/100$, the last factor differs from $1$ by at most
$A|u|$; subsequent displacements contract by
$q=|aL_+|e^{A/100}<1$.  Thus the multiplicative tail after
any certified iterate is bounded by
$\exp(A|u_n|/(1-q))-1$.  The outward-rounded derivative and
tail disks after $1200$ upper iterates give the stated slope
inequalities.

Since $S(1)=E_b(S(0))$, logarithm continuation gives
$T(1)-T(0)=1+k\varpi$ for some integer $k$.  The certified
bound $\operatorname{Re}\varpi>50$ and the integral of the slope
over $[0,1]$ force $k=0$.  For the seam location, compare the
same-depth ratio
$\mathcal R_n=(E_b^{\circ n}(S(0))-L_+)/
(E_b^{\circ n}(1)-L_+)$.  On the invariant disk, the factor
$h(u)=(e^{au}-1)/(au)$ gives the ratio-tail bound
\[
 \left|\frac{\mathcal R_\infty}{\mathcal R_n}-1\right|
 \le \exp\!\left(
 \frac{A(|u_n^{(0)}|+|u_n^{(1)}|)}{1-q}\right)-1.
\]
The certificate encloses the principal logarithm of this ratio,
including its tail, and locates
$\operatorname{Im}(T(0)+\varpi)<-2.5-0.056$.
The imaginary-slope bound then places the entire shifted circle
below $-2.5$.  Positive real slope and degree one make it an
embedded circle.  A thin analytic collar persists for parameters
near $\theta_0$, and the three-section compactification and
normalisation argument of Proposition~\ref{prop:weld-parameter}
produces the stated local holomorphic sewn family.
\end{proof}

\Needspace{12\baselineskip}
\begin{proposition}[A transverse crossing on the straight interior path]
\label{prop:interior-path-crossing}
Set
$\theta_\dagger=23541/2560000+0.01\ii$ and let
$S=S_{\theta_\dagger}$ be the lower regular superfunction.
There are distinct real heights
\[
\begin{aligned}
x&\in 0.04563349481384156132545513+[-10^{-4},10^{-4}],\\
y&\in 0.2767500343864400645621332+[-10^{-4},10^{-4}]
\end{aligned}
\]
for which $S(x)=S(y)$.  This crossing is transverse.  The physical
subarc $S([0.04,0.28])$ lies in a simply connected disc contained in
the upper attracting chart and avoiding $L_+$.  Hence the same two
heights have equal upper Koenigs times on one logarithmic branch.
Any upper-time sewing curve containing this subarc therefore
self-intersects at $\theta_\dagger$.  Moreover,
$\operatorname{Re}T'(0)<-0.1$.  In particular, the straight path
$\theta=p+0.01\ii$ cannot transport the real inner sewing by
maintaining this lower regular circle as an embedded, positive-slope
seam at every parameter.  This does not rule out continuation by a
different seam or parameter path.
\end{proposition}

\begin{proof}
Both accompanying certificates use the exact rational parameter and
the lower Poincar\'e limit from the preceding proposition.  For the
crossing, \nolinkurl{certify_interior_self_intersection.py} takes
$\rho=1/2000$, Poincar\'e depth $4500$, and $800$ further
exponential iterates.  Its uniform Poincar\'e value tail is bounded by
\[
 \frac{A\rho^2e^{A\rho}}{M-1}M^{-4500},\qquad
 A=|\log b|,\quad M=|E_b'(L_-)|,
\]
and Cauchy's estimate bounds the derivative tail.  It propagates
value and derivative discs over the two stated height intervals.
For $F(x,y)=S(x)-S(y)$, a rational preconditioner of the real
$2\times2$ Jacobian gives a Krawczyk image within
$[-1.3\cdot10^{-6},1.3\cdot10^{-6}]\times[-10^{-6},10^{-6}]$
relative to the box centre.  The two Jacobian defect row sums are
$<0.1$.  Thus the Krawczyk map contracts the box strictly into
itself; its unique zero is transverse, and the two height intervals
are disjoint.

Eight more rational height cells enclose $S([0.04,0.28])$ in
$B(2.69185+0.019515\ii,0.001)$; the maximum certified distance
from its centre is $<1.1\cdot10^{-6}$.  This disc lies in
$B(L_+,0.02)$ and avoids $L_+$.  The latter disc is invariant since
$|E_b'(L_+)|e^{A/50}<1$; the upper Koenigs coordinate has no zero
there except $L_+$, by the local linearisation, the absence of other
preimages in this small disc, and Hurwitz's theorem.  Its logarithm
therefore has a single branch on the smaller disc containing the
whole subarc.  Because $S(x)=S(y)$, that branch gives $T(x)=T(y)$;
the nonzero conformal multiplier from $S'$ to $T'$ preserves
transversality.

Independently, \nolinkurl{certify_interior_path_fold.py} propagates
the value and derivative discs at height $0$ to the invariant upper
disc and bounds the remaining Koenigs product.  It gives
$\operatorname{Re}T'(0)<-0.11272<-0.1$.
\end{proof}

\Needspace{11\baselineskip}
\begin{proposition}[The lower sewing circle in the exterior quotient]
\label{prop:interior-quotient-loop}
At $\theta_0=0.1000435842+0.01\ii$, let $H_{\theta_0}$ be the
exterior initial Jordan cell and set
$w_n(x)=E_{b(\theta_0)}^{\circ n}(S_{\theta_0}(x))$.
There is a unique $x_*\in(0,1)$ such that
$w_{76}(x_*)$ lies on the chord boundary of $H_{\theta_0}$.
For $0\le x<x_*$ one has $w_{77}(x)\in H_{\theta_0}$, and for
$x_*<x\le1$ one has $w_{76}(x)\in H_{\theta_0}$.
At $x_*$ these two representatives are the boundary pair
$w_{77}(x_*)=E_b(w_{76}(x_*))$; at the other join,
$w_{77}(0)=w_{76}(1)$.  The resulting closed curve is an embedded
essential analytic circle in the exterior quotient cylinder.
Its annular physical chart extends jointly holomorphically to
parameters near $\theta_0$.  Thus the local lower sewing circle and the
continued exterior quotient have a common annular physical chart.
This does not identify their global uniformising coordinates.
\end{proposition}

\begin{proof}
Use the chord coordinate
$\zeta(w)=(w/r-\cos\theta_0)/(\ii\sin\theta_0)$.
The chord is $[-1,1]$, its exponential image is
$s([-1,1])$, and direct substitution gives the exact identity
\[
 \zeta(E_b(w))=s(\zeta(w)),\qquad
 s(t)=\frac{e^{\ii\theta_0t}-\cos\theta_0}
              {\ii\sin\theta_0}.
\]
On the central rectangle
$\mathcal R=[0.45,0.55]+\ii[-0.05,0.10]$, the outward-rounded
certificate verifies
$\operatorname{Re}s'>0.98$ and
$\operatorname{Im}s(t)<-0.03$ for real
$t\in[0.45,0.55]$, while
$\operatorname{Re}s(0.45)<0.48<0.52<\operatorname{Re}s(0.55)$.
Thus $s$ is injective on this convex rectangle,
its real arc lies below the chord, and for
$\zeta\in\mathcal R\cap s(\mathcal R)$
\[
 \zeta\in\zeta(H_{\theta_0})
 \quad\Longleftrightarrow\quad
 \operatorname{Im}\zeta<0<\operatorname{Im}s^{-1}(\zeta).
\]
The inverse here is the branch on $s(\mathcal R)$.

The same certificate uses the Poincar\'e approximation and explicit
tail from Proposition~\ref{prop:interior-sewing-circle}, followed by
$75$--$77$ exponential iterates.  It partitions $[0,1]$ into $64$
rational cells and propagates both value and derivative disks.
Throughout the interval it verifies that
$\zeta(w_{75}),\zeta(w_{76}),\zeta(w_{77})$ belong to
$\mathcal R$, all with real parts in $(0.48,0.52)$, and
\[
 \operatorname{Im}\zeta(w_{75})>0.03,
 \qquad \operatorname{Im}\zeta(w_{77})<-0.005,
 \qquad \frac{d}{dx}\operatorname{Im}\zeta(w_{76})<-0.03.
\]
Its endpoint disks put
$\operatorname{Im}\zeta(w_{76}(0))>0$ and
$\operatorname{Im}\zeta(w_{76}(1))<0$, so the middle curve crosses
the chord exactly once and transversely.  The criterion above,
applied to $w_{77}=E_b(w_{76})$ before the crossing and to
$w_{76}=E_b(w_{75})$ afterwards, proves the asserted membership
in $H_{\theta_0}$.

For both $n=76,77$ the certificate also gives
$d\operatorname{Re}\zeta(w_n(x))/dx>10^{-4}$.
Since $w_{77}(0)=w_{76}(1)$, the two interior arcs have disjoint
real-coordinate ranges except at this common endpoint.  The only
other join is the single identified boundary pair at $x_*$.
Hence the quotient curve is simple and analytic, and crossing the
glued boundary once makes it essential.  All inequalities are
strict and all arcs avoid the two fixed endpoints, so compactness
gives a jointly holomorphic annular collar for nearby parameters.
\end{proof}

\begin{corollary}[Two distinct return circles in the exterior quotient]
\label{cor:interior-two-returns}
At the same parameter $\theta_0$, the lower regular circle has a second
embedded essential return circle in the exterior quotient, represented
by $w_{15}(x)$ before a unique chord crossing of $w_{14}(x)$ and by
$w_{14}(x)$ after it.  This circle is disjoint from the return circle
of Proposition~\ref{prop:interior-quotient-loop}.  In particular, a
forward return of the lower circle to the exterior initial cell does
not by itself specify a canonical comparison map between the two
quotient cylinders.
\end{corollary}

\begin{proof}
The independent outward-rounded certificate
\texttt{certify\_interior\_early\_return.py} uses the same Poincar\'e
tail and $64$ rational cells as the preceding proposition.  On
$[0.25,0.33]+\ii[-0.10,0.08]$ it verifies
$\operatorname{Re}s'>0.98$ and
$\operatorname{Im}s(t)<-0.03$ for real $t\in[0.25,0.33]$.
All three curves $\zeta(w_{13})$, $\zeta(w_{14})$, and
$\zeta(w_{15})$ lie in this rectangle, with real parts in
$(0.26,0.31)$, and the certificate gives
\[
 \operatorname{Im}\zeta(w_{13})>0.005,\qquad
 \operatorname{Im}\zeta(w_{15})<-0.03,\qquad
 \frac{d}{dx}\operatorname{Im}\zeta(w_{14})<-0.03.
\]
The endpoint signs of $\operatorname{Im}\zeta(w_{14})$ are positive
and negative.  Thus $w_{14}$ crosses the chord once, and the same
initial-cell criterion as above puts $w_{15}$ before the crossing
and $w_{14}$ afterwards inside $H_{\theta_0}$.  For $n=14,15$,
$d\operatorname{Re}\zeta(w_n)/dx>10^{-4}$; the endpoint identity
$w_{15}(0)=w_{14}(1)$ and the boundary gluing make a simple
essential circle.  Its real-coordinate strip $(0.26,0.31)$ is
disjoint from the strip $(0.48,0.52)$ certified for the earlier
proposition, even after boundary identification.
\end{proof}

\begin{corollary}[Degree-one regular germs at the exterior ends]
\label{cor:interior-end-germs}
Let $Q_{\theta_0}$ be the three-point-normalised coordinate of the
exterior quotient at $\theta_0$.  On the corresponding quotient end
germs, the lower regular superfunction $S$ and the upper regular
superfunction $\Rr$ satisfy
\[
 Q_{\theta_0}(S(z))^{-1}
   =C_-e^{-2\pi\ii z}
       \bigl(1+O(e^{-2\pi\ii z})\bigr)
       \quad(\operatorname{Im}z\to-\infty),
\]
and
\[
 Q_{\theta_0}(\Rr(z))
   =C_+e^{2\pi\ii z}
       \bigl(1+O(e^{2\pi\ii z})\bigr)
       \quad(\operatorname{Im}z\to+\infty),
\]
where $C_\pm\ne0$.  These are degree-one local puncture maps.
Together with a chosen return circle from
Proposition~\ref{prop:interior-quotient-loop} or
Corollary~\ref{cor:interior-two-returns}, this locates the remaining
chart-identification problem in the compact middle bands between
that circle and the two puncture germs.
It does not assert that the two local identifications continue
through those bands or agree with one another.
\end{corollary}

\begin{proof}
Let $\tau_-$ be the Koenigs coordinate at the repelling lower fixed
point, normalised by $\tau_-'(L_-)=1$, and let
$\mu=E_b'(L_-)$.  On a deep lower time half-cylinder,
$\tau_-(S(z))=-\mu^z$ exactly.  The lower exterior puncture
coordinate is, on its logarithmic sector quotient,
\[
 \xi_-=
 \exp\!\left(-\frac{2\pi\ii}{\log\mu}
                      \log(-\tau_-(w))\right).
\]
The choice of logarithmic branch is fixed by the lower regular
curve; multiplication by $\mu$ changes its exponent by an integer
multiple of $2\pi\ii$, so $\xi_-$ descends to the glued end.
Consequently $\xi_-(S(z))=e^{-2\pi\ii z}$.
The exterior quotient was compactified by this simple puncture
coordinate, so its global coordinate has the expansion
$Q_{\theta_0}^{-1}=C_-\xi_-(1+O(\xi_-))$ with $C_-\ne0$.

At the upper end let $\sigma_+$ be the attracting Koenigs
coordinate and $\lambda_+=E_b'(L_+)$.  On a deep upper time
half-cylinder,
$\sigma_+(\Rr(z))=\sigma_+(1)\lambda_+^z$.
The upper exterior puncture coordinate
\[
 \xi_+=\exp\!\left(\frac{2\pi\ii}{\log\lambda_+}
           \log\frac{\sigma_+(w)}{\sigma_+(1)}\right)
\]
therefore obeys $\xi_+(\Rr(z))=e^{2\pi\ii z}$.
The global coordinate has a simple zero there, giving
$Q_{\theta_0}=C_+\xi_+(1+O(\xi_+))$ with $C_+\ne0$.
The local end degrees are one.  The statements concern germs in
the respective quotient end charts; they do not supply analytic
continuation from the middle circle to those germs.
\end{proof}

\Needspace{10\baselineskip}
\begin{proposition}[A single-crossing criterion for a full height segment]
\label{prop:exterior-single-crossing}
Fix the prescribed parameter $p_c$ of
Proposition~\ref{prop:exterior-brjuno-path} and abbreviate
$Q_q=Q_{p_c,q}$ for the three-point-normalised coordinates of
Proposition~\ref{prop:exterior-brjuno-germ-limit}.  Suppose the curve
$t\mapsto Q_{q_c}(\gamma_{q_c}(t))$, $-1<t<1$, meets the unit circle
exactly once, and the meeting is transverse; equivalently,
\[
 \log|Q_{q_c}(\gamma_{q_c}(t))|=0
\]
has exactly one zero $t_*$ and its derivative there is nonzero.
Then there are a complex parameter neighbourhood $U_*$ of
$\theta_c=p_c+\ii q_c$ and a connected open height neighbourhood
$W\supset[0,1]$ such that the exterior germ of
Proposition~\ref{prop:exterior-analytic-germ-corridor} extends
jointly holomorphically to $U_*\times W$.  In particular, it
crosses a non-cusp boundary subarc and takes the values
$F_\theta(0)=1$, $F_\theta(1)=b(\theta)$ there.  The
single-crossing hypothesis is verified by
Proposition~\ref{prop:exterior-seam-monotonicity}.
\end{proposition}

\begin{proof}
The normalised coordinates $Q_\theta$ and their inverses vary
continuously in $C^1$ on compact collars away from the two
punctures by Propositions~\ref{prop:exterior-brjuno-germ-limit}
and~\ref{prop:exterior-analytic-germ-corridor}.  In puncture
coordinates, $Q_\theta$ tends uniformly to $0$ at the upper end
and to $\infty$ at the lower end.  Thus the unique transverse
intersection with $|Q_{\theta_c}|=1$ persists, and remains the
only such intersection, throughout a sufficiently small complex
neighbourhood $U_*$ of $\theta_c$.

The marked point $Q_\theta=1$ and the seam point split
$|Q_\theta|=1$ into two compact subarcs.  On each subarc the
inverse of $Q_\theta$ has a representative in $H_\theta$.
Starting at the interior marked point $w_{28}(\theta)$, pull the
first representative back by the conformal branch $\Phi_\theta$ of
Proposition~\ref{prop:exterior-pullback-cell}.  At the unique seam
crossing the representative passes from the image arc to the chord;
this is the orientation of the positive Abel deck generator.
Continue on the second subarc using
$E_{b(\theta)}\circ\Phi_\theta$.  The seam identity in that proposition makes
the two physical charts agree on their common analytic collar.
They therefore define a holomorphic developing map at every height
$z\in[0,1]$ with $Q_\theta=e^{2\pi\ii z}$, taking the values $1$ and
$E_{b(\theta)}(1)=b(\theta)$ at its endpoints.

The path is compact and avoids both punctures.  Finitely many of the
collars just described cover it.  Their domains, transition maps
and physical developing maps vary jointly holomorphically in the
parameter by Proposition~\ref{prop:exterior-analytic-germ-corridor}
and the finite-log construction of
Proposition~\ref{prop:exterior-pullback-cell}.
The seam identity of that proposition extends to the complex
parameter neighbourhood by the identity theorem, since both sides
are holomorphic in $\theta$ there and agree along the real $q$ path.
After shrinking $U_*$ and these finitely many charts, they cover
one connected complex neighbourhood $W$ of $[0,1]$ for every
$\theta\in U_*$.  The analytic overlap identities provide a jointly
holomorphic family on $U_*\times W$.  Near $z=0$ it is the
previously transported normalised Kneser germ, so the identity
theorem identifies the family throughout $W$.
\end{proof}

\begin{corollary}[A full exterior height corridor]
\label{cor:exterior-full-height}
At the prescribed non-cusp Brjuno parameter $\theta_c$ there
are a complex parameter neighbourhood $U_*$ and a connected
open height neighbourhood $W\supset[0,1]$ on which the exterior
Kneser family continued from real $b>\eta$ is jointly holomorphic
and satisfies $F_\theta(0)=1$ and $F_\theta(1)=b(\theta)$.
The parameter neighbourhood crosses the Shell--Thron boundary.
No comparison with the inner sewn family $\Kk^W$ is asserted.
\end{corollary}

\begin{proof}
Proposition~\ref{prop:exterior-seam-monotonicity} verifies the
unique transverse seam crossing at $\theta_c$.
Apply Proposition~\ref{prop:exterior-single-crossing}.
\end{proof}

\begin{corollary}[Full exterior height along an interior-reaching path]
\label{cor:exterior-wide-full-height}
Put $p_0=0.1000435842$ and
$J=\{p_0+\ii q:0\le q\le0.01\}$.  There are a connected
complex parameter neighbourhood $\Omega\supset J$ and a connected
height neighbourhood $W_J\supset[0,1]$ such that the exterior
Kneser branch continued from $q=0$ is jointly holomorphic on
$\Omega\times W_J$, the map $\theta\mapsto b(\theta)$ is
injective on $\Omega$, and the family satisfies
$F_\theta(0)=1$, $F_\theta(1)=b(\theta)$.
In particular its full-height branch is defined at the strictly
interior parameter $\theta_0=p_0+0.01\ii$ of
Proposition~\ref{prop:interior-chord-overlap}.
\end{corollary}

\begin{proof}
The wide-box certificate of
Proposition~\ref{prop:exterior-analytic-germ-corridor} supplies
the marked pullback cell, its analytic seam identity, and a jointly
holomorphic normalised germ throughout a neighbourhood of $J$.
Proposition~\ref{prop:exterior-seam-monotonicity} supplies a unique
transverse seam crossing at every point of $J$.  The compact-arc
development in the proof of
Proposition~\ref{prop:exterior-single-crossing} therefore applies
locally at each such parameter; it uses the Jordan cell and its
puncture tails, not the sign of $|\lambda_+|-1$.
Finitely many resulting parameter charts cover $J$.  Their
full-height families agree on overlaps because all extend the
same jointly holomorphic germ at $z=0$; the identity theorem
then extends the agreement along $[0,1]$.  Choose a thin connected
parameter neighbourhood $\Omega$ of $J$ inside their union and
inside the injectivity domain $U$ of
Proposition~\ref{prop:exterior-analytic-germ-corridor}, and
a common thin connected tube $W_J$ about $[0,1]$ inside their
finitely many height domains.  Patching gives the assertion.
\end{proof}

\begin{corollary}[The exterior zero label at the boundary]
\label{cor:exterior-boundary-zero}
After shrinking $U_*$ in Corollary~\ref{cor:exterior-full-height},
the exterior family extends jointly holomorphically to
$U_*\times W_-$ for a connected height neighbourhood
$W_-\supset[-1,1]$.  Throughout this parameter neighbourhood,
\[
 F_\theta(-1)=0,\qquad
 \partial_zF_\theta(-1)=
   \frac{\partial_zF_\theta(0)}{\log b(\theta)}\ne0.
\]
Its next local backward branch at $z=-2$ has the logarithmic
monodromy $2\pi\ii/\log b(\theta)$ of
Proposition~\ref{prop:backward-log-monodromy}.
Thus the exterior boundary germ has the same zero backward label
as the real inner sewn family; this label alone does not identify
the two families.
\end{corollary}

\begin{proof}
The proof of Proposition~\ref{prop:exterior-single-crossing} applies
locally at every point of the compact marked path, since
Proposition~\ref{prop:exterior-seam-monotonicity} gives the same
strict crossing throughout it.  The resulting full-height charts
agree on overlaps: they contain the transported germ at $z=0$,
and the identity theorem extends that agreement along $[0,1]$.
They therefore give a continuous family of height paths
$F_\theta([0,1])$ from the real seed to $\theta_c$.

These paths avoid $0$.  Indeed, each developing representative
lies either in $\Phi_\theta(H_\theta)$ or in
$E_{b(\theta)}(\Phi_\theta(H_\theta))$.
The second set omits $0$ because the exponential map does.
If $\Phi_\theta(w)=0$ for $w\in\overline{H_\theta}$, then
$w=E_{b(\theta)}^{\circ28}(0)$, contrary to the certified
singular-orbit clearance of
Proposition~\ref{prop:exterior-pullback-cell}.
The compact height paths and the strict clearance allow a common
small simply connected tube $V\supset[0,1]$, after restricting
the parameter to a neighbourhood of $\theta_c$, on which
$F_\theta$ has no zeros.

Choose its jointly holomorphic logarithm $L_\theta(z)$ on
$U_*\times V$ by $L_\theta(0)=0$, and choose the local
holomorphic branch $\ell(\theta)=\log b(\theta)$ from the real
exterior seed.  Since $F_\theta(1)=b(\theta)$,
$[L_\theta(1)-\ell(\theta)]/(2\pi\ii)$ is an integer, constant
along the marked path.  It is zero at the real seed, where
$F_\theta([0,1])$ is positive.  Hence $L_\theta(1)=\ell(\theta)$
near $\theta_c$.

For $z\in V-1$ put
$\widetilde F_\theta(z)=L_\theta(z+1)/\ell(\theta)$.
On the connected overlap with $V$, the Abel equation and the
choice at $z=0$ give $\widetilde F_\theta=F_\theta$;
shrink the tube if needed so this overlap is connected.
They glue on $W_-=V\cup(V-1)$ and give
$F_\theta(-1)=L_\theta(0)/\ell(\theta)=0$.
Differentiation gives the displayed nonzero derivative because
the exterior height chart is locally conformal at $z=0$.
The final assertion is
Proposition~\ref{prop:backward-log-monodromy}.
\end{proof}

\begin{corollary}[Exterior root isolation and local normality]
\label{cor:exterior-boundary-normal}
There are a parameter disc $U_0\Subset U_*$ about $\theta_c$,
a radius $r>0$, and constants $0<c<C$ such that, for every
$\theta\in U_0$ and $|z|\le r$,
\[
 c|z|\le |F_\theta(z)-1|\le C|z|,
 \qquad c\le |F_\theta(z)|\le C.
\]
After reducing $r$, the same constants can be chosen so that
$c|z+1|\le |F_\theta(z)|\le C|z+1|$ for $|z+1|\le r$.
Joint holomorphy also gives the exterior half of condition (P1)
on every $V\Subset W_-$ across this boundary subarc, while the
displayed estimates isolate its normalised root uniformly.
No corresponding inner estimate is asserted.
\end{corollary}

\begin{proof}
The jointly holomorphic quotient
$A(\theta,z)=(F_\theta(z)-1)/z$ extends across $z=0$ with
$A(\theta_c,0)=\partial_zF_{\theta_c}(0)\ne0$.
Shrink $U_0$ and $r$ so that $A$ is bounded above and away from
zero on $\overline{U_0}\times\overline{B(0,r)}$.
The same compactness argument applies to $F_\theta$ near
$F_{\theta_c}(0)=1$ and to
$F_\theta(z)/(z+1)$ near $z=-1$, using
Corollary~\ref{cor:exterior-boundary-zero} for its nonzero value there.
Enlarge $C$ and decrease $c$ to cover all three inequalities.
\end{proof}

\Needspace{18\baselineskip}
\begin{lemma}[The integer label at either vertical end]\label{lem:end-spectrum}
Let $E_b(w)=b^w$ have a fixed point $L$ admitting a local Koenigs
coordinate $\sigma$, with multiplier $\lambda$.  Choose
$\varepsilon\in\{+1,-1\}$ and $\ell=\log\lambda$ so that
$0<\varepsilon\operatorname{Im}\ell<2\pi$.
Suppose $F$ is nonconstant and holomorphic on a sufficiently high
$\varepsilon$-half-plane, satisfies $F(z+1)=E_b(F(z))$, and
$F(x+\varepsilon\ii t)\to L$ uniformly for $0\le x\le1$ as
$t\to+\infty$.  Then, deep in that end and in a fundamental vertical
strip, there are an integer $m\ge0$ and a holomorphic function $h$
near $q=0$, with $h(0)\ne0$, such that
\[
 \sigma(F(z))=\eee^{\ell z}q^m h(q),\qquad
 q=\eee^{\varepsilon2\pi\ii z}.
\]
The Koenigs logarithm consequently has the form
\[
 \frac{\log\sigma(F(z))}{\ell}
 =\left(1+\frac{\varepsilon2\pi\ii m}{\ell}\right)z
  +\frac{\log h(q)}{\ell}.
\]
The degree-one form $z+$ a bounded holomorphic $1$-periodic
correction at this end is equivalent to $m=0$.
For a jointly holomorphic parameter family with a common such end
and holomorphically chosen $L$, $\sigma$ and $\ell$, the constant
Laurent coefficient $a_0(b)$ is holomorphic in $b$ (and equals
$h_b(0)$ when $m=0$).  On a connected parameter domain, if nonzero
at one base, its zeros are isolated.
\end{lemma}

\begin{proof}
For large $t$, the values $F(x+\varepsilon\ii t)$, $0\le x\le1$,
lie in the Koenigs chart.  There
$g(z)=\eee^{-\ell z}\sigma(F(z))$ is $1$-periodic,
so it descends to a holomorphic function on a punctured $q$-disk.
Its Laurent coefficient at index $n$ is
\[
 a_n=\int_0^1\eee^{-\ell(x+\varepsilon\ii t)}
       \sigma(F(x+\varepsilon\ii t))
       \eee^{-\varepsilon2\pi\ii n(x+\varepsilon\ii t)}\,dx.
\]
For $n<0$ its modulus is bounded by a constant times
$\sup_{0\le x\le1}|\sigma(F(x+\varepsilon\ii t))|
 \eee^{(\varepsilon\operatorname{Im}\ell+2\pi n)t}$, which tends to zero.
Thus every negative Laurent coefficient vanishes and $g$ extends
holomorphically to $q=0$.  It is not identically zero, so it has
finite order $m\ge0$ there and $g=q^m h(q)$ with $h(0)\ne0$.
Taking a logarithm on the deep end gives the displayed formula.
For the parameter assertion, use the same coefficient integral on a
fixed sufficiently deep line locally in $b$; it is holomorphic in
$b$.  The identity theorem makes its zeros isolated unless it
vanishes identically, which the nonzero seed excludes.
\end{proof}

A fixed-point limit at either end therefore does not, by itself,
select the degree-one time used in Definition~\ref{def:class}.
This is an actual ambiguity for a repelling fixed point: the inverse
Koenigs map $P$ extends to an entire Poincar\'e function, and for every
$m\ge0$ and $c\ne0$ the entire solution
$F_m(z)=P(c\eee^{(\ell+\varepsilon2\pi\ii m)z})$ satisfies the
functional equation and tends to $L$ uniformly on
$0\le\operatorname{Re}z\le1$ as
$\varepsilon\operatorname{Im}z\to+\infty$, with exactly the label $m$.
These examples impose one end condition; they do not assert the
normalisation or the limit at the other end.  The two integer labels
and their nonvanishing leading coefficients are branch data for a
two-regular-chart identification.

\begin{corollary}[The exterior branch has degree-one time at both ends]
\label{cor:exterior-end-labels}
At every real $b>\eta$, Kneser's classical solution has $m=0$ in
Lemma~\ref{lem:end-spectrum} at both conjugate fixed-point ends.
Its two leading coefficients are nonzero and vary holomorphically
and without zeros in the local parameter families of
Proposition~\ref{prop:exterior-real-parameter}.
\end{corollary}

\begin{proof}
In the compactified cylinder of that proposition, the local puncture
coordinates are
$\xi_+=\exp(2\pi\ii\log\sigma_+/\ell_+)$ and
$\xi_-=\exp(-2\pi\ii\log\sigma_-/\ell_-)$.  The global sphere
coordinate $q_b$ has a simple zero at the upper end and a simple pole
at the lower end.  Consequently
$q_b=c_+(b)\xi_+(1+O(\xi_+))$ and
$q_b^{-1}=c_-(b)\xi_-(1+O(\xi_-))$, with $c_\pm(b)$ holomorphic and
nonzero.  Taking logarithms and using
$z=(2\pi\ii)^{-1}\log q_b$ shows at either end that
$\log\sigma_\pm/\ell_\pm=z+C_\pm(b)+O(e^{\varepsilon2\pi\ii z})$,
where $\varepsilon=+1$ at the upper end and $-1$ at the lower end.
Thus the coefficient of $z$ is one, so the label is zero.  The same
local expansions give nonvanishing holomorphic leading coefficients.
\end{proof}

\begin{corollary}[Parameter-uniform exterior tails]
\label{cor:exterior-uniform-tails}
Fix $b_0>\eta$.  There are a complex disc $\Delta$ about $b_0$,
$\delta,T,\rho>0$, and functions
$h_\pm\in\mathcal O(\Delta\times B(0,\rho))$ which are nowhere zero
after shrinking the disc and radius, such that the local exterior
family satisfies
\[
 \sigma_{\pm,b}(\kappa_b(z))
   =\eee^{\ell_{\pm,b}z}
      h_\pm(b,\eee^{\pm2\pi\ii z})
 \quad\text{for }-\delta<\operatorname{Re}z<1+\delta,
       \quad\pm\operatorname{Im}z>T.
\]
For every smaller closed base disc $\Delta'\Subset\Delta$ and every
$0<c<\inf_{b\in\Delta'}\{\operatorname{Im}\ell_{+,b},
-\operatorname{Im}\ell_{-,b}\}$, a constant $C$ gives
\[
 \sup_{\substack{b\in\Delta'\\0\le x\le1}}
 |\kappa_b(x\pm\ii t)-L_\pm(b)|\le C\eee^{-ct}
 \qquad(t>T).
\]
The same estimate holds uniformly over a complex neighbourhood of
any compact subinterval of $(\eta,\infty)$, with a possibly smaller
common $c>0$.
\end{corollary}

\begin{proof}
In the proof of Corollary~\ref{cor:exterior-end-labels}, the global
coordinate and the upper puncture coordinate satisfy
$q_b=\phi_+(b,\xi_+)$, where $\phi_+(b,0)=0$ and
$\partial_{\xi}\phi_+(b,0)\ne0$.  The joint inverse function theorem
therefore gives $\xi_+=q_b\psi_+(b,q_b)$ with $\psi_+$ holomorphic and
nonzero near $q_b=0$.  Since $q_b=\eee^{2\pi\ii z}$, taking a
holomorphic logarithm of $\psi_+$ and then exponentiating the relation
between $\xi_+$ and $\sigma_{+,b}$ yields the asserted $h_+$; the
constant integer choice of the lift is absorbed into it.  At the
lower end apply the same argument to
$q_b^{-1}=\phi_-(b,\xi_-)$, with
$\xi_-=\eee^{-2\pi\ii\log\sigma_{-,b}/\ell_{-,b}}$.
Choose $\rho$ and $T$ uniformly after shrinking $\Delta$.  On
$\Delta'$, both $h_\pm$ and their reciprocals are bounded on a
smaller closed $q$-disc, and
$|\eee^{\ell_{\pm,b}(x\pm\ii t)}|\le C_0\eee^{-c_0t}$
for $0\le x\le1$ and some $c_0>c$.  The inverse Koenigs maps satisfy
$\sigma_{\pm,b}^{-1}(\zeta)=L_\pm(b)+\zeta+O(\zeta^2)$ uniformly
there, giving the estimate.  A finite subcover of a compact real base
interval gives common parameters and a positive minimum decay rate.
\end{proof}

\Needspace{17\baselineskip}
\begin{proposition}[A normal-family criterion for base continuation]
\label{prop:approx-continuation}
Let $\Omega$ and $W$ be connected open sets in the base and height planes,
respectively, and suppose $I=\Omega\cap(\eta,\infty)$ contains a nonempty
interval.  Let $D\subset W$ have an accumulation point in $W$.
Suppose $A_N\in\mathcal O(\Omega\times W)$ is locally uniformly bounded
in $N$ and, for every $(b,z)\in I\times D$,
$A_N(b,z)\to\kappa_b(z)$, where $\kappa_b$ is Kneser's classical solution
on $W$.  Then $A_N$ converges locally uniformly on $\Omega\times W$
to a jointly holomorphic family $A$, and $A_b=\kappa_b$ on $W$ for every
$b\in I$.  If, in addition, $U_0\subset\Omega$ is a nonempty open set
where the sewn family $\Kk^W$ is jointly holomorphic on $U_0\times W$,
and $A_N(b,z)\to\Kk_b^W(z)$ for $(b,z)\in U_0\times D$, then
$A=\Kk^W$ on $U_0\times W$.
\end{proposition}

\begin{proof}
Local boundedness makes the sequence normal on $\Omega\times W$.
Any two subsequential limits agree on $I\times D$ by the prescribed
pointwise convergence.  For each $b\in I$ the identity theorem in $z$
makes them agree on $\{b\}\times W$; for each $z\in W$ the identity
theorem in $b$ then makes them agree throughout $\Omega\times\{z\}$.
Hence the normal sequence has a unique cluster limit and converges
locally uniformly to it.  Applying the identity theorem in $z$ once
more identifies $A_b$ with $\kappa_b$ for $b\in I$ and, under the
additional assumption, with $\Kk_b^W$ for $b\in U_0$.
\end{proof}

This gives an alternative to checking (P0)--(P2) one boundary point at a
time: a single locally bounded sequence of jointly holomorphic
approximants on a corridor crossing the boundary would provide the
continuation and, if its inner limit were verified, the branch
identification.  At each fixed finite stage, the quadrature formulas are locally
holomorphic in $b$ on a cell where all logarithm and inverse-chart
branches are fixed and the normalising denominator is nonzero.  A
common branch-safe corridor for every stage, and a bound uniform in
quadrature order and iteration count, have not been established for
the approximants of~\cite[\S5]{Paulsen2019}.  Real-base convergence by itself
cannot supply the bound: for real $b_0$ the entire functions
$a_N(b)=\exp[-N-N^3(b-b_0)^2]$ tend to zero at every real $b$, but
are unbounded along every vertical line $b=b_0+\ii y$ with $y\ne0$.

\subsection{Reduction of the continuation to Brjuno boundary points}\label{sec:brjuno}
We record how the one remaining hypothesis reduces to statements at boundary
points with Brjuno rotation number.  Let $\Gamma$ be an open arc of the
Shell--Thron boundary in the upper half-plane.  Parametrise it by
$\lambda_c=\eee^{2\pi\ii\alpha}$, where $\lambda_c$ is the multiplier of the fixed
point $L_1$.  Let $U=U_{\mathrm{out}}\cup\Gamma\cup U_{\mathrm{in}}$ be a
neighbourhood of $\Gamma$ split by it.

\begin{lemma}[Radial small divisors]\label{lem:radial}
For $\lambda=r\lambda_c$ with $r>0$ and every $m\ge1$,
$|\lambda^m-1|\ge\min\bigl(\tfrac1{\sqrt2}|\lambda_c^m-1|,\ \tfrac12\bigr)$.
\end{lemma}

\begin{proof}
Write $\lambda_c^m=\eee^{\ii\theta}$ with $|\theta|\le\pi$.  The distance from $1$
to the ray $\{t\eee^{\ii\theta}:t>0\}$ equals $|\sin\theta|$ if
$|\theta|\le\pi/2$ and $1$ otherwise.  In the first case,
$|\sin\theta|=|\eee^{\ii\theta}-1|\cos(\theta/2)\ge|\eee^{\ii\theta}-1|/\sqrt2$.  In
the second, $|\eee^{\ii\theta}-1|\le2$.
\end{proof}

The local germ has a useful exact form.  Since
$b=\exp(\lambda\eee^{-\lambda})$ and $L=\eee^\lambda$, the coordinate
$v=(\log b)(w-L)$ transforms $w\mapsto b^w$ into
\begin{equation}\label{eq:radial-normal-form}
  f_\lambda(v)=\lambda(\eee^v-1).
\end{equation}

\begin{lemma}[Radial convergence of local charts]\label{lem:radial-chart}
Let $\lambda_c=\eee^{2\pi\ii\alpha}$ with $\alpha$ Brjuno, and put
$\lambda=r\lambda_c$.  The normalised inverse Schr\"oder charts
$H_\lambda$, defined by
$f_\lambda\circ H_\lambda=H_\lambda\circ(\lambda\,\cdot)$ and
$H_\lambda(0)=0$, $H_\lambda'(0)=1$, are holomorphic and uniformly
bounded on one fixed disk $|v|<\rho$ for $r$ sufficiently close to $1$
from either side.  They converge locally uniformly there to the Siegel
chart $H_{\lambda_c}$.  On a possibly smaller fixed disk, their inverses
have the same convergence property.
\end{lemma}

\begin{proof}
The map $v\mapsto\eee^v-1$ is univalent on $|v|<1$.  Thus the
cone-uniform Brjuno majorant estimate of~\cite{MarmiCarminati2008}
applies directly to~\eqref{eq:radial-normal-form}; the radial paths lie in
such a cone.  Lemma~\ref{lem:radial} gives the requisite uniform lower
bounds for the small divisors.  Each Taylor coefficient of $H_\lambda$ is
rational in $\lambda$ with poles only at roots of unity, hence converges to
its value at the irrational $\lambda_c$.  The uniform majorant makes this
coefficientwise convergence locally uniform on a fixed disk.  Since
$H_\lambda'(0)=1$, Cauchy estimates on that disk give a smaller disk on
which the $H_\lambda$ are uniformly close to the identity and univalent;
their inverses then converge locally uniformly as well.
\end{proof}

This lemma concerns germs at $L$.  Neither the Siegel disk nor its
finite-forward-pullback domain supplies a complete sewing line
(Proposition~\ref{prop:siegel-pullback}).

\begin{lemma}[The other fixed points are repelling]
\label{lem:boundary-repelling}
For $b=\exp(\lambda\eee^{-\lambda})$ with $0<|\lambda|\le1$ and
$\lambda\ne1$, every fixed point of $E_b$ other than $L_1=\eee^\lambda$
is simple and has a multiplier of modulus greater than $1$.
\end{lemma}

\begin{proof}
Put $a=\log b=\lambda\eee^{-\lambda}$ and
$g(\mu)=\mu\eee^{-\mu}$.  Fixed points $w=E_b(w)$ correspond to
$g(\mu)=a$ through $\mu=aw$, $w=\eee^\mu$; $\mu$ is their multiplier.
The chosen logarithm is indeed $a$, since $|a|\le\eee<\pi$.
We show that $g$ is injective on the closed unit disk.  In polar
coordinates,
\[
 \arg g(r\eee^{\ii\theta})=\theta-r\sin\theta,\qquad
 |g(r\eee^{\ii\theta})|=r\eee^{-r\cos\theta}.
\]
For each $0<r\le1$, the first expression is strictly increasing in
$\theta$ modulo $2\pi$.  Holding its value fixed and varying $r<1$,
implicit differentiation gives
\[
 \frac{d}{dr}\log|g(r\eee^{\ii\theta(r)})|
 =\frac{|1-r\eee^{\ii\theta(r)}|^2}
        {r(1-r\cos\theta(r))}>0.
\]
Thus different radii cannot have the same image, including by
continuity at $r=1$.  Since $\lambda$ is already a preimage of $a$ in
the closed unit disk, every other preimage has modulus greater than $1$.
Finally $g'(\mu)=\eee^{-\mu}(1-\mu)$ vanishes only at $\mu=1$,
which cannot be another preimage when $\lambda\ne1$.
\end{proof}

\begin{proposition}[Transport of the lower fixed point]
\label{prop:lower-transport}
Fix $\lambda_0\in(0,1)$ and its real second multiplier $\mu_0>1$,
where $\mu_0\eee^{-\mu_0}=\lambda_0\eee^{-\lambda_0}$.
Let $\gamma:[0,1]\to\overline{\mathbb D}\setminus\{0,1\}$ be a
piecewise $C^1$ path with $\gamma(0)=\lambda_0$,
$\gamma([0,1))\subset\mathbb D$, and
$\gamma(1)=\lambda_c\in\partial\mathbb D\setminus\{1\}$.
Then $\mu_0$ continues analytically along $\gamma$ to a multiplier
$\mu_c$ with $|\mu_c|>1$.  The associated lower fixed point and its
regular superfunction continue along a narrow parameter corridor and
extend locally jointly holomorphically in the base and $z$ through
$b_c=\exp(\lambda_c\eee^{-\lambda_c})$, with a transported branch of
$\log\mu$.
\end{proposition}

\begin{proof}
Put $g(t)=t\eee^{-t}$.  At every point of the path, the implicit
equation $g(\mu)=g(\lambda)$ has a locally unique continuation of the
second root as long as $\mu\ne1$.  Lemma~\ref{lem:boundary-repelling}
keeps it outside the closed unit disk, so it cannot meet the root
$\lambda$.  On the compact trace $K=\gamma([0,1])$, which avoids $1$,
there is a $\delta>0$ such that every such second root obeys
$|\mu-1|\ge\delta$: otherwise $\mu_j\to1$ and
$\lambda_j\to\lambda_*\in K$ would give $g(\lambda_*)=g(1)$,
contradicting injectivity of $g$ on the closed unit disk.
Along each smooth piece of $\gamma$,
\[
 \frac{d\mu}{d\lambda}
 =\frac{g'(\lambda)}{g'(\mu)}
 =\frac{1-\lambda}{\lambda}\,\frac{\mu}{1-\mu}.
\]
Here $|\lambda|$ is bounded below on $K$, and
$|\mu/(1-\mu)|\le1+1/\delta$.  The derivative is therefore bounded
by a constant times the path speed.  At any hypothetical first stopping
point, finite path length gives a finite limit $\mu\ne1$; the
implicit-function theorem would extend the branch.  Thus it reaches
$\lambda_c$.  It remains outside the unit disk by
Lemma~\ref{lem:boundary-repelling}: it cannot merge with $\lambda_c$,
because $g'(\lambda_c)\ne0$ would make the local root uniquely equal to
$\lambda$ near $\lambda_c$, contrary to $|\mu|>1$ inside the disk.
The lower Koenigs construction and
the proof of Lemma~\ref{lem:lower-chart} apply in neighbourhoods along
the path and at its endpoint.  Their normalisation and the continued
branch of $\log\mu$ patch these local superfunctions.  Finally
$d\exp(g(\lambda))/d\lambda\ne0$ at $\lambda_c\ne1$, so the endpoint
chart may be expressed in the base $b$.
\end{proof}

This transports the lower regular chart only.  It does not transport
the sewing correction or the solution $\Kk^W$ to the boundary arc.

\begin{corollary}[A global lower regular family in the inner parameter disk]
\label{cor:global-lower}
Let $\Omega=\mathbb D\setminus(-1,0]$.  The real second multiplier
$\mu_2>1$ extends to a single-valued holomorphic function on $\Omega$
with $|\mu_2(\lambda)|>1$.  It has a holomorphic logarithm fixed by the
positive real value on $(0,1)$.  With this choice, the lower regular
superfunction $S_\lambda(z)$ is jointly holomorphic on
$\Omega\times\mathbb C$ and extends locally across every point of
$\partial\mathbb D\setminus\{1,-1\}$.
\end{corollary}

\begin{proof}
The continuation argument of Proposition~\ref{prop:lower-transport}
works along every finite path in $\Omega$.  Since $\Omega$ is simply
connected, the monodromy theorem makes the continuation single-valued.
Lemma~\ref{lem:boundary-repelling} gives $|\mu_2|>1$, hence a global
holomorphic logarithm.  At each parameter, the Koenigs coordinate at
$L_2=\eee^{\mu_2}$, normalised by derivative $1$, is unique and varies
locally holomorphically; so the local constructions of
Lemma~\ref{lem:lower-chart} patch.  Its forward-iterate formula makes
$S_\lambda$ entire in $z$ and jointly holomorphic locally in
$(\lambda,z)$.  Proposition~\ref{prop:lower-transport} gives the stated
local boundary extensions.
\end{proof}

For $0<|\lambda|<1$ the normalised attracting Koenigs germ also varies
locally holomorphically in $\lambda$, since its coefficient equations
have no interior resonances.  Thus both local fixed-point charts can be
transported through $\Omega$; persistence of the overlap needed for
their sewing remains unproved.

\begin{proposition}[Infinite monodromy around $b=1$]
\label{prop:lower-monodromy}
The map $\lambda\mapsto b=\exp(\lambda\eee^{-\lambda})$ is injective on
$\mathbb D$.  On the universal cover of $0<|\lambda|<1$, the lower
multiplier is single-valued and holomorphic, and its regular
superfunction is jointly holomorphic in the parameter and height.
The lower multiplier branch of
Corollary~\ref{cor:global-lower} has infinite-order monodromy around
$b=1$: continuing it around any positively oriented loop that winds
once around $b=1$ in the image of $0<|\lambda|<1$ changes its value,
and its values after different integer winding numbers are distinct.
Consequently neither the lower fixed point nor its normalised lower
regular superfunction is single-valued on that punctured base domain.
\end{proposition}

\begin{proof}
Injectivity of $g(\lambda)=\lambda\eee^{-\lambda}$ on $\mathbb D$ was
proved in Lemma~\ref{lem:boundary-repelling}.  Since $|g(\lambda)|<\eee<\pi$,
the exponential is injective on $g(\mathbb D)$ as well.  Moreover
$b'(\lambda)=b(\lambda)\eee^{-\lambda}(1-\lambda)\ne0$ in $\mathbb D$,
so this is a conformal parameterisation.
The continuation proof of Proposition~\ref{prop:lower-transport}
applies to every finite path in the punctured disk.  Monodromy on its
simply connected universal cover gives a holomorphic root $\mu$.
Since $|\mu|>1$, it has a global holomorphic logarithm there.
The local construction of Lemma~\ref{lem:lower-chart} then patches
to give the asserted superfunction.

Choose $0<\varepsilon<1$ so small that
$R:=-\log\varepsilon>\varepsilon$, and continue the real lower root
$\mu_0>1$ along $\lambda(t)=\varepsilon\eee^{\ii t}$, with $t$
running from $0$ to $2\pi n$ for any integer $n$.  The
path-continuation proof of
Proposition~\ref{prop:lower-transport} applies to every finite number
of turns.  All continued roots obey $|\mu(t)|>1$.  Their common
modulus equation gives
\[
 \operatorname{Re}\mu(t)
 =R+\operatorname{Re}\lambda(t)+\log|\mu(t)|
 >R-\varepsilon>0.
\]
Thus the principal $\operatorname{Log}\mu(t)$ remains defined throughout.  Taking
continuous logarithms in $g(\mu(t))=g(\lambda(t))$ and using the real
equation at $t=0$ yields
\[
 \mu(t)-\operatorname{Log}\mu(t)
   =R-\ii t+\varepsilon\eee^{\ii t}.
\]
At the end of $n$ turns, the right side is
$R+\varepsilon-2\pi\ii n$.  Since the same single-valued function
$\mu\mapsto\mu-\operatorname{Log}\mu$ is used at every endpoint, roots obtained
after different winding numbers cannot coincide.  A radial path from
any real $\lambda_0\in(0,1)$ to $\varepsilon$ conjugates this
monodromy to the loop based at $\lambda_0$.

Distinct endpoint roots give distinct fixed points $L_2=\eee^\mu$,
because $\mu=(\log b)L_2$, and distinct regular superfunctions,
because $S(x)\to L_2$ as $x\to-\infty$.  Homotopy invariance extends
the conclusion to every loop of the same winding number.
\end{proof}

\begin{proposition}[The upper chart at the normalising orbit]
\label{prop:upper-orbit}
Let $0<|\lambda|<1$, $a=\lambda\eee^{-\lambda}$,
$b=\eee^a$, and $L_1=\eee^\lambda$.  Then
$E_b^{\circ n}(1)\to L_1$ locally uniformly in $\lambda$.
The attracting Koenigs coordinate $\sigma_\lambda$, normalised by
$\sigma_\lambda(L_1)=0$ and $\sigma_\lambda'(L_1)=1$, extends to the
attracting basin of $L_1$ and depends locally jointly holomorphically
on $(\lambda,w)$ there.  In particular
$c(\lambda)=\sigma_\lambda(1)$ is holomorphic and nonzero on the
punctured unit disk, and $\sigma_\lambda'(1)\ne0$.  It extends across
$\lambda=0$ with
\[
 c(\lambda)=-\lambda-\tfrac12\lambda^2+O(\lambda^3).
\]
Consequently there is a unique holomorphic function $h$ on $|\lambda|<1$
with $h(0)=0$ and $c(\lambda)=-\lambda\exp h(\lambda)$; in particular
$h(\lambda)=\lambda/2+O(\lambda^2)$.

On the universal cover of $0<|\lambda|<1$, the normalised upper regular
superfunction germ
\[
 R_\lambda(z)=\sigma_\lambda^{-1}
       \bigl(c(\lambda)\exp(z\log\lambda)\bigr),
 \qquad R_\lambda(0)=1,
\]
is jointly holomorphic near each parameter and $z=0$.  Along every
compact parameter path avoiding $0$ it has a common positive
height radius on a sufficiently narrow corridor.  On a possibly
smaller common disk, $R_\lambda(z)=1$ only at $z=0$, and
$R_\lambda$ and $1/R_\lambda$ are uniformly bounded along the path.
The deck transformation around $\lambda=0$ acts with infinite order on
these normalised upper germs.
\end{proposition}

\begin{proof}
Under the change of coordinate $v=aw$, the map $E_b$ becomes
$v\mapsto a\eee^v$, whose attracting fixed point is $v=\lambda$ and
whose only finite singular value is $0$.  The singular-value theorem
for exponential maps places $0$ in the immediate attracting basin
of this fixed point~\cite[\S2]{Schleicher2003}.  Since $0\mapsto a$
in the $v$ coordinate, $w=1$ belongs to the same basin.
Attraction is locally uniform in the parameter: at any fixed
$\lambda_0$, a finite iterate of $1$ enters a small invariant
contraction disk about $L_1(\lambda_0)$, and this persists for nearby
parameters.

The local Koenigs coordinate varies holomorphically with $\lambda$.
For a basin point $w$, choose $N$ with $E_b^{\circ N}(w)$ in its local
domain and define
\[
 \sigma_\lambda(w)=\lambda^{-N}
       \sigma_\lambda\bigl(E_b^{\circ N}(w)\bigr).
\]
This formula is independent of larger $N$ and locally jointly
holomorphic.  Its $w$ derivative never vanishes, since $E_b$ has no
critical points and the local Koenigs coordinate is univalent.
The local coordinate has only the zero $L_1$, so $c(\lambda)=0$
exactly when a finite iterate of $1$ lands on $L_1$.  In that case
the singular orbit of $v\mapsto a\eee^v$ would be finite.
For a postsingularly finite exponential map, a preperiodic dynamic
ray lands at its singular value~\cite[Thm.~3.1]{LaubnerSchleicherVicol2008};
its points escape, so it cannot land at a point in the open attracting
basin.  This contradicts the membership of $0$ in that basin proved above.
Thus $c$ never vanishes.

To analyse the removable endpoint $\lambda=0$, write $u=w-L_1$ and
\[
 G_\lambda(u)=E_b(L_1+u)-L_1
   =\eee^\lambda(\exp(\lambda\eee^{-\lambda}u)-1)
   =\lambda u+O(\lambda^2u^2).
\]
On a fixed small $u$-disk and for sufficiently small $|\lambda|$,
$|G_\lambda(u)|\le C|\lambda||u|$ with $C^2|\lambda|<1$.
Put $u_n=G_\lambda^{\circ n}(u)$.  Each
$\lambda^{-n}u_n$ has a removable singularity at $\lambda=0$, and
\[
 \left|\lambda^{-n-1}u_{n+1}-\lambda^{-n}u_n\right|
 \le C'|\lambda|(C^2|\lambda|)^n|u|^2.
\]
The telescoping limit is the normalised Koenigs coordinate and extends
jointly holomorphically through $\lambda=0$ with
$\sigma_\lambda(L_1+u)=u+O(\lambda u^2)$.
Since $1-L_1=-\lambda-\lambda^2/2+O(\lambda^3)$, this gives the
stated expansion of $c$.  The function $-c(\lambda)/\lambda$ is now
holomorphic and zero-free on the whole unit disk, with value $1$ at the
origin.  It therefore has the asserted normalised holomorphic logarithm.

Choose a holomorphic branch of $\log\lambda$ on the
universal cover.  The inverse-function theorem at $w=1$, applied to
$\sigma_\lambda'(1)\ne0$, gives the displayed germ and its joint
holomorphy.  A finite cover of a compact lifted parameter path gives
one positive height radius after narrowing its corridor.  Since
$R_\lambda'(0)=c(\lambda)\log\lambda/\sigma_\lambda'(1)\ne0$,
the same compactness argument gives a smaller common root-isolation
disk.  Shrink it once more to avoid zeros of $R_\lambda$ and obtain
the stated bounds.  On one positive turn around $0$, $c(\lambda)$ and
$\sigma_\lambda'(1)$ return unchanged, whereas $\log\lambda$ gains
$2\pi\ii$.  After $n$ turns, $R_\lambda'(0)$ therefore changes by
$2\pi\ii n c(\lambda)/\sigma_\lambda'(1)\ne0$ for $n\ne0$.
\end{proof}

\begin{corollary}[Diverging chart-entry depth at a Brjuno boundary point]
\label{cor:entry-depth}
Let $\lambda_c=\eee^{2\pi\ii\alpha}\ne1$ with $\alpha$ Brjuno,
$\lambda=r\lambda_c$, and
$b(r)=\exp(\lambda\eee^{-\lambda})$.  Choose $\rho>0$ so that the
charts of Lemma~\ref{lem:radial-chart} are univalent near
$\{|\zeta|\le\rho\}$ for $r$ close to $1$ from either side, and let
$0<s<\rho$.  Put
\[
 \mathcal D_{r,s}=L_1(b(r))+(\log b(r))^{-1}
                         H_{r\lambda_c}(B(0,s)).
\]
For every fixed $N\ge0$, the point $E_{b(r)}^{\circ N}(1)$ lies outside
$\mathcal D_{r,s}$ for all $r\ne1$ sufficiently close to $1$.
For $r<1$ the first-entry time
\[
 m_s(r)=\min\{n\ge0:E_{b(r)}^{\circ n}(1)\in\mathcal D_{r,s}\}
\]
is finite and tends to infinity as $r\uparrow1$.
\end{corollary}

\begin{proof}
For the boundary map $f_{\lambda_c}(v)=\lambda_c(\eee^v-1)$,
the singular value $-\lambda_c$ belongs to the Julia set whenever
the map has a Siegel disk~\cite[Introduction, remark after
Thm.~2]{Rempe2008}.  Under the coordinate
$v=(\log b_c)(w-L_1(b_c))$, this singular value is $w=0$.
The Julia set is forward invariant, and $E_{b_c}(0)=1$; hence
every $E_{b_c}^{\circ N}(1)$ belongs to the Julia set and lies outside
the boundary Siegel chart $\mathcal D_{1,\rho}$.

If the first assertion failed, there would be $r_j\to1$, $r_j\ne1$,
and $\zeta_j\in B(0,s)$ with
\[
 E_{b(r_j)}^{\circ N}(1)
 =L_1(b(r_j))+(\log b(r_j))^{-1}H_{r_j\lambda_c}(\zeta_j).
\]
After taking a subsequence, $\zeta_j\to\zeta$ with $|\zeta|\le s<\rho$.
Lemma~\ref{lem:radial-chart} and finite-iterate continuity put
$E_{b_c}^{\circ N}(1)$ in $\mathcal D_{1,\rho}$, a contradiction.
For each $r<1$, Proposition~\ref{prop:upper-orbit} says the orbit of $1$
converges to $L_1(b(r))$, so it eventually enters
$\mathcal D_{r,s}$.  The first assertion excludes every bounded entry
time as $r\uparrow1$.
\end{proof}

This concerns finite pullbacks of a fixed local upper chart at the
normalising orbit.  It does not preclude continuation of the sewn
solution by a different mechanism.

\begin{proposition}[Collapse of the attracting basin at a Brjuno boundary point]
\label{prop:basin-inradius}
Let $\lambda_c=\eee^{2\pi\ii\alpha}\ne1$ with $\alpha$ Brjuno,
$b(\lambda)=\exp(\lambda\eee^{-\lambda})$, and let
$\mathcal A_\lambda$ be the immediate attracting basin of
$L_1(\lambda)=\eee^\lambda$ for $0<|\lambda|<1$.  Then
$1\in\mathcal A_\lambda$ and
\[
 d(\lambda):=\operatorname{dist}
       (1,\mathbb C\setminus\mathcal A_\lambda)
 \longrightarrow0
 \qquad(\lambda\to\lambda_c,\ |\lambda|<1).
\]
The Poincar\'e distance within the immediate basin also satisfies
\[
 \operatorname{dist}^{\mathrm{hyp}}_{\mathcal A_\lambda}
       (1,b(\lambda))\longrightarrow\infty
 \qquad(\lambda\to\lambda_c,\ |\lambda|<1).
\]
Moreover, let $\lambda_j\to\lambda_c$ from the disk and let $V$ be a
connected height domain containing $0$.  If holomorphic maps
$G_j:V\to\mathcal A_{\lambda_j}$ satisfy $G_j(0)=1$, then
$G_j\to1$ locally uniformly on $V$, without a boundedness assumption.
Consequently, if $1\in V$, there is a neighbourhood $U$ of $\lambda_c$
such that no $\lambda\in U\cap\{0<|\lambda|<1\}$ admits a holomorphic
$G:V\to\mathcal A_\lambda$ with $G(0)=1$ and $G(1)=b(\lambda)$.
\end{proposition}

\begin{proof}
The proof of Proposition~\ref{prop:upper-orbit} puts $0$ in the
immediate attracting basin for every inner parameter; since
$E_{b(\lambda)}(0)=1$, this basin also contains $1$.  The proof of
Corollary~\ref{cor:entry-depth} puts $1$ in the Julia set of the
boundary map $E_{b(\lambda_c)}$.  Repelling periodic points are dense
in this Julia set~\cite[Thm.~4]{Bergweiler1993}.  Given $\epsilon>0$,
choose such a point $p_c$ with $|p_c-1|<\epsilon/2$.
If its period is $m$, then
$|(E_{b(\lambda_c)}^{\circ m})'(p_c)|>1$, so the implicit-function
theorem continues it to repelling periodic points $p_\lambda$ for
all $\lambda$ near $\lambda_c$.  They lie outside
$\mathcal A_\lambda$ and satisfy $|p_\lambda-1|<\epsilon$.
Thus $d(\lambda)<\epsilon$ near $\lambda_c$.

For the second assertion, first note that $b(\lambda_c)\ne1$: indeed
$0<|\lambda_c\eee^{-\lambda_c}|\le\eee<2\pi$, so its exponential
cannot equal $1$.  Both $1$ and $b(\lambda_c)=E_{b(\lambda_c)}(1)$
belong to the Julia set.  Choose a repelling periodic point
$q_c\ne1$ near $b(\lambda_c)$; it continues to repelling periodic
points $q_j\to q_c$ of $E_{b(\lambda_j)}$.
By the first assertion, choose
$p_j\in\mathbb C\setminus\mathcal A_{\lambda_j}$ with $p_j\to1$.
For all large $j$ the functions
\[
 H_j(z)=\frac{G_j(z)-p_j}{q_j-p_j}
\]
omit both $0$ and $1$, while $H_j(0)\to0$.
Montel's theorem gives spherical normality of this family.  Every
subsequential limit takes the value $0$ at the origin; Hurwitz's theorem
for the zero-free $H_j$ forces that limit to be identically zero.
Hence $H_j\to0$ and $G_j\to1$ locally uniformly on $V$.
Since $\mathcal A_{\lambda_j}$ omits the two distinct points
$p_j,q_j$, it is a hyperbolic domain.
If the asserted Poincar\'e distances stayed bounded along a sequence,
take universal covering maps
$\phi_j:\mathbb D\to\mathcal A_{\lambda_j}$ with $\phi_j(0)=1$.
The lifts of $b(\lambda_j)$ can then be chosen as points $t_j$ in a
fixed compact subdisk of $\mathbb D$.  Applying the preceding collapse
to $\phi_j$ gives
$b(\lambda_j)=\phi_j(t_j)\to1$, contrary to $b(\lambda_c)\ne1$.
If $1\in V$ and $G_j(1)=b(\lambda_j)$, this gives
$b(\lambda_c)=1$, a contradiction.  The neighbourhood assertion follows
by choosing a contrary sequence $\lambda_j\to\lambda_c$.
\end{proof}

This collapse rules out a fixed basin-valued upper chart carrying
both normalising heights $0$ and $1$ near the boundary.  If a normalised
holomorphic solution exists on such a fixed connected height domain,
its image must meet
$\partial\mathcal A_\lambda\subset J(E_{b(\lambda)})$
for all sufficiently close inner parameters.  This does not
decide whether the sewn family reaches such a fixed height domain
near $\lambda_c$; continuation through another chart remains possible.

\begin{lemma}[Parameter stability of the lower chart]
\label{lem:lower-chart}
Let $b_c$ be a point of the boundary arc at which the chosen lower fixed
point $L_2(b_c)$ is simple and its multiplier
$\mu_c=E_{b_c}'(L_2(b_c))$ satisfies $|\mu_c|>1$.  Choose the local
branches of $\log b$ and $\log\mu$ used to define $E_b$ and the lower
regular superfunction.  Near $b_c$, the normalised Koenigs coordinate
$\tau_b$ at $L_2(b)$ and its local inverse depend jointly holomorphically
on the base and their arguments on fixed disks.  Moreover the entire
superfunction
\[
 S_b(z)=E_b^{\circ m}\!\left(\tau_b^{-1}
          \bigl(-\exp((z-m)\log\mu(b))\bigr)\right)
\]
is independent of every sufficiently large integer $m$, locally jointly
holomorphic in $(b,z)$, and converges with all $z$-derivatives locally
uniformly to $S_{b_c}$ as $b\to b_c$ from either side.
\end{lemma}

\begin{proof}
The implicit-function theorem gives holomorphic $L_2(b)$ and $\mu(b)$
near $b_c$.  In the coordinate $u=w-L_2(b)$, set
$G_b(u)=E_b(L_2(b)+u)-L_2(b)$ and let $h_b$ be its local inverse.
Then $h_b'(0)=\mu(b)^{-1}$.  Choose
$|\mu_c|^{-1}<q<|\mu_c|^{-1/2}$.  After shrinking the parameter
neighbourhood and a fixed disk $|u|<r$, we have
$|h_b(u)|\le q|u|$, $|\mu(b)h_b(u)-u|\le C|u|^2$, and
$M q^2<1$, where $M=\sup_b|\mu(b)|$.  Hence
\[
 \tau_b(L_2(b)+u):=\lim_{n\to\infty}\mu(b)^n h_b^{\circ n}(u)
\]
converges normally in $(b,u)$: successive terms differ by at most
$C(Mq^2)^n|u|^2$.  It has derivative $1$ at the fixed point and obeys
$\tau_b(E_b(w))=\mu(b)\tau_b(w)$ wherever both sides are local.
The inverse-function theorem, uniformly after one more shrinking, gives
jointly holomorphic inverses on a fixed disk about $0$.

Since $\operatorname{Re}\log\mu(b)=\log|\mu(b)|$ stays positive, for any
compact set of $(b,z)$ one integer $m$ puts
$-\exp((z-m)\log\mu(b))$ in that inverse disk.  The local conjugacy
makes the displayed formula independent of increasing $m$.  These
formulas therefore define a locally jointly holomorphic function for all
$z\in\mathbb C$.  Holomorphy and Cauchy's formula give the asserted
locally uniform convergence, including all $z$-derivatives.
\end{proof}

\begin{proposition}[Backward-value label of the lower regular family]
\label{prop:lower-backward-label}
In the parameter neighbourhood of Lemma~\ref{lem:lower-chart}, put
$a(b)=\log b\ne0$.  For every $k\in\mathbb Z\setminus\{0\}$ there is a
simple root $c$ of $S_{b_c}(c)=1$ with
$S_{b_c}(c-1)=2\pi\ii k/a(b_c)$.  Each such root continues locally
holomorphically as $c(b)$, and
$G_b(z):=S_b(z+c(b))$ is a jointly holomorphic, entire-in-$z$
solution with $G_b(0)=1$ and
\[
 G_b(-1)=\frac{2\pi\ii k}{\log b}.
\]
Conversely every normalised translate of $S_b$ has one of these nonzero
integer labels, which is constant along local parameter continuation.
In fact every entire-in-$z$ solution with $F(0)=1$ has a nonzero
backward-value label.
For real $1<b<\eta$, the geometric sewn solution satisfies
$\Kk_b^W(-1)=0$; this identity persists on its local complex parameter
discs.  Consequently a normalised lower regular translate cannot
agree with the sewn solution on a common connected regular $z$-domain
containing $-1$.
\end{proposition}

\begin{proof}
The functional equation gives $S_b(z)=E_b(S_b(z-1))\ne0$, so the
nonconstant entire function $S_b$ omits $0$.  Picard's theorem makes
it attain every nonzero value, including $2\pi\ii k/a(b_c)$.
Choose $t$ with that value and put $c=t+1$; then $S_{b_c}(c)=1$.
The formula in Lemma~\ref{lem:lower-chart}, differentiated after
choosing a sufficiently large $m$, shows $S_b'(z)\ne0$: the derivatives
of $E_b^{\circ m}$ and of the local inverse $\tau_b^{-1}$, and the
exponential factor, are all nonzero.  Thus the root $c$ is simple and
the implicit-function theorem continues it holomorphically in $b$.
The displayed value at $-1$ persists by the functional equation and
continuity.  Conversely $S_b(c)=1$ forces
$a(b)S_b(c-1)\in2\pi\ii\mathbb Z$; omission of $0$ excludes the zero
integer.  The integer is locally constant because it depends
holomorphically on $b$.
For an arbitrary entire solution $F$, the same argument applies:
$F(z)=E_b(F(z-1))$ makes $F$ zero-free, so $F(0)=1$ forces
$F(-1)=2\pi\ii k/a(b)$ with $k\ne0$.

For real $b\in(1,\eta)$, the marked upper point $w=0$ of
Proposition~\ref{prop:weld} corresponds to $z=0$, and its deck translate
$w=-1$ corresponds to $z=-1$.  The real attracting regular branch has
$\Rr_b(-1)=0$, since $E_b(0)=1=\Rr_b(0)$ and this inverse branch is real.
Hence $\Kk_b^W(-1)=0$.  Proposition~\ref{prop:weld-parameter} and the
identity theorem preserve the equality on each local parameter disc.
The two values at $-1$ exclude the asserted agreement.
\end{proof}

\begin{corollary}[Rigidity of the backward branch]
\label{cor:backward-rigidity}
Let $U$ be a connected base domain on which $a(b)=\log b$ is holomorphic
and nonzero.  Suppose $F_b(z)$ is jointly holomorphic in $(b,z)$ near
$U\times\{-1,0\}$, satisfies $F_b(0)=1$, and obeys the tetration equation
near $z=-1$.  Then
\[
 k(b):=\frac{a(b)F_b(-1)}{2\pi\ii}\in\mathbb Z
\]
is constant on $U$.  If $k=0$, none of the $F_b$ can extend
holomorphically through $z=-2$ while preserving the tetration equation;
in particular they cannot be entire solutions.  The inner sewn family
has $k=0$ on its real-axis parameter neighbourhood, as does any
continuation of Kneser's real-normalised family from $b>\eta$ that
remains regular at $z=-1$.  Thus their values and all base derivatives
at $z=-1$ agree automatically wherever both branches are defined.
\end{corollary}

\begin{proof}
The equation at $z=-1$ gives $\exp(a(b)F_b(-1))=1$.
The displayed function is holomorphic and integer-valued, hence
constant.  If $k=0$ and $F_b$ extended regularly through $-2$, the
equation would give $E_b(F_b(-2))=F_b(-1)=0$, impossible because an
exponential never vanishes.  Proposition~\ref{prop:lower-backward-label}
gives $k=0$ for the sewn family.  Kneser's family has $F_b(-1)=0$ for
real $b>\eta$ by its real normalisation; the identity theorem transports
this value on any stated parameter continuation.
\end{proof}

\begin{proposition}[Logarithmic monodromy at the first backward singularity]
\label{prop:backward-log-monodromy}
Let $U$ be a simply connected parameter disc on which $a(b)=\log b\ne0$.
Suppose $F_b(z)$ is jointly holomorphic near $U\times\{-1\}$,
$F_b(-1)=0$, and $F_b'(-1)\ne0$.  Locally in $(b,t)$ write
$F_b(-1+t)=t\,u_b(t)$ with $u_b(0)\ne0$.  On any simply connected
sector $0<|t|<r$, all holomorphic backward branches satisfying
$\exp(a(b)G_b(-2+t))=F_b(-1+t)$ have the form
\[
 G_b(-2+t)=\frac{\operatorname{Log}t+\log u_b(t)+2\pi\ii k}{a(b)},
 \qquad k\in\mathbb Z.
\]
Continuation once around $t=0$ adds $2\pi\ii/a(b)$.  Hence no
single-valued holomorphic or meromorphic backward branch satisfying
the equation exists on a full punctured disc about $z=-2$.
For every real $b\in(1,\eta)$, the sewn solution $\Kk_b^W$ has a
simple zero at $z=-1$, so this description applies to it and to its
local holomorphic parameter family.
\end{proposition}

\begin{proof}
The quotient $u_b(t)$ extends jointly holomorphically through $t=0$.
After shrinking $U$ and $r$, it has a joint holomorphic logarithm.
Two logarithms of $t u_b(t)$ on a connected sector differ by a constant
multiple of $2\pi\ii$, giving the formula.  A turn about $t=0$ changes
$\operatorname{Log}t$ by $2\pi\ii$.  If a single-valued backward branch existed on
the punctured disc, integrating the logarithmic derivative of
$\exp(aG)=t u$ around a small circle would give both $0$ and $2\pi\ii$.
A meromorphic branch has no pole away from the centre, since its
exponential would have an essential singularity there; it is therefore
subject to the same contradiction.

For the sewn family, the upper chart of Proposition~\ref{prop:weld}
contains $w=-1$ and its inverse uniformising coordinate $w(z)$ is
conformal, with $w(-1)=-1$.  Proposition~\ref{prop:upper-orbit} gives
$\Rr_b'(0)\ne0$, and differentiating the functional equation at
$z=-1$ gives
$\Rr_b'(-1)=\Rr_b'(0)/a(b)\ne0$.  Thus
$(\Kk_b^W)'(-1)=\Rr_b'(-1)w'(-1)\ne0$.
Proposition~\ref{prop:weld-parameter} preserves this simple zero on
a sufficiently small complex parameter disc about each real base.
\end{proof}

This monodromy is shared by every simple-zero branch with $F(-1)=0$;
it does not by itself identify the two families in (P2).

The families $G_b$ show that the functional equation and the normalisation
$F(0)=1$ allow entire solutions jointly holomorphic across a non-cusp
boundary point.  The integer at $z=-1$ distinguishes these families
from the inner sewn branch.  Thus (P0) still requires continuation of
the \emph{specified} branches with backward value $0$.

\begin{proposition}[Siegel pullbacks have no complete sewing band]
\label{prop:siegel-pullback}
Let $\lambda_c=\eee^{2\pi\ii\alpha}$ be Brjuno and different from $1$,
and let $b_c=\exp(\lambda_c\eee^{-\lambda_c})$.  Take $\rho>0$ small enough
that the chart of Lemma~\ref{lem:radial-chart} is univalent on a
neighbourhood of $\{|\zeta|\le\rho\}$ and its image lies in that
lemma's common inverse-chart disk, and put
\[
 \mathcal D_\rho=
 L_1(b_c)+(\log b_c)^{-1}H_{\lambda_c}(\{|\zeta|<\rho\}),
 \qquad
 \chi_c(w)=H_{\lambda_c}^{-1}
       \bigl((\log b_c)(w-L_1(b_c))\bigr).
\]
Then $E_{b_c}(\mathcal D_\rho)=\mathcal D_\rho$ and
$\chi_c\circ E_{b_c}=\lambda_c\chi_c$ there.  Define
\[
 \Omega_\rho=
 \bigl\{z\in\mathbb C:\ S_{b_c}(z+n)\in\mathcal D_\rho
                         \text{ for some }n\in\mathbb Z_{\ge0}\bigr\},\qquad
 V_c(z)=\lambda_c^{-n}\chi_c(S_{b_c}(z+n)).
\]
The following assertions hold.
\begin{enumerate}
\item $\Omega_\rho$ is nonempty, open and invariant under $z\mapsto z+1$ and
  $z\mapsto z-1$.  The formula for $V_c$ is independent of the admissible
  integer $n$ and defines a holomorphic function with
  $V_c(z+1)=\lambda_c V_c(z)$.  It has zeros, while
  $S_{b_c}^{-1}(1)\ne\varnothing$ and
  $S_{b_c}^{-1}(1)\cap\Omega_\rho=\varnothing$.
\item $\Omega_\rho$ contains no nonempty path-connected set $C$ with
  $C+m=C$ for any nonzero integer $m$.  In particular it contains no complete horizontal line,
  horizontal strip, or lift of an essential closed curve in
  $\Omega_\rho/\mathbb Z$.  The same statement holds if
  $\mathcal D_\rho$ is replaced by the maximal Siegel disk.
\item For every compact $K\subset\Omega_\rho$ there is one integer $N$
  such that the pulled-back upper coordinates
  \[
   V_{b,N}(z)=\lambda(b)^{-N}
      H_{\lambda(b)}^{-1}\!\left(
        (\log b)(S_b(z+N)-L_1(b))\right)
  \]
  are defined near $K$ for $b=\exp(\lambda\eee^{-\lambda})$,
  $\lambda=r\lambda_c$, and $r$ sufficiently close to $1$ from either
  side.  They converge to $V_c$ locally uniformly near $K$, together
  with all $z$-derivatives.
\end{enumerate}
\end{proposition}

\begin{proof}
The Schr\"oder identity and $|\lambda_c|=1$ make
$\mathcal D_\rho$ invariant in both directions.  The admissible sets
$\{z:S_{b_c}(z+n)\in\mathcal D_\rho\}$ increase with $n$.
If $m\ge n$ are both admissible, the Schr\"oder identity gives
$\chi_c(S_{b_c}(z+m))=\lambda_c^{m-n}\chi_c(S_{b_c}(z+n))$.
This proves independence, holomorphy and the translation identity;
admissibility after either integer translation follows by increasing
the pullback depth when needed.  The nonconstant entire function $S_{b_c}$
omits $0$, since its defining formula has at least one exponential
iterate.  Little Picard's theorem therefore makes it attain
$L_1(b_c)\ne0$; any such preimage lies in $\Omega_\rho$ and is a zero
of $V_c$.  The same argument shows that $S_{b_c}$ attains $1$.
By Corollary~\ref{cor:entry-depth}, $1$ belongs to the Julia
set of $E_{b_c}$.  Every point in its forward orbit remains in that set,
whereas $\mathcal D_\rho$ lies in the Fatou set.  Hence no $z$ with
$S_{b_c}(z)=1$ can satisfy the defining condition of $\Omega_\rho$.

Put $a=\log b_c$ and
$\mathcal B=\bigcup_{n\ge0}E_{b_c}^{-n}(\mathcal D_\rho)$.
Because $E_{b_c}$ maps $\mathcal D_\rho$ bijectively onto itself,
\[
 E_{b_c}^{-1}(\mathcal D_\rho)
 =\bigcup_{k\in\mathbb Z}
    (\mathcal D_\rho+2\pi\ii k/a),
\]
a disjoint union.  Thus $\mathcal D_\rho$ is a component of its first
preimage.  Inductively it is a component of every
$E_{b_c}^{-n}(\mathcal D_\rho)$: a path at level $n$ starting in
$\mathcal D_\rho$ maps to a path at level $n-1$ starting there, so the
induction hypothesis and the first-preimage description keep it in
$\mathcal D_\rho$.  It is also a component of $\mathcal B$, since any
path in the increasing union has compact image and is contained in one
finite level.

Suppose $C\subset\Omega_\rho$ were nonempty, path-connected and
$C+m=C$ for some $m>0$.  The set $S_{b_c}(C)$ is path-connected in
$\mathcal B$ and meets $\mathcal D_\rho$: for any $z\in C$, choose an
admissible depth $n$ and then a multiple $N$ of $m$ with $N\ge n$;
$z+N\in C$ and $S_{b_c}(z+N)\in\mathcal D_\rho$.  Hence
$S_{b_c}(C)\subset\mathcal D_\rho$.  But for a fixed $z\in C$,
$S_{b_c}(z-jm)\to L_2(b_c)$ as $j\to\infty$, whereas
$\chi_c(S_{b_c}(z-jm))=\lambda_c^{-jm}\chi_c(S_{b_c}(z))$ lies on a
compact circle inside the chart disk.  Its limit would be a fixed
point inside $\mathcal D_\rho$, necessarily $L_1(b_c)$, not $L_2(b_c)$.
The same argument uses the maximal Siegel disk in place of
$\mathcal D_\rho$.

Compactness of $K$ and the increasing admissible sets give one $N$
with $S_{b_c}(K+N)\Subset\mathcal D_\rho$.  Lemma~\ref{lem:radial-chart} and
Lemma~\ref{lem:lower-chart} give the convergence in (3); Cauchy's
formula gives all $z$-derivatives.
\end{proof}

\begin{proposition}[No bounded periodic sewing collar in the attracting basin]
\label{prop:no-basin-collar}
Let $\lambda_c=\eee^{2\pi\ii\alpha}\ne1$ with $\alpha$ Brjuno, and
continue the chosen lower regular family $S_\lambda$ to a
neighbourhood of $\lambda_c$ as in Lemma~\ref{lem:lower-chart}.
For $0<|\lambda|<1$ let
$\mathcal A_\lambda^{\rm all}
 =\bigcup_{n\ge0}E_{b(\lambda)}^{-n}(\mathcal A_\lambda)$
be the full attracting basin of $L_1(\lambda)$.
Let $V\subset\mathbb C$ be any nonempty connected open set with
$V+1=V$, and let $\lambda_j\to\lambda_c$ from $0<|\lambda|<1$
on the selected lower branch.  There is no locally bounded sequence
of holomorphic $1$-periodic functions $Q_j$ on $V$ such that
\[
 S_{\lambda_j}\bigl(z+Q_j(z)\bigr)
       \in\mathcal A_{\lambda_j}^{\rm all}
 \qquad(z\in V,\ j\ge1).
\]
In particular, taking $Q_j=0$ excludes any fixed lower-time collar
whose entire image stays in the attracting basin up to this boundary
point.
\end{proposition}

\begin{proof}
Suppose the sequence existed.  Montel's theorem gives a subsequence
$Q_j\to Q$ locally uniformly on $V$.  Set
$\phi_j(z)=z+Q_j(z)$ and $\phi(z)=z+Q(z)$.
The identity $\phi(z+1)=\phi(z)+1$ makes $\phi$ nonconstant and
$W=\phi(V)$ a nonempty connected open set with $W+1=W$.
Lemma~\ref{lem:lower-chart} gives
$H_j:=S_{\lambda_j}\circ\phi_j\to
 H:=S_{\lambda_c}\circ\phi$ locally uniformly on $V$.

We claim $H(V)$ avoids the Julia set.  Otherwise its open image
contains a repelling periodic point $p_c$, by the density
theorem~\cite[Thm.~4]{Bergweiler1993}.  Continue it to repelling periodic
points $p_j$ of the nearby maps.  Choose a closed disk
$\overline D\Subset V$ containing a zero of $H-p_c$ and none on its
boundary.  Rouch\'e's theorem gives a zero of $H_j-p_j$ in $D$ for
large $j$, contrary to
$H_j(V)\subset\mathcal A_{\lambda_j}^{\rm all}$.

Thus the connected set $H(V)=S_{\lambda_c}(W)$ lies in one Fatou
component $U$.  The identity $W+1=W$ and the functional equation
make $U$ invariant.  The Fatou-component classification for
exponential maps~\cite[\S1, Remarks on notation]{Rempe2006Topological}
says that, in the presence of this Siegel disk, every Fatou component
is an iterated preimage of it.  It applies here because, with
$a=\log b(\lambda_c)\ne0$, the affine map
$w\mapsto aw+\log a$ conjugates $E_{b(\lambda_c)}$ to
$u\mapsto\eee^u+\log a$.  Hence $U$ is the maximal Siegel disk.
Thus $W$ lies in the eventual maximal-Siegel domain of
Proposition~\ref{prop:siegel-pullback}(2), whose maximal-disk version
forbids a nonempty connected translation-invariant set.
\end{proof}

\begin{corollary}[No uniformly controlled boundary sewing]
\label{cor:no-uniform-sewing}
Keep the Brjuno parameter and disk $\mathcal D_\rho$ of
Proposition~\ref{prop:siegel-pullback}.  Let $r_j\to1$ with $r_j\ne1$,
$\lambda_j=r_j\lambda_c$, and
$b_j=\exp(\lambda_j\eee^{-\lambda_j})$.  For $0<s\le\rho$ put
\[
 \mathcal D_{j,s}=L_1(b_j)+(\log b_j)^{-1}H_{\lambda_j}(B(0,s)).
\]
On a complete horizontal strip $\Sigma$, there is no locally bounded
sequence of holomorphic $1$-periodic functions $Q_j$ with the following
property: for every compact $K\subset\Sigma$ there are an integer
$N_K\ge0$ and a radius $s_K<\rho$ such that, for all large $j$,
\[
 S_{b_j}(z+Q_j(z)+N_K)\in\mathcal D_{j,s_K}
 \qquad(z\in K).
\]
\end{corollary}

\begin{proof}
By Montel's theorem a subsequence of $Q_j$ converges locally uniformly
on $\Sigma$ to a holomorphic $1$-periodic $Q$.  The lower-chart lemma and
the radial-chart lemma imply, for each compact $K$, that
$S_{b_c}(z+Q(z)+N_K)$ belongs to the image of the closed disk
$\{|\zeta|\le s_K\}$ under the boundary chart.  This image is contained
in $\mathcal D_\rho$, since $s_K<\rho$.  Hence
$z+Q(z)\in\Omega_\rho$ for every $z\in\Sigma$.
The image $(\mathrm{id}+Q)(\Sigma)$ is nonempty and path-connected, and
periodicity gives $(\mathrm{id}+Q)(\Sigma)+1=(\mathrm{id}+Q)(\Sigma)$.
This contradicts Proposition~\ref{prop:siegel-pullback}(2).
\end{proof}

In particular, a boundary solution cannot have overlapping upper-Siegel
and lower-repelling regular representations on a complete horizontal
line with $1$-periodic time corrections: the lower times
$z+Q(z)$ would form a path-connected, translation-invariant subset of
the corresponding eventual Siegel domain, contradicting part~(2).

\begin{proposition}[The local chart has a natural parameter boundary]
\label{prop:chart-boundary}
The family $H_\lambda$ for~\eqref{eq:radial-normal-form} has no jointly
holomorphic extension in $(\lambda,v)$ from either side of the unit circle
through any root of unity, on any fixed nonzero disk in $v$.  Consequently
it has no such extension across any open arc of the unit circle.
\end{proposition}

\begin{proof}
If an extension existed at a primitive $q$th root $\lambda_0$, the
identity $f_\lambda\circ H_\lambda=H_\lambda\circ(\lambda\,\cdot)$
and $H_\lambda'(0)=1$ would persist at $\lambda_0$.  Iterating $q$ times
would give $f_{\lambda_0}^{\circ q}\circ H_{\lambda_0}=H_{\lambda_0}$.
Local invertibility of $H_{\lambda_0}$ would make
$f_{\lambda_0}^{\circ q}$ the identity near $0$, hence everywhere by the
identity theorem for entire functions.  This would make $f_{\lambda_0}$
injective, contrary to its $2\pi\ii$-periodicity.  Roots of unity are dense
on the unit circle.
\end{proof}

This obstruction concerns the chart, not the sewn Kneser-type function:
parameter singularities in the separate charts might cancel in the sewing.
The monogenic parameter regularity of~\cite{MarmiCarminati2008} is therefore
the relevant comparison for the local step.

\begin{lemma}[Root isolation implies normality]\label{lem:root-normal}
Let $\mathcal F\subset\mathcal O(B(0,R))$ satisfy $F(0)=1$, $F(z)\ne0$
throughout $B(0,R)$, and
$F^{-1}(1)\cap B(0,R)=\{0\}$ for every $F\in\mathcal F$.
Then the normalised logarithms $g_F$ with $g_F(0)=0$, the functions $F$,
and their reciprocals $1/F$ are uniformly
bounded on every $|z|\le r<R$, with bounds depending only on $r/R$.
In particular $|F'(0)|\le C/R$ for an absolute constant $C$.
If $F(-1)=0$ and $F(z)=\exp(aF(z-1))$ on $B(0,R)$, then
$g_F(z)=aF(z-1)$; when $|a|$ is bounded below, the family is also
uniformly bounded on every smaller disk $B(-1,r)$.  For these
solutions the zero-free hypothesis is automatic.
\end{lemma}

\begin{proof}
Choose the holomorphic logarithm $g_F$ with $g_F(0)=0$.  Root isolation
makes $h_F=\tfrac12+g_F/(4\pi\ii)$ omit $0$ and $1$, while
$h_F(0)=\tfrac12$.  Schottky's theorem bounds $h_F$, hence $g_F$,
$F=\exp g_F$, and $1/F$, uniformly on compact subdisks.  Cauchy's
estimate at radius $R/2$ gives the derivative bound.  The functional
equation makes $aF(z-1)$ a logarithm of $F(z)$ normalised at $0$,
so it equals $g_F$.
\end{proof}

For a zero-free $F\in\mathcal O(B(0,R_0))$ with $F(0)=1$ and
$F\not\equiv1$, let $d_F$ be the distance from $0$ to the nearest
other preimage of $1$, or $\infty$ if there is none.  Applying the lemma
on every disk of radius $s<\min(R_0,d_F)$ gives
$|F'(0)|\min(R_0,d_F)\le C$.  Thus derivative blow-up on a fixed disk
forces another preimage of $1$ to approach $0$.

\begin{proposition}[Reduction]\label{prop:brjuno}
Let $W$ be a nonempty open subset of a common regular $z$-domain for both
families, with $\operatorname{Re}z>-2$ on $W$.  Suppose the following hold.
\begin{enumerate}
\item[(P0)] The branch $\kappa$ from $b>\eta$ has a jointly holomorphic
  continuation on a connected exterior parameter corridor reaching
  $U_{\mathrm{out}}\times W$.  The sewing family of
  Theorem~\ref{thm:W1} has a jointly holomorphic continuation on a
  connected interior parameter corridor reaching $U_{\mathrm{in}}\times W$.
  Both continuations use the stated branches and have $W$ in their common
  regular $z$-domain.
\item[(P1)] The base-continued family $\kappa_b(z)$ on
  $U_{\mathrm{out}}\times W$ and the two-sided regular family
  $\Kk^W_b(z)$ on $U_{\mathrm{in}}\times W$ are locally uniformly bounded
  up to $\Gamma$: for every compact subarc $\Gamma'\subset\Gamma$ and
  $V\Subset W$, they are bounded on a neighbourhood of $\Gamma'$ on their
  respective sides, uniformly for $z\in V$.
\item[(P2)] Let $D\subset W$ be a countable set with an accumulation point
  in $W$.  For each $z\in D$, the two one-sided radial limits exist and
  are equal at almost every point of $\Gamma$ (in arc length).
\end{enumerate}
Then $\kappa$ continues jointly holomorphically in $(b,z)$ across $\Gamma$
on $W$ and coincides with the continued $\Kk^W$ on
$U_{\mathrm{in}}\times W$.
The equality extends to any common connected regular $z$-domain by the
identity theorem.
\end{proposition}

\begin{proof}
Intersect the full-measure subsets in (P2) over the countable set $D$.
At a point $b_c$ in the intersection, (P1) makes each radial family normal
as a family of holomorphic functions on $W$.  Every subsequential limit
has the values prescribed on $D$; the identity theorem makes it unique.
Thus the two radial families converge locally uniformly on $W$ and have
the same limit.  Fix $z\in W$.  On each side, bounded holomorphic functions
have non-tangential limits almost everywhere (Fatou); the radial limits
just obtained agree with these.
The $H^\infty$ gluing theorem across an analytic arc therefore gives a
holomorphic function of $b$ through $\Gamma$.  The bounds in (P1) are locally
uniform in $z$.  At a point of $\Gamma$, approach from $U_{\mathrm{in}}$:
the functions of $z$ form a locally bounded family and converge pointwise to
the glued values.  Montel's theorem shows that the limit is holomorphic in
$z$.  The glued function is thus separately holomorphic and locally bounded,
so Hartogs' theorem makes it jointly holomorphic.  It equals $\kappa$ on the
outer side and $\Kk^W$ on the inner side; the identity theorem extends the
equality in $z$ wherever both are defined on a common connected domain.
\end{proof}

\begin{proposition}[$H^1$ alternative to (P1)]
\label{prop:hardy-reduction}
In Proposition~\ref{prop:brjuno}, (P1) may be replaced by the following
condition.  In every local conformal coordinate $\xi=s+\ii t$ straightening
$\Gamma$, for each compact interval $I$ on the straightened arc and each
$V\Subset W$, there are $\epsilon>0$ and $C<\infty$ such that
on both sides
\begin{equation}\label{eq:hardy-bound}
 \sup_{0<|t|<\epsilon}\int_I
 \sup_{z\in V}|F_\pm(s+\ii t,z)|\,ds\le C,
\end{equation}
where $F_+$ and $F_-$ denote the respective families from (P0).
The conclusion of Proposition~\ref{prop:brjuno} still holds.
\end{proposition}

\begin{proof}
Shrink $I$ slightly to avoid the ends of the coordinate rectangle.
By Fubini and~\eqref{eq:hardy-bound}, the two vertical sides of a smaller
rectangle can be chosen with integrable $\sup_{z\in V}|F_\pm|$; its
opposite horizontal edge is away from the arc.  The local $H^1$ boundary theorem
(or the disk theorem after conformally mapping this rectangle)~\cite{Duren1970}
applied for each fixed $z$ gives
$L^1(I)$ traces $f_\pm(\cdot,z)$, with convergence in $L^1$ on the
smaller interval as $t\to0^\pm$.  The radial parameter paths in (P2)
meet the analytic arc transversely, so their limits agree almost everywhere
with these Hardy traces.  The bound~\eqref{eq:hardy-bound}
and Cauchy's formula in $z$ show that
$z\mapsto f_\pm(\cdot,z)$ are holomorphic as $L^1(I)$-valued maps:
on each compact $z$-circle the $L^1$ norms are uniformly bounded, so
Bochner dominated convergence passes the boundary limit through the
Cauchy integral.
By (P2), their difference vanishes in $L^1(I)$ at every $z\in D$.
The Banach-valued identity theorem makes the difference zero on $W$.

The two piecewise defined families are locally integrable in $(s,t,z)$
by~\eqref{eq:hardy-bound}.  Equality of their boundary traces removes
the jump term in the distributional $\bar\partial_\xi$ derivative.
The distributional $\bar\partial_z$ derivative also vanishes, since the
families are jointly holomorphic away from the measure-zero arc.
Weyl's lemma for $\bar\partial$ therefore extends the glued family
jointly holomorphically across the arc.  The identity theorem gives the
same conclusion on the common connected regular $z$-domain as before.
\end{proof}

\begin{remark}
The exponent $1$ is sharp for this abstract boundary-gluing argument.
Indeed, on opposite sides of the real $\xi$-axis let
$F_\pm(\xi,z)=1+z/\xi$.  The two boundary values agree almost
everywhere and $F_\pm(\xi,0)=1$.  For every $0<p<1$ and compact
$I,V$, the integrals
$\int_I\sup_{z\in V}|F_\pm(s+\ii t,z)|^p\,ds$ stay uniformly bounded
as $t\to0^\pm$, but the common function has a pole at $\xi=0$ for
$z\ne0$.  Thus replacing~\eqref{eq:hardy-bound} by a bare $H^p$
bound with $p<1$ would not prove the gluing statement; this example
does not address additional constraints of the tetration equation.
\end{remark}

Lemma~\ref{lem:root-normal} gives a geometric sufficient condition for (P1)
when the common $z$-domain contains $B(0,R)\cup B(-1,R)$: if both families
have no preimage of $1$ in $B(0,R)$ other than $0$ for every base in
the two collars, then (P1) holds on every smaller disk $W=B(0,r)$.
Along a connected parameter corridor, the argument principle transports
this root count from one base where $0$ is the only, simple root, provided
$F_b$ is defined near the closed disk $\overline{B(0,R)}$ and
$F_b(z)\ne1$ on the fixed circle $|z|=R$ throughout the corridor.
This contour exclusion near $\Gamma$ has not been proved.
Corollary~\ref{cor:sewn-cusp-local} supplies a starting base and a fixed
root-isolation disk for the inner sewing branch near the real cusp;
transport to $\Gamma$ still requires the contour exclusion.

The open set $W$ in this proposition is essential.  At the single point
$z_0=0$, both families equal $1$ by normalisation.  More generally, at every
$n\in\mathbb Z_{\ge0}$ they equal $E_b^{\circ n}(1)$, including all derivatives
with respect to $b$ wherever the families are holomorphic.  These integer
heights therefore cannot test (P2).  Corollary~\ref{cor:backward-rigidity}
adds $z=-1$ and all its base derivatives when both specified branches
retain their zero backward label; that point distinguishes them from
entire lower regular translates, but not from each other.
An accumulating set of noninteger heights
is needed.  The monodromy in Proposition~\ref{prop:upper-orbit} gives an
explicit illustration: after any nonzero number of parameter turns, the upper
regular germs still take the same values $E_b^{\circ n}(1)$ at every
nonnegative integer $n$, although their derivatives at $z=0$ differ.
A pointwise version of (P1)--(P2) would be vacuous and could not
establish continuation of the family.
Even finitely many $z$-derivatives at every integer do not remove this
ambiguity.  For any $m\ge0$ and any solution $F$ regular near the
integers, the local solution
$F_\epsilon(z)=F(z+\epsilon(\eee^{2\pi\ii z}-1)^{m+1})$ obeys the same
functional equation and has the same $z$-jets through order $m$ at
each integer where both are defined.  For small nonzero $\epsilon$ it
generally differs at noninteger heights; this observation makes no
claim that $F_\epsilon$ remains in the two-chart sewing class.
Condition (P0) is also substantive: Proposition~\ref{prop:weld-parameter}
constructs $\Kk^W$ on discs along the whole real segment $(1,\eta)$, and
Theorem~\ref{thm:W1} adds quantitative bounds near $\eta$; neither reaches
an arbitrary Shell--Thron boundary arc.  Proposition~\ref{prop:exterior-real-parameter} supplies a jointly
holomorphic exterior germ along the entire real ray $b>\eta$.
Proposition~\ref{prop:exterior-curve-collar} gives an explicit
\emph{geometric} base corridor on which its unmarked initial curve and
two repelling ends persist up to a non-cusp boundary arc.
Proposition~\ref{prop:exterior-marked-path} certifies transport of the
normalising point $w=1$ along one path in that corridor, hence of the
exterior germ at $z=0$; Proposition~\ref{prop:exterior-brjuno-path}
places its endpoint at a prescribed Brjuno parameter.
Proposition~\ref{prop:exterior-pullback-cell} pulls the whole initial
region back through $28$ iterates to a Jordan cell containing $w=1$,
even at that endpoint.
Proposition~\ref{prop:exterior-flat-cylinder} gives a uniform
quasiconformal model of these quotient cylinders, and
Proposition~\ref{prop:exterior-brjuno-germ-limit} gives one-sided
limits on a fixed small height disk about $z=0$ along a short
boundary subarc, without an arithmetic restriction.
Proposition~\ref{prop:exterior-analytic-germ-corridor} upgrades these
limits to a jointly holomorphic germ across the subarc in a genuine
complex base corridor, certified on a rectangle with
$|\operatorname{Im}\theta|\le0.01$.
Proposition~\ref{prop:exterior-real-seam-anchor} proves one transverse
seam crossing at the real seed and reduces its possible loss along
the path to a scalar elliptic tangency test.
Lemma~\ref{lem:exterior-seam-end-slope} confines any such tangency
to a finite-height cylinder band, without yet giving its explicit
width.
Propositions~\ref{prop:exterior-frozen-abel}
and~\ref{prop:exterior-frozen-correction} give an explicit
positive-slope approximate Abel coordinate and a global $L^4$
bound on its correction to the exact uniformisation.
Proposition~\ref{prop:exterior-seam-monotonicity} differentiates the
equation vertically and upgrades this to a pointwise derivative
bound, excluding all seam tangencies on a wider box that also
contains interior bases.
Proposition~\ref{prop:interior-chord-overlap} certifies at one such
base that a nontrivial subarc of the exterior chord lies in the
immediate attracting basin.  This physical overlap does not
identify the two quotient coordinates.
The single-crossing criterion of
Proposition~\ref{prop:exterior-single-crossing} is therefore met.
Corollary~\ref{cor:exterior-full-height} supplies a common regular
height neighbourhood of $[0,1]$ across a non-cusp boundary subarc
for the exterior family.  This proves the exterior full-height
substep of (P0); the inner sewn family has not been transported
to the same subarc, so (P0) as a whole remains open.
Corollary~\ref{cor:exterior-wide-full-height} carries that same
full-height branch along the certified path to a strictly interior
base where Proposition~\ref{prop:interior-chord-overlap} supplies a
physical attracting overlap.
Proposition~\ref{prop:interior-sewing-circle} separately certifies an
embedded lower-to-upper sewing circle and a locally holomorphic sewn
family at this interior base.  Its parameter neighbourhood has not
been connected to the real inner sewn family, and the two local
sewing constructions have not been identified.
Proposition~\ref{prop:interior-path-crossing} certifies a transverse
self-intersection of that lower regular curve at another interior base
on the proposed straight parameter path.  Thus the same circle cannot
remain an embedded positive-slope seam throughout that path; a
parameter detour or a change of seam remains possible.
Proposition~\ref{prop:interior-quotient-loop} now places the entire
lower sewing circle, after $76$--$77$ forward iterates, as one
embedded essential loop in the exterior quotient at that same base.
This supplies a common annular physical chart, but does not yet
identify the compactified quotient surfaces or their normalised
coordinates.
Corollary~\ref{cor:interior-two-returns} certifies another,
disjoint return loop after $14$--$15$ iterates, so the physical
return alone does not select a unique quotient comparison.
Corollary~\ref{cor:interior-end-germs} also matches both regular
charts to degree-one puncture germs of the exterior quotient.
For either chosen return loop, the unproved part is its analytic
connection through the two compact middle bands to those germs.
Corollary~\ref{cor:exterior-boundary-zero} continues this exterior
family through height $z=-1$, where it has a simple zero and the
same backward label as the inner sewn family.  That label does not
by itself identify the two branches.
Corollary~\ref{cor:exterior-boundary-normal} gives a fixed disc of
uniform root isolation and the exterior local bounds required by
(P1) on this boundary subarc.  The inner bounds remain open.
This exterior substep is not established by the parameter argument
of~\cite[Prop.~3 and \S5]{Paulsen2019}.
Proposition~\ref{prop:upper-orbit} transports the upper normalised germ
along every compact inner path avoiding $\lambda=0$,
but gives neither a boundary-uniform height disk nor the sewing correction.
Corollary~\ref{cor:entry-depth} shows that no fixed finite pullback depth
at the normalising orbit can supply such a boundary chart: the required
depth diverges at every Brjuno point.
Proposition~\ref{prop:basin-inradius} further shows that any chart on a
fixed connected height domain whose image stays in the immediate
attracting basin and whose value at $z=0$ is $1$ must collapse to the
constant $1$ there, even without a prior bound.  Such a chart cannot
carry the normalised value at $z=1$ near a Brjuno boundary point.
Proposition~\ref{prop:no-basin-collar} also rules out a locally bounded
periodic lower-time correction on a fixed translation-invariant
collar whose image remains in the full attracting basin up to that point.

Non-Brjuno numbers have measure zero, so (P2) may be verified at Brjuno points
only.  There, Lemma~\ref{lem:radial-chart} makes the upper charts converge
locally at $L_1$ from both sides to the Siegel chart.
Lemmas~\ref{lem:boundary-repelling} and~\ref{lem:lower-chart} give locally
uniform convergence of the entire lower superfunction at every point of
$\Gamma$ away from the cusp.
Proposition~\ref{prop:siegel-pullback}(2) rules out a complete sewing
band obtained from the Siegel chart even after arbitrary finite forward
pullbacks.  Corollary~\ref{cor:no-uniform-sewing} further shows that, on a
fixed complete strip, locally bounded periodic lower times cannot have
locally uniform finite-depth entry into a fixed upper chart disk.
The chart convergence in part (3) is only local on
$\Omega_\rho$ and cannot furnish an essential sewing curve.  Thus (P2)
needs control of the sewn solutions that does not assume a boundary
transition map on such a band.  The pointwise \texttt{fatou.gp} values in
Section~\ref{sec:bndtest} are compatible with a smooth join at the
Diophantine point $\alpha=(\sqrt2-1)/5$, but are not identified with the
two families in (P2).  We have not proved (P0),
(P1), its Hardy-norm alternative~\eqref{eq:hardy-bound}, or (P2);
in particular the proposed boundary-band argument cannot
establish (P2).

\section{What is claimed}\label{sec:claims}

\paragraph{Proved.}
\begin{itemize}
\item Elementary facts: Lemma~\ref{lem:param}, the conditional first-mode
  expansion of Proposition~\ref{prop:firstmode}, the scale $\Lambda$ and constant
  $C_\eta$ of Proposition~\ref{prop:saddle}, the formal
  identities~\eqref{eq:Bn}, \eqref{eq:germzeta} and~\eqref{eq:kapparho}, and
  translation invariance of every $\kappa_n$.
\item The Kneser-free reformulation: Lemma~\ref{lem:realT}, the welding identity
  (Theorem~\ref{thm:welding}), and the sewing construction
  (Proposition~\ref{prop:weld}).
\item The parabolic limit: the confluence statement (C) (Theorem~\ref{thm:C}).
  It gives $\tau_n(b)\to B_n\eee^{2\pi\ii na}$ for every $n$ and $B_0=0$;
  Proposition~\ref{prop:horn-nonzero} gives $B_1\ne0$, hence the ratio limit
  and limiting phase (Corollary~\ref{cor:tauconv}).
  The supporting lemmas are~\ref{lem:res}--\ref{lem:petal},
  \ref{lem:band} and~\ref{lem:modes}.
\item The Kneser side: the decay estimate with explicit constants
  (Theorem~\ref{thm:W2}), and the characterisation
  $\mathcal C_b(Y)=\{\Kk^W_b\}$ (Theorem~\ref{thm:char}).  Hence
  $\hat c_n\to B_n\eee^{2\pi\ii na}$ for every $n$ holds for the unique solution;
  the ratio limit of Conjecture~\ref{conj:horn2} follows.  The
  solution has an attracting-regular representation above, a repelling-regular
  representation below, and a seam half a period below the real axis
  (Corollary~\ref{cor:Kw}).
\item Local and uniform statements in the base: Theorems~\ref{thm:W1},
  \ref{thm:bv} and~\ref{thm:unif}.
\end{itemize}
These proofs were checked by two independent adversarial reviews and a
second-round review of the repairs.

\paragraph{Conditional.}
The same conclusions for the base-continued Kneser family $\Kk_b$ follow once
that family belongs to $\mathcal C_{b'}$ for $b'$ in some open set of bases near
$(b_1,\eta)$ (Theorem~\ref{thm:unif}), or once it satisfies the ladder
hypothesis of Theorem~\ref{thm:bv}.  This is a qualitative statement about the
definition of the family in~\cite{Paulsen2019} (Remark~\ref{rem:open}).  The
complex-base construction computes $\kappa$ in exactly this form, and the
ladder data are consistent with it.

\paragraph{Numerical, without certified error bounds.}
The horn coefficients~\eqref{eq:A1} and~\eqref{eq:B12}--\eqref{eq:kappainf}, the
constant $a$, and all ladder values.  The deficit law~\eqref{eq:kappa1}, i.e.\
analyticity of $\log\tau_n$ in the unfolding parameter $p$ with
$C=0.0101990063$.  Agreement across settings is not a bound on total error.

\paragraph{Not proved.}
A closed form for the first-order coefficients $\kappa^{(n)}$.  The absence of
every negative mode for the ladder solutions, which is a statement about
$\kappa$ rather than $\Kk^W$ (for $\Kk^W$ it holds by construction).  Large-order
growth of the $c_n$ beyond the bounds of Theorem~\ref{thm:W2}.  Global
uniqueness of Kneser-type solutions across the Shell--Thron boundary.

\paragraph{Prior art.}
The qualitative statement $\Kk_b\ne \Rr_b$ is
Paulsen's~\cite[Prop.~4]{Paulsen2019}, whose quantifier is over open
subintervals; the pointwise measurements here are complementary, not a
reduplication.  We have not found the scale $\Lambda$, its mechanism, or the
quantity $\chat$ discussed in the tetration literature; the search
term ``$4\pi^{2}$'' does not occur in~\cite{Paulsen2019}.  We have \emph{not}
carried out a systematic search of the exponentially-small-splitting literature
for an equivalent statement about the exponential family, and regard the
attribution of Section~\ref{sec:parabolic} as open.

\section{Reproducibility}\label{sec:repro}

The computation of $\Kk$ uses the $\theta$-mapping construction with two fixed
points, seeded along ladders $b+\ii y$ with $y\downarrow0$, at requested
precisions up to $52$ digits; the regular solution $\Rr$ is computed
independently from the Schr\"oder series.  Both are part of the
\texttt{kneser} Python library.  The analyses use stored ladder data and separate parabolic computations.
The relevant scripts include:
\texttt{first\_mode\_check.py} (Table~\ref{tab:qz}),
\texttt{c1\_invariant.py} (Table~\ref{tab:chat}),
\texttt{theta\_spectrum.py} (Table~\ref{tab:spectrum},
Observation~\ref{obs:spectrum}),
\texttt{lambda\_law\_analysis.py} (Lemma~\ref{lem:param},
Proposition~\ref{prop:saddle}),
\texttt{parabolic\_horn.py} (Section~\ref{sec:horn}; this one is independent of
all the ladder data),
\texttt{parabolic\_horn\_inverse.py}, \texttt{kappa\_ladder.py},
\texttt{kappa\_gauge\_check.py}, \texttt{deficit\_law.py}, and, for
Section~\ref{sec:transition}, \texttt{transition\_map.py},
\texttt{eta\_approach.py}, \texttt{eta\_sweep.py}, \texttt{confluence\_check.py} and \texttt{weld\_solve.py} (output in
\texttt{docs/data/eta-sweep.txt}).
The mathematical interpretation of the last two follows the corrected
gauge distinction and sign calculation above, not their earlier explanatory
comments.  Reproducing historical fit values also requires the recorded
data selection, correction order, and external-point phase convention.
The additional \texttt{deficit\_audit.py} records those choices and matched-order
fits, omission checks, and archived-ladder sensitivity in
\texttt{docs/data/deficit-audit.json}.
The external values of $\Kk$ are the published \texttt{fatou.gp} reference set
of the tetration forum and the number~\eqref{eq:paulsen}
from~\cite{Paulsen2019}.

One instrument-level caveat is worth recording because it invalidated an earlier
version of these measurements: in the complex-base builder the fixed points were
computed at a hard-coded precision of $30$ digits, and since the repelling-side
super-function is evaluated as $E^{d}(L+\lambda^{z-d})$ with depth $d$
proportional to the requested precision, that error was amplified by
$|\lambda|^{d}$ --- so that asking for more digits made the residual
\emph{worse}.  After the fix, residuals fall to $\sim5\cdot10^{-22}$ at $20$
digits and the signal-to-noise ratio of $D$ reaches $10^{10}$--$10^{13}$.

\begin{thebibliography}{99}

\bibitem{BanuelosJanakiraman}
R.~Ba\~nuelos and P.~Janakiraman,
\emph{$L^p$-bounds for the Beurling--Ahlfors transform},
Trans.\ Amer.\ Math.\ Soc.\ \textbf{360} (2008), no.~7, 3603--3612.
\href{https://www.math.purdue.edu/~banuelos/Papers/Beur_Ahl.pdf}{Author PDF}.

\bibitem{Bergweiler1993}
W.~Bergweiler,
\emph{Iteration of meromorphic functions},
Bull.\ Amer.\ Math.\ Soc.\ (N.S.)\ \textbf{29} (1993), no.~2, 151--188.
\href{https://arxiv.org/abs/math/9310226}{arXiv:math/9310226}.

\bibitem{BuffEpstein2002}
X.~Buff and A.~L.~Epstein,
\emph{A parabolic Pommerenke--Levin--Yoccoz inequality},
Fund.\ Math.\ \textbf{172} (2002), 249--289.
\href{https://www.impan.pl/shop/en/publication/transaction/download/product/89261}{Publisher PDF}.

\bibitem{Cheritat2022}
A.~Ch\'eritat,
\emph{Near parabolic renormalization for unicritical holomorphic maps},
Arnold Math.\ J.\ \textbf{8} (2022), 169--270.
\href{https://doi.org/10.1007/s40598-020-00172-6}{doi:10.1007/s40598-020-00172-6}.

\bibitem{ChristopherRousseau}
C.~Christopher and C.~Rousseau,
\emph{The moduli space of germs of generic families of analytic
diffeomorphisms unfolding a parabolic fixed point},
Int.\ Math.\ Res.\ Not.\ IMRN (2014); arXiv:0809.2167.

\bibitem{Duren1970}
P.~L.~Duren, \emph{Theory of $H^p$ Spaces}, Academic Press, New York, 1970.

\bibitem{Ecalle1981}
J.~\'Ecalle,
\emph{Les fonctions r\'esurgentes}, Publ.\ Math.\ d'Orsay, 1981.

\bibitem{FischerGrauert1965}
W.~Fischer and H.~Grauert,
\emph{Lokal-triviale Familien kompakter komplexer Mannigfaltigkeiten},
Nachr.\ Akad.\ Wiss.\ G\"ottingen Math.-Phys.\ Kl.\ II (1965), 89--94.

\bibitem{Glutsyuk2001}
A.~A.~Glutsyuk,
\emph{Confluence of singular points and the nonlinear Stokes phenomenon},
Trans.\ Moscow Math.\ Soc.\ \textbf{62} (2001), 49--95.

\bibitem{Kneser1950}
H.~Kneser,
\emph{Reelle analytische L\"osungen der Gleichung
$\vartheta(\vartheta(x))=\eee^{x}$ und verwandter Funktionalgleichungen},
J.\ Reine Angew.\ Math.\ \textbf{187} (1950), 56--67.

\bibitem{Kouznetsov2009}
D.~Kouznetsov,
\emph{Solution of $F(z+1)=\exp(F(z))$ in complex $z$-plane},
Math.\ Comp.\ \textbf{78} (2009), 1647--1670.

\bibitem{KouznetsovTrappmann2010}
D.~Kouznetsov and H.~Trappmann,
\emph{Portrait of the four regular super-exponentials to base $\sqrt2$},
Math.\ Comp.\ \textbf{79} (2010), 1727--1756.

\bibitem{KouznetsovTrappmann2012}
H.~Trappmann and D.~Kouznetsov,
\emph{Computation of the two regular super-exponentials to base
$\exp(1/\eee)$},
Math.\ Comp.\ \textbf{81} (2012), 2207--2227.

\bibitem{LaubnerSchleicherVicol2008}
B.~Laubner, D.~Schleicher and V.~Vicol,
\emph{A combinatorial classification of postsingularly finite complex
exponential maps},
Discrete Contin.\ Dyn.\ Syst.\ \textbf{22} (2008), no.~3, 663--682.
\href{https://doi.org/10.3934/dcds.2008.22.663}{doi:10.3934/dcds.2008.22.663}.

\bibitem{Lavaurs1989}
P.~Lavaurs,
\emph{Syst\`emes dynamiques holomorphes: explosion de points p\'eriodiques
paraboliques}, Th\`ese, Universit\'e Paris-Sud, Orsay, 1989.

\bibitem{MarmiCarminati2008}
S.~Marmi and C.~Carminati,
\emph{Linearization of germs: regular dependence on the multiplier},
Bull.\ Soc.\ Math.\ France \textbf{136} (2008), no.~4, 533--564.
\href{https://doi.org/10.24033/bsmf.2565}{doi:10.24033/bsmf.2565}.

\bibitem{AhlforsBers1960}
L.~Ahlfors and L.~Bers,
\emph{Riemann's mapping theorem for variable metrics},
Ann.\ of Math.\ (2) \textbf{72} (1960), 385--404.

\bibitem{MRR2004}
P.~Marde\v{s}i\'c, R.~Roussarie and C.~Rousseau,
\emph{Modulus of analytic classification for unfoldings of generic parabolic
diffeomorphisms},
Mosc.\ Math.\ J.\ \textbf{4} (2004), no.~2, 455--502.

\bibitem{Oudkerk1999}
R.~Oudkerk,
\emph{The parabolic implosion for
$f_0(z)=z+z^{\nu+1}+O(z^{\nu+2})$},
Ph.D.\ thesis, University of Warwick, 1999;
see also \emph{The parabolic implosion: Lavaurs maps and strong convergence for
rational maps}, in: Value Distribution Theory and Complex Dynamics,
Contemp.\ Math.\ \textbf{303}, Amer.\ Math.\ Soc., 2002, 79--105.

\bibitem{PaulsenCowgill2017}
W.~Paulsen and S.~Cowgill,
\emph{Solving $F(z+1)=b^{F(z)}$ in the complex plane},
Adv.\ Comput.\ Math.\ \textbf{43} (2017), 1261--1282.
\href{https://doi.org/10.1007/s10444-017-9524-1}{doi:10.1007/s10444-017-9524-1}.

\bibitem{Paulsen2019}
W.~Paulsen,
\emph{Tetration for complex bases},
Adv.\ Comput.\ Math.\ \textbf{45} (2019), 243--267.
\href{https://doi.org/10.1007/s10444-018-9615-7}{doi:10.1007/s10444-018-9615-7}.

\bibitem{Rempe2006Topological}
L.~Rempe,
\emph{Topological dynamics of exponential maps on their escaping sets},
Ergodic Theory Dynam.\ Systems \textbf{26} (2006), no.~6, 1939--1975.
\href{https://arxiv.org/abs/math/0309107}{arXiv:math/0309107}.

\bibitem{Rempe2008}
L.~Rempe,
\emph{Siegel disks and periodic rays of entire functions},
J.\ Reine Angew.\ Math.\ \textbf{624} (2008), 81--102.
\href{https://doi.org/10.1515/CRELLE.2008.081}{doi:10.1515/CRELLE.2008.081}.

\bibitem{RousseauTeyssier2008}
C.~Rousseau and L.~Teyssier,
\emph{Analytical moduli for unfoldings of saddle-node vector fields},
Mosc.\ Math.\ J.\ \textbf{8} (2008), no.~3, 547--614.

\bibitem{Schleicher2003}
D.~Schleicher,
\emph{Attracting dynamics of exponential maps},
Ann.\ Acad.\ Sci.\ Fenn.\ Math.\ \textbf{28} (2003), no.~1, 3--34.
\href{https://www.math.stonybrook.edu/preprints/ims00-04.pdf}{Preprint PDF}.

\bibitem{Shishikura2000}
M.~Shishikura,
\emph{Bifurcation of parabolic fixed points},
in: \emph{The Mandelbrot Set, Theme and Variations},
London Math.\ Soc.\ Lecture Note Ser.\ \textbf{274},
Cambridge Univ.\ Press, 2000, 325--363.

\bibitem{TrappmannKouznetsov2011}
H.~Trappmann and D.~Kouznetsov,
\emph{Uniqueness of holomorphic Abel functions at a complex fixed point pair},
Aequationes Math.\ \textbf{81} (2011), 65--76.

\bibitem{Voronin1981}
S.~M.~Voronin,
\emph{Analytic classification of germs of conformal mappings
$(\mathbb C,0)\to(\mathbb C,0)$},
Funct.\ Anal.\ Appl.\ \textbf{15} (1981), 1--13.

\end{thebibliography}

\end{document}
